% !TeX root = ../../Book1.tex
% 扩充编号：A01、A02、A03。八节正文、章首、小结与八道习题初稿完成；待正式审核。
% §29.1：沿用第8章的谱定义和内积约定；Neumann级数与一般有界自伴谱界在此直接证明。
% 二次型证明衔接 b1:lem:compact-self-adjoint-extreme，不使用其紧性结论。
% 文献实读：Axler2020Measure 印刷 pp.288--296（10.22、10.25、10.32--10.37），
% 对应几何级数求逆、可逆集开性、谱的基本性质及对角／乘法模型。
% 谱积分、无界算子与酉群的实读定位另见各节注释及参考文献原文件/来源记录.md。
\OptionalChapter{一般谱理论}{b1:ch:29}
\BookChapterNumber{b1:ch:29}

\begin{chapterabstract}
紧自伴算子可以沿特征方向展开，但乘法算子可能没有任何特征向量。
本章研究去掉紧性以后，自伴算子的谱怎样描述方程的可解性，
以及怎样以乘法表示和谱积分代替特征向量求和。
有界算子的讨论从预解估计开始；转向无界算子时，
还须把定义域纳入伴随、乘积和极限的每一次运算。
最后研究自伴算子与强连续酉群的对应。
\end{chapterabstract}

\begin{learninggoals}
\begin{itemize}
\item\label{b1:item:29-goal-resolvent}用预解算子研究线性方程，区分特征值、连续谱与离散谱，并能计算基本模型的谱。
\item\label{b1:item:29-goal-calculus}理解连续函数演算怎样由多项式逼近得到，以及谱测度怎样把演算扩展到 Borel 函数。
\item\label{b1:item:29-goal-domain}区分稠密定义算子的对称性与自伴性，能核对谱积分和微分算子的定义域。
\item\label{b1:item:29-goal-group}掌握自伴算子与强连续酉群之间的对应，明确生成元的符号及定义域条件。
\end{itemize}
\end{learninggoals}

本章以\chapref{b1:ch:08}的有界算子、Hilbert 空间投影和伴随为起点。
其中有界逆定理把双射性转为解的稳定性，正交分解则用来判断值域是否稠密。
\secref{b1:sec:2901}只需要这些工具与算子空间的完备性；
读者可以先在对角算子和 \(L^2(0,1)\) 上的乘法算子中理解谱的不同形态。

\ProseParagraphBreak

\secref{b1:sec:2902}先证明本章所需的 Stone--Weierstrass 定理，
再由多项式的一致逼近建立连续函数演算。
\secref{b1:sec:2903}使用\chapref{b1:ch:11}的正测度表示定理，
把循环子空间写成 \(L^2\) 空间；\secref{b1:sec:2904}将这些表示合为唯一的谱测度。
这条路线先完成有界自伴情形，再补出后面 Cayley 变换所需的酉算子谱表示。

\ProseParagraphBreak

\secref{b1:sec:2905}至\secref{b1:sec:2906}把定义域纳入伴随和函数演算。
\secref{b1:sec:2907}的微分算子例子还使用一维 Sobolev 函数的绝对连续代表、
端点迹及分部积分，相关版本在该节开头恢复。
最后的\secref{b1:sec:2908}由谱指数构造酉群，并从群的差商恢复自伴生成算子。
阅读这四节时，应把每个算子的定义域与作用法则一起记下。

\ProseParagraphBreak

关于算子几何级数、可逆集的开性和谱的基本模型，
可对照 Sheldon Axler 的《Measure, Integration \& Real Analysis》\cite[pp.~288--296]{Axler2020Measure}。
谱积分与无界算子的经典论述可对照 Walter Rudin 的《Functional Analysis》
\cite[Sections~12.17--12.21, 13.1--13.12, 13.24, 13.30, 13.35, 13.38]{Rudin1991Functional}；
各节在实际使用处给出更具体的定位。
本章沿用\secref{b1:sec:0805}的 \(T-zI\) 约定。

% !TeX root = ../../Book1.tex
\section{谱与预解算子}
\label{b1:sec:2901}

有限维矩阵的谱由特征值组成，因而求谱可以归结为求解齐次方程。
在无穷维空间中，\(T-zI\) 即使没有非零核，也可能不能把整个空间映满。
本节从非齐次方程
\[
 (T-zI)x=y
\]
出发，把解的存在、唯一性和对数据的连续依赖放在一起考察。
随后用对角算子和乘法算子说明：谱既可能由孤立的特征值及其聚点组成，
也可能充满一个区间，而其中没有任何特征值。

\subsection{谱、预解集与可逆性}
\label{b1:subsec:290101}

\secref{b1:sec:0805}已经对复 Hilbert 空间上的有界算子定义了谱。
这个定义不依赖内积，可以直接用于复 Banach 空间。
以下用 \(\mathcal B(X)=\mathcal B(X,X)\) 表示 \(X\) 上的有界线性算子空间，
用 \(I\) 表示恒等算子；可逆始终指存在定义在整个空间上的有界双侧逆。

\begin{definition}[预解集与预解算子]
\label{b1:def:29-resolvent}
设 \(X\ne\{0\}\) 是复 Banach 空间，\(T\in\mathcal B(X)\)。集合
\[
 \rho(T)=\{z\in\mathbb C:T-zI\text{ 可逆}\}
\]
称为 \(T\) 的\term[yujie-ji]{预解集}[resolvent set]；它的补集
\(\sigma(T)=\mathbb C\setminus\rho(T)\) 是 \(T\) 的谱。
对 \(z\in\rho(T)\)，算子
\[
 R_T(z)=(T-zI)^{-1}
\]
称为 \(T\) 在 \(z\) 处的\term[yujie-suanzi]{预解算子}[resolvent operator]。
\footnote{对同一个算子，\((zI-T)^{-1}=-R_T(z)\)。第二卷的半群理论采用
\((zI-T)^{-1}\) 的写法，因而相应预解恒等式中的系数也要换号；谱和预解集不受影响。}
\symidx{\(\rho(T)\)：有界算子的预解集}
\symidx{\(R_T(z)\)：预解算子}
\end{definition}

由有界逆定理\cref{b1:cor:bounded-inverse}，这里的可逆等价于
\(T-zI:X\to X\) 是双射。把逆的有界性写入定义，
是为了同时记录解的稳定性：对任意 \(y\in X\)，唯一解满足
\[
 x=R_T(z)y,\qquad \|x\|\leq\|R_T(z)\|\cdot\|y\|.
\]
若数据从 \(y\) 改为 \(y+h\)，解的变化至多为
\(\|R_T(z)\|\cdot\|h\|\)。预解算子的范数因而衡量该参数下方程对误差的放大。

\ProseParagraphBreak

单射只保证解至多有一个。例如，在复空间 \(\ell^2\) 上定义右移算子
\[
 S(x_1,x_2,\ldots)=(0,x_1,x_2,\ldots).
\]
它满足 \(\|Sx\|_2=\|x\|_2\)，但值域恰为第一坐标等于零的闭子空间，
所以 \(0\in\sigma(S)\)，尽管 \(Sx=0\) 没有非零解。
即使已知统一下界，也还需要判断值域是否足够大。

\subsection{Neumann 级数与预解恒等式}
\label{b1:subsec:290102}

实数的几何级数可以在算子范数中使用，前提是所在算子空间完备。
\cref{b1:thm:operator-space-complete}保证：若 \(X\) 是 Banach 空间，
则 \(\mathcal B(X)\) 也是 Banach 空间。
下面的算子级数称为\term[neumann-jishu]{Neumann 级数}[Neumann series]。

\begin{lemma}[Neumann 级数]
\label{b1:lem:29-neumann}
设 \(X\ne\{0\}\) 是复 Banach 空间，\(A\in\mathcal B(X)\)，且 \(\|A\|<1\)。
则 \(I-A\) 可逆，且
\[
 (I-A)^{-1}=\sum_{n=0}^{\infty}A^n,\qquad
 \|(I-A)^{-1}\|\leq\frac{1}{1-\|A\|},
\]
其中级数在算子范数中收敛。
\end{lemma}

\begin{proof}
算子范数的乘积估计给出 \(\|A^n\|\leq\|A\|^n\)，
故 \(\sum_{n=0}^{\infty}\|A^n\|<\infty\)。
由算子空间的完备性，部分和 \(B_N=\sum_{n=0}^N A^n\) 收敛于某个
\(B\in\mathcal B(X)\)，并且 \(\|B\|\leq(1-\|A\|)^{-1}\)。
有限和满足
\[
 (I-A)B_N=B_N(I-A)=I-A^{N+1}.
\]
由于 \(\|A^{N+1}\|\to0\)，而与一个固定有界算子相乘在算子范数下连续，
令 \(N\to\infty\) 得 \((I-A)B=B(I-A)=I\)。
\end{proof}

Neumann 级数既能在大参数处求逆，也能从一个已知可逆的参数出发，处理邻近参数。

\begin{proposition}[预解集的开性与局部展开]
\label{b1:prop:29-resolvent-local}
设 \(X\ne\{0\}\) 是复 Banach 空间，\(T\in\mathcal B(X)\)。
若 \(|z|>\|T\|\)，则 \(z\in\rho(T)\)，且
\[
 R_T(z)=-\frac1z\sum_{n=0}^{\infty}\frac{T^n}{z^n},\qquad
 \|R_T(z)\|\leq\frac{1}{|z|-\|T\|}.
\]
若 \(z_0\in\rho(T)\) 且 \(|z-z_0|\cdot\|R_T(z_0)\|<1\)，则
\(z\in\rho(T)\)，并且
\[
 R_T(z)=\sum_{n=0}^{\infty}(z-z_0)^nR_T(z_0)^{n+1}.
\]
两个级数均在算子范数中收敛。因此 \(\rho(T)\) 是开集，
\(\sigma(T)\) 是包含在 \(\{z:|z|\leq\|T\|\}\) 中的闭集。
\end{proposition}

\begin{proof}
对 \(|z|>\|T\|\)，分解
\[
 T-zI=-z\bigl(I-z^{-1}T\bigr)
\]
并把\cref{b1:lem:29-neumann}用于 \(z^{-1}T\)，即得第一个展开及范数估计。
对第二个结论，记 \(R_0=R_T(z_0)\)，则
\[
 T-zI=(T-z_0I)\bigl(I-(z-z_0)R_0\bigr).
\]
右边两个因子都可逆，故
\[
 R_T(z)=\bigl(I-(z-z_0)R_0\bigr)^{-1}R_0
       =\sum_{n=0}^{\infty}(z-z_0)^nR_0^{n+1}.
\]
这还给出
\[
 \|R_T(z)\|\leq
 \frac{\|R_0\|}{1-|z-z_0|\cdot\|R_0\|}.
\]
每个预解点都有一个仍包含在预解集中的开圆盘，故预解集开、谱闭。
第一个展开则给出谱的有界性。
\end{proof}

不同参数处的逆还满足一个代数恒等式。
在下面的约定中，逆取的是 \(T-zI\)，所以右端系数为 \(z-w\)。

\begin{proposition}[预解恒等式]
\label{b1:prop:29-resolvent-identity}
设 \(X\ne\{0\}\) 是复 Banach 空间，\(T\in\mathcal B(X)\)，
\(z,w\in\rho(T)\)。则
\[
 R_T(z)-R_T(w)=(z-w)R_T(z)R_T(w),\qquad
 R_T(z)R_T(w)=R_T(w)R_T(z).
\]
特别地，映射 \(z\mapsto R_T(z)\) 在 \(\rho(T)\) 上按算子范数连续。
\end{proposition}

\begin{proof}
利用双侧逆直接计算，有
\[
 \begin{aligned}
 R_T(z)-R_T(w)
 &=R_T(z)\bigl((T-wI)-(T-zI)\bigr)R_T(w)\\
 &=(z-w)R_T(z)R_T(w).
 \end{aligned}
\]
交换 \(z,w\) 再计算一次，得到同一差等于
\((z-w)R_T(w)R_T(z)\)。当 \(z\ne w\) 时可约去标量因子；
当 \(z=w\) 时交换性直接成立。
固定 \(w\in\rho(T)\)，由\cref{b1:prop:29-resolvent-local}，
在 \(w\) 的充分小邻域内有 \(\|R_T(z)\|\leq2\|R_T(w)\|\)。因此
\[
 \|R_T(z)-R_T(w)\|
 \leq2|z-w|\cdot\|R_T(w)\|^2\longrightarrow0.
\]
\end{proof}

\subsection{自伴算子的谱界}
\label{b1:subsec:290103}

从这一小节起，空间取为复 Hilbert 空间，内积仍对第一变量线性。
对于有界自伴算子，内积不仅限制特征值，还能控制全部谱。
下面先补出从下界估计通向可逆性的步骤。

\begin{lemma}[下界与伴随核]
\label{b1:lem:29-invertibility}
设 \(H\ne\{0\}\) 是复 Hilbert 空间，\(A\in\mathcal B(H)\)。
若存在 \(c>0\)，使 \(\|Ax\|\geq c\|x\|\) 对所有 \(x\in H\) 成立，
并且 \(\ker A^*=\{0\}\)，则 \(A\) 可逆，且 \(\|A^{-1}\|\leq c^{-1}\)。
\end{lemma}

\begin{proof}
下界给出单射性。若 \(Ax_n\to y\)，则
\[
 c\|x_n-x_m\|\leq\|Ax_n-Ax_m\|,
\]
所以 \((x_n)\) 是 Cauchy 列。空间完备，故 \(x_n\to x\) 于某个 \(x\in H\)；
由 \(A\) 连续，\(y=Ax\)。因此 \(\operatorname{ran}A\) 闭。
另一方面，伴随的定义给出
\[
 y\in(\operatorname{ran}A)^\perp
 \quad\Longleftrightarrow\quad
 \langle Ax,y\rangle=0\ (x\in H)
 \quad\Longleftrightarrow\quad A^*y=0.
\]
因而值域的正交补为零。由 Hilbert 空间的投影定理，
一个闭子空间与其正交补之和是整个空间，故 \(\operatorname{ran}A=H\)。
对 \(y=Ax\)，原下界给出 \(\|A^{-1}y\|\leq c^{-1}\|y\|\)。
\end{proof}

\begin{proposition}[非实参数处的预解估计]
\label{b1:prop:29-self-adjoint-real-spectrum}
设 \(H\ne\{0\}\) 是复 Hilbert 空间，\(T\in\mathcal B(H)\) 且 \(T=T^*\)。
若 \(z\in\mathbb C\setminus\mathbb R\)，则 \(z\in\rho(T)\)，且
\[
 \|R_T(z)\|\leq\frac1{|\operatorname{Im}z|}.
\]
因此 \(\sigma(T)\subset\mathbb R\)。
\end{proposition}

\begin{proof}
自伴性保证 \(\langle Tx,x\rangle\) 为实数。由于内积对第一变量线性，
\[
 \operatorname{Im}\langle(T-zI)x,x\rangle
 =-\operatorname{Im}z\cdot\|x\|^2.
\]
取绝对值，再用 Cauchy--Schwarz 不等式，可得
\[
 |\operatorname{Im}z|\cdot\|x\|^2
 \leq\|(T-zI)x\|\cdot\|x\|.
\]
对非零 \(x\) 约去 \(\|x\|\)，得到统一下界
\(\|(T-zI)x\|\geq|\operatorname{Im}z|\cdot\|x\|\)；零向量也满足它。
又有 \((T-zI)^*=T-\overline zI\)，而 \(\overline z\) 的虚部同样非零，
所以上述下界也适用于 \(T-\overline zI\)，其核为零。
\cref{b1:lem:29-invertibility}于是给出可逆性和逆的范数估计。
\end{proof}

还需要排除谱为空的可能，并说明谱界与算子范数的关系。
紧自伴情形的\cref{b1:lem:compact-self-adjoint-extreme}曾利用二次型找到近似特征向量，
再利用紧性取到真正的特征向量。下面只用前一个步骤，因此不要求紧性。

\begin{theorem}[有界自伴算子的谱界]
\label{b1:thm:29-self-adjoint-spectral-bound}
设 \(H\ne\{0\}\) 是复 Hilbert 空间，\(T\in\mathcal B(H)\) 且 \(T=T^*\)。
则 \(\sigma(T)\) 是非空紧集，且
\[
 \sigma(T)\subset[-\|T\|,\|T\|],\qquad
 \max_{t\in\sigma(T)}|t|=\|T\|.
\]
此外，\(T\) 的范数可以由二次型恢复：
\[
 \|T\|=\sup_{\|x\|=1}|\langle Tx,x\rangle|.
\]
\end{theorem}

\begin{proof}
记 \(q(x)=\langle Tx,x\rangle\) 和
\(m=\sup_{\|x\|=1}|q(x)|\)。Cauchy--Schwarz 不等式给出 \(m\leq\|T\|\)。
对任意单位向量 \(x,y\)，自伴性与平行四边形恒等式给出
\[
 \begin{aligned}
 4\operatorname{Re}\langle Tx,y\rangle&=q(x+y)-q(x-y),\\
 |\operatorname{Re}\langle Tx,y\rangle|
 &\leq\frac m4\bigl(\|x+y\|^2+\|x-y\|^2\bigr)=m.
 \end{aligned}
\]
若 \(Tx\ne0\)，取单位向量 \(y=Tx/\|Tx\|\)，得到 \(\|Tx\|\leq m\)；
若 \(Tx=0\)，这个不等式也成立。对单位向量 \(x\) 取上确界，便有 \(\|T\|=m\)。

若 \(T=0\)，则 \(0\in\sigma(T)\)，而所有非零参数都在预解集中，结论成立。
以下设 \(C=\|T\|>0\)。由上确界的定义，可取单位向量列 \((x_n)\)，使
\(|q(x_n)|\to C\)。这些数为实数，取子列后可设
\(q(x_n)\to\lambda\)，其中 \(\lambda=C\) 或 \(\lambda=-C\)。于是
\[
 \begin{aligned}
 \|(T-\lambda I)x_n\|^2
 &=\|Tx_n\|^2-2\lambda\cdot q(x_n)+\lambda^2\\
 &\leq2C^2-2\lambda\cdot q(x_n)\longrightarrow0.
 \end{aligned}
\]
若 \(T-\lambda I\) 有有界逆，则
\[
 1=\|x_n\|\leq\|(T-\lambda I)^{-1}\|\cdot\|(T-\lambda I)x_n\|
 \longrightarrow0,
\]
产生矛盾。因此 \(\lambda\in\sigma(T)\)。
由\cref{b1:prop:29-resolvent-local,b1:prop:29-self-adjoint-real-spectrum}，
谱是闭集且包含在实区间 \([-C,C]\) 内，故紧；刚找到的谱点又满足 \(|\lambda|=C\)。
\end{proof}

谱中的其他点也可以用单位向量上的近似方程来判断。

\begin{corollary}[近似特征向量判别]
\label{b1:cor:29-approximate-eigenvectors}
设 \(H\ne\{0\}\) 是复 Hilbert 空间，\(T\in\mathcal B(H)\) 且 \(T=T^*\)，
\(t\in\mathbb R\)。则 \(t\in\sigma(T)\) 当且仅当存在单位向量列 \((x_n)\)，使
\[
 \|(T-tI)x_n\|\longrightarrow0.
\]
\end{corollary}

\begin{proof}
若存在这样的向量列，有界逆会导出上一证明中的矛盾，故 \(t\in\sigma(T)\)。
反过来，若不存在这样的向量列，则
\(c=\inf_{\|x\|=1}\|(T-tI)x\|>0\)，否则逐个取使范数小于 \(1/n\) 的单位向量即可。
由齐次性，\(\|(T-tI)x\|\geq c\|x\|\) 对所有 \(x\in H\) 成立。
因为 \(t\) 为实数，\(T-tI\) 自伴，其核和伴随的核都为零。
\cref{b1:lem:29-invertibility}给出 \(t\in\rho(T)\)。
\end{proof}

上述单位向量列称为 \(t\) 的\term[jinsi-tezheng-xiangliang]{近似特征向量列}[approximate eigenvector sequence]。
它不必有范数收敛子列，因此存在这种向量列并不保证存在真正的特征向量。

\subsection{离散谱与连续谱的模型}
\label{b1:subsec:290104}

为了比较不同的失去可逆性的方式，先明确两个名称。

\begin{definition}[点谱与连续谱]
\label{b1:def:29-point-continuous-spectrum}
设 \(X\ne\{0\}\) 是复 Banach 空间，\(T\in\mathcal B(X)\)，\(z\in\mathbb C\)。
若 \(\ker(T-zI)\ne\{0\}\)，则称 \(z\) 属于 \(T\) 的\term[dian-pu]{点谱}[point spectrum]，
记作 \(z\in\sigma_{\mathrm p}(T)\)。
若 \(T-zI\) 单射，值域在 \(X\) 中稠密但不等于 \(X\)，
则称 \(z\) 属于 \(T\) 的\term[lianxu-pu]{连续谱}[continuous spectrum]，
记作 \(z\in\sigma_{\mathrm c}(T)\)。
\symidx{\(\sigma_{\mathrm p}(T)\)：点谱}
\symidx{\(\sigma_{\mathrm c}(T)\)：连续谱}
\end{definition}

点谱就是特征值的集合。一般有界算子还可能出现单射而值域不稠密的情形，
前面的右移算子在参数零处就是例子。
对于有界自伴算子和实数 \(t\)，则有
\[
 \bigl(\operatorname{ran}(T-tI)\bigr)^\perp=\ker(T-tI).
\]
所以只要 \(t\) 不是特征值，值域就稠密；若它仍是谱点，就属于连续谱。
结合谱为实集的结论，有界自伴算子的谱由点谱和连续谱组成。

\begin{definition}[离散谱]
\label{b1:def:29-discrete-spectrum}
设 \(H\ne\{0\}\) 是复 Hilbert 空间，\(T\in\mathcal B(H)\) 且 \(T=T^*\)。
若 \(t\in\sigma(T)\) 是 \(\sigma(T)\) 的孤立点，且
\(0<\dim\ker(T-tI)<\infty\)，则称 \(t\) 是离散谱点。
这些点组成 \(T\) 的\term[lisan-pu]{离散谱}[discrete spectrum]。
\end{definition}

这个定义同时要求孤立性和有限重数，不能把任意特征值都称为离散谱点。
这里只对自伴算子使用这一名称。

\begin{proposition}[趋于零的对角算子]
\label{b1:examp:29-compact-diagonal-spectrum}
在复 Hilbert 空间 \(\ell^2\) 上，令
\[
 D(x_1,x_2,\ldots)=(x_1,\tfrac12x_2,\tfrac13x_3,\ldots).
\]
则 \(D\) 紧且自伴，并且
\[
 \sigma(D)=\{0\}\cup\{1/n:n\geq1\},\qquad
 \sigma_{\mathrm p}(D)=\{1/n:n\geq1\},\qquad
 \sigma_{\mathrm c}(D)=\{0\}.
\]
每个非零谱点都属于离散谱。
\end{proposition}

\begin{proof}
\secref{b1:sec:0805}已经证明它紧且自伴。
每个 \(1/n\) 都是特征值，相应特征空间为 \(\operatorname{span}\{e_n\}\)。
若 \(z\notin K=\{0\}\cup\{1/n:n\geq1\}\)，则 \(K\) 紧，
故 \(\delta=\operatorname{dist}(z,K)>0\)。此时
\[
 R_D(z)y=\left(\frac{y_n}{1/n-z}\right)_{n\geq1},\qquad
 \|R_D(z)y\|_2\leq\delta^{-1}\|y\|_2,
\]
直接代入可验证它是双侧逆。另一方面，\(D\) 单射，值域包含全部有限支撑序列，
因而稠密；但 \(y=(1/n)_{n\geq1}\in\ell^2\) 不在值域内，因为它的唯一可能原像
\((1,1,\ldots)\) 不属于 \(\ell^2\)。所以零属于连续谱，且没有 \(K\) 以外的谱点。
每个 \(1/n\) 都是孤立的单重特征值，故属于离散谱。
\end{proof}

谱点零虽然不是特征值，却是这些离散谱点的聚点；
单位向量 \(e_n\) 满足 \(\|De_n\|_2=1/n\to0\)，给出了它的近似特征向量列。

\begin{proposition}[区间上的乘法算子]
\label{b1:examp:29-multiplication-spectrum}
在复空间 \(H=L^2((0,1),\mathrm{d}t)\) 上令 \((Mf)(t)=t\cdot f(t)\)。
则 \(M\) 有界且自伴，并且
\[
 \sigma(M)=\sigma_{\mathrm c}(M)=[0,1],\qquad
 \sigma_{\mathrm p}(M)=\varnothing.
\]
\end{proposition}

\begin{proof}
由 \(0<t<1\) 得 \(\|Mf\|_2\leq\|f\|_2\)，且
\(\langle Mf,g\rangle=\langle f,Mg\rangle\)，所以 \(M\) 有界且自伴。
若 \(z\notin[0,1]\)，则 \(\delta=\operatorname{dist}(z,[0,1])>0\)，而
\[
 (R_M(z)g)(t)=\frac{g(t)}{t-z},\qquad
 \|R_M(z)g\|_2\leq\delta^{-1}\|g\|_2
\]
给出有界双侧逆。

固定 \(\lambda\in[0,1]\)。对 \(0<\varepsilon<1\)，记
\[
 E_{\lambda,\varepsilon}=(0,1)\cap(\lambda-\varepsilon,\lambda+\varepsilon),
 \qquad
 f_\varepsilon=|E_{\lambda,\varepsilon}|^{-1/2}\cdot\mathbf1_{E_{\lambda,\varepsilon}},
\]
其中 \(|E|\) 表示 Lebesgue 测度。这个集合即使在 \(\lambda=0,1\) 时也有正测度，且
\[
 \|f_\varepsilon\|_2=1,\qquad
 \|(M-\lambda I)f_\varepsilon\|_2^2
 =\frac1{|E_{\lambda,\varepsilon}|}
   \int_{E_{\lambda,\varepsilon}}|t-\lambda|^2\,\mathrm{d}t
 \leq\varepsilon^2.
\]
取 \(\varepsilon\downarrow0\)，由\cref{b1:cor:29-approximate-eigenvectors}可知
\(\lambda\in\sigma(M)\)。

若 \((M-\lambda I)f=0\)，则 \(f\) 只能在零测集 \(\{t=\lambda\}\) 上非零，
故 \(f=0\) 于 \(L^2\)。因而 \(M-\lambda I\) 单射，且由自伴性知值域稠密。
它却不满射：常函数 \(1\) 的原像只能是 \((t-\lambda)^{-1}\)，
而此函数的平方在 \(\lambda\) 附近不可积，在两个端点处的单侧积分也发散。
故区间内每个点都属于连续谱，而区间外每个点都属于预解集。
\end{proof}

这两个模型都可以看作乘法：一个在坐标序列上乘以 \(1/n\)，另一个在函数空间中乘以 \(t\)。
后续的谱表示将把一般自伴算子化为适当测度空间上的乘法算子，
从而以同一种运算处理特征向量展开和连续谱。

% !TeX root = ../../Book1.tex
\section{Stone–Weierstrass 定理与连续函数演算}
\label{b1:sec:2902}

若 \(T\) 是矩阵，多项式 \(p(T)\) 由有限次加法和乘法定义。
对于连续函数 \(f\)，要赋予 \(f(T)\) 意义，就需要一种与算子范数相容的逼近。
本节先证明连续函数的代数逼近定理，再对有界自伴算子建立
\(\|p(T)\|=\max_{\sigma(T)}|p|\)。这个等式既保证逼近所得的极限存在，
也保证极限不依赖所选多项式。

\subsection{连续函数代数与分离点条件}
\label{b1:subsec:290201}

\begin{definition}[函数代数与分离点]
\label{b1:def:29-function-algebra}
设 \(K\ne\varnothing\) 是紧 Hausdorff 空间，\(\mathbb K=\mathbb R\) 或 \(\mathbb C\)。
若 \(\mathcal A\subset C(K;\mathbb K)\) 是线性子空间，且对逐点乘法封闭，
则称 \(\mathcal A\) 为\term[hanshu-daishu]{函数代数}[function algebra]。
若对任意不同的 \(x,y\in K\)，都有 \(f\in\mathcal A\) 使 \(f(x)\ne f(y)\)，
则称 \(\mathcal A\) \term[fenli-dian]{分离点}[separates points]。
对于复函数代数，若 \(f\in\mathcal A\) 蕴含 \(\overline f\in\mathcal A\)，
则称它对复共轭封闭。
\end{definition}

这里的拓扑取一致范数 \(\|f\|_\infty=\max_{x\in K}|f(x)|\) 所给出的拓扑。
含常数函数的代数可以调整一个函数的常数项和比例；配合分离点条件，
便能同时指定任意两个不同点处的函数值。
但仅有逐点取值的自由仍不够，下一步还要把这些两点插值拼成全空间的一致逼近。

\ProseParagraphBreak

例如，实直线紧子集上的多项式限制组成含常数且分离点的代数。
相反，在 \([-1,1]\) 上只取偶多项式，所得函数在 \(x\) 与 \(-x\) 处总相等；
它们不能一致逼近函数 \(x\mapsto x\)，因为在两个端点中至少有一处误差不小于 \(1\)。
分离点条件排除了这种无法区分位置的限制。

\subsection{Stone–Weierstrass 逼近}
\label{b1:subsec:290202}

证明中需要把两个近似函数取最大值或最小值。
下面先说明，一致闭代数中的绝对值运算可以从多项式运算得到。

\begin{lemma}[绝对值的多项式逼近]
\label{b1:lem:29-absolute-polynomial}
对每个 \(M>0\)，函数 \(s\mapsto|s|\) 在 \([-M,M]\) 上是实多项式的一致极限。
因此，若 \(K\) 是紧 Hausdorff 空间，\(\mathcal A\subset C(K;\mathbb R)\)
是含常数函数的一致闭代数，则 \(f\in\mathcal A\) 蕴含 \(|f|\in\mathcal A\)。
\end{lemma}
\begin{proof}
在 \([0,1]\) 上递归定义多项式
\[
 p_0(u)=0,\qquad p_{n+1}(u)=p_n(u)+\frac{u-p_n(u)^2}{2}.
\]
若 \(0\leq p_n(u)\leq\sqrt u\)，则
\[
 \sqrt u-p_{n+1}(u)
 =\bigl(\sqrt u-p_n(u)\bigr)
  \left(1-\frac{\sqrt u+p_n(u)}2\right)\geq0,
\]
且 \(p_{n+1}(u)\geq p_n(u)\)。由归纳，所有这些不等式成立。
给定 \(0<\delta<1\)，当 \(u\leq\delta^2\) 时误差不超过 \(\delta\)；
当 \(u\geq\delta^2\) 时，每步误差至多乘以 \(1-\delta/2\)，故
\[
 0\leq\sqrt u-p_n(u)\leq
 \max\{\delta,(1-\delta/2)^n\}\qquad(0\leq u\leq1).
\]
先令 \(n\to\infty\)，再令 \(\delta\downarrow0\)，得到一致收敛。
于是 \(Mp_n(s^2/M^2)\) 在 \([-M,M]\) 上一致趋于 \(|s|\)。
对 \(f\in\mathcal A\)，取 \(M>\|f\|_\infty\)，这些多项式复合 \(f\) 都属于 \(\mathcal A\)，
再用闭性即得 \(|f|\in\mathcal A\)。
\end{proof}

因而上述实闭代数也对
\[
 \max(f,g)=\frac{f+g+|f-g|}{2},\qquad
 \min(f,g)=\frac{f+g-|f-g|}{2}
\]
封闭。这使有限个局部逼近可以保持单侧误差地合在一起。
以下结果称为\term[stone-weierstrass-dingli]{Stone--Weierstrass 定理}[Stone--Weierstrass theorem]。

\begin{theorem}[Stone--Weierstrass 定理]
\label{b1:thm:29-stone-weierstrass}
设 \(K\ne\varnothing\) 是紧 Hausdorff 空间。
若实函数代数 \(\mathcal A\subset C(K;\mathbb R)\) 含常数函数并分离点，
则它在一致范数下稠密。
若复函数代数 \(\mathcal A\subset C(K;\mathbb C)\) 含常数函数、分离点且对复共轭封闭，
则它在 \(C(K;\mathbb C)\) 中一致稠密。
\end{theorem}
\begin{proof}
先考虑实情形。代数的闭包仍是代数，因为一致收敛保持加法与乘法；
所以可把 \(\mathcal A\) 换成闭包，证明这个闭包含任意 \(h\in C(K;\mathbb R)\)。
固定 \(\varepsilon>0\)。对每对 \(x,y\in K\)，存在 \(f_{x,y}\in\mathcal A\)，使
\(f_{x,y}(x)=h(x)\)、\(f_{x,y}(y)=h(y)\)：当 \(x\ne y\) 时，
对一个分离它们的函数作仿射变换；当 \(x=y\) 时取常数 \(h(x)\)。

固定 \(x\)。开集
\(\{z:f_{x,y}(z)>h(z)-\varepsilon\}\) 随 \(y\) 变化覆盖 \(K\)。
取有限子覆盖，对对应的 \(f_{x,y}\) 取最大值，得到 \(g_x\in\mathcal A\)，满足
\[
 g_x>h-\varepsilon\ \text{于 }K,\qquad g_x(x)=h(x).
\]
现在开集 \(\{z:g_x(z)<h(z)+\varepsilon\}\) 随 \(x\) 变化又覆盖 \(K\)。
取有限子覆盖，对相应 \(g_x\) 取最小值，得到 \(g\in\mathcal A\)，使
\(h-\varepsilon<g<h+\varepsilon\) 于 \(K\)。因此 \(h\) 属于原代数的闭包。

复情形中，\(\mathcal A\) 内实值函数组成实代数。
若 \(f(x)\ne f(y)\)，则 \(\operatorname{Re}f\) 或 \(\operatorname{Im}f\) 分离这两点；
复共轭封闭性保证它们仍属于 \(\mathcal A\)。实情形可分别逼近任意连续复函数的实部和虚部，
把两个近似相加便得到复情形。
\end{proof}

紧性在证明中使用了两次：先把固定一点的两点插值合成一个全局下界，
再把不同固定点所给的上界合在一起。定理不要求空间可度量。

\subsection{多项式函数演算}
\label{b1:subsec:290203}

以下回到非零复 Hilbert 空间 \(H\) 上的有界自伴算子 \(T\)，
记 \(K=\sigma(T)\subset\mathbb R\)。
对复多项式 \(p(t)=\sum_{j=0}^m a_jt^j\)，定义
\[
 p(T)=\sum_{j=0}^m a_jT^j.
\]
有限和的运算给出 \((pq)(T)=p(T)q(T)\) 和
\(p(T)^*=\overline p(T)\)，其中 \(\overline p(t)=\sum_j\overline{a_j}t^j\)。
这种把函数送到算子的对应称为\term[hanshu-yansuan]{函数演算}[functional calculus]。
下面的谱计算只需要代数基本定理对一元复多项式的因式分解。

\begin{proposition}[多项式谱映射与范数]
\label{b1:prop:29-polynomial-calculus}
设 \(H\ne\{0\}\) 是复 Hilbert 空间，\(T\in\mathcal B(H)\) 且 \(T=T^*\)。
对任意复多项式 \(p\)，有
\[
 \sigma(p(T))=p(\sigma(T)),\qquad
 \|p(T)\|=\max_{t\in\sigma(T)}|p(t)|.
\]
\end{proposition}
\begin{proof}
先证谱等式。常多项式的情形由标量算子的逆直接得到。
若 \(p\) 非常数，对每个 \(w\in\mathbb C\)，由代数基本定理可写
\[
 p(t)-w=a\prod_{j=1}^m(t-z_j),\qquad a\ne0.
\]
于是 \(p(T)-wI=a\prod_j(T-z_jI)\)。这些因子相互交换。
有限个相互交换的有界算子的乘积可逆，当且仅当每个因子可逆。
为看清反向，若 \(AB=BA\) 且 \(AB\) 可逆，则 \(A\) 与 \((AB)^{-1}\) 交换，
故 \(B(AB)^{-1}\) 是 \(A\) 的双侧逆；对有限个因子重复这个论证即可。
因此 \(w\notin\sigma(p(T))\) 等价于所有 \(z_j\notin\sigma(T)\)，即 \(w\notin p(\sigma(T))\)。

对任意 \(B\in\mathcal B(H)\)，伴随的范数等式与 Cauchy--Schwarz 不等式给出
\[
 \|B\|^2=\sup_{\|x\|=1}\langle B^*Bx,x\rangle
 \leq\|B^*B\|\leq\|B^*\|\cdot\|B\|=\|B\|^2.
\]
取 \(B=p(T)\)，令实系数多项式 \(q(t)=\overline p(t)p(t)\)。
算子 \(q(T)=p(T)^*p(T)\) 自伴，由\cref{b1:thm:29-self-adjoint-spectral-bound}和刚证的谱映射等式，
\[
 \|p(T)\|^2=\|q(T)\|
 =\max_{s\in\sigma(q(T))}|s|
 =\max_{t\in\sigma(T)}|p(t)|^2.
\]
\end{proof}

这说明即使多项式在谱以外变化很大，它作用在 \(T\) 上的范数仍只由谱上的值决定。
若两个多项式在谱上相同，它们也给出同一个算子。

\subsection{连续函数演算的构造与唯一性}
\label{b1:subsec:290204}

\begin{theorem}[连续函数演算]
\label{b1:thm:29-continuous-calculus}
设 \(H\ne\{0\}\) 是复 Hilbert 空间，\(T\in\mathcal B(H)\) 且 \(T=T^*\)，
\(K=\sigma(T)\)。存在唯一的连续线性映射
\[
 C(K;\mathbb C)\longrightarrow\mathcal B(H),\qquad f\longmapsto f(T),
\]
它保持乘法，把常函数 \(1\) 送到 \(I\)，把坐标函数 \(t\mapsto t\) 送到 \(T\)。
该映射还满足
\[
 \overline f(T)=f(T)^*,\qquad
 \|f(T)\|=\max_K|f|,\qquad
 \sigma(f(T))=f(K).
\]
若 \(f\geq0\) 于 \(K\)，则 \(\langle f(T)x,x\rangle\geq0\) 对所有 \(x\in H\) 成立。
\end{theorem}
\begin{proof}
多项式在实紧集 \(K\) 上分离点、含常数且对共轭封闭，
故由\cref{b1:thm:29-stone-weierstrass}在 \(C(K;\mathbb C)\) 中稠密。
给定 \(f\)，选多项式 \(p_n\to f\) 一致收敛。
\cref{b1:prop:29-polynomial-calculus}说明 \((p_n(T))\) 在算子范数中是 Cauchy 列，
所以可定义 \(f(T)=\lim_n p_n(T)\)。若换另一列 \(q_n\to f\)，则
\(\|p_n(T)-q_n(T)\|=\|p_n-q_n\|_{C(K)}\to0\)，故定义与选择无关。
有限和、乘积和伴随在算子范数下连续，所以多项式的相应性质都传到极限。
同一范数等式还给出 \(\|f(T)\|=\|f\|_{C(K)}\)。
任一满足题设的连续映射在多项式上必须取值为 \(p(T)\)，由稠密性知其唯一。

若 \(w\notin f(K)\)，函数 \(g=(f-w)^{-1}\) 在 \(K\) 上连续；
乘法性质说明 \(g(T)\) 是 \(f(T)-wI\) 的双侧逆。
反之，取 \(t\in K\) 和单位向量列 \(x_n\)，使 \((T-tI)x_n\to0\)。
对固定多项式 \(p\)，因 \(p(s)-p(t)=(s-t)q(s)\)，有
\((p(T)-p(t)I)x_n\to0\)。再选 \(p\) 使 \(\|f-p\|_{C(K)}<\varepsilon\)，得到
\[
 \limsup_n\|(f(T)-f(t)I)x_n\|\leq2\varepsilon.
\]
令 \(\varepsilon\downarrow0\)，便得近似特征向量列；有界逆不可能存在，
故 \(f(t)\in\sigma(f(T))\)。最后，若 \(f\geq0\)，令 \(g=\sqrt f\in C(K)\)，则
\(f(T)=g(T)^*g(T)\)，其二次型非负。
\end{proof}

\begin{definition}[正算子]
\label{b1:def:29-positive-operator}
设 \(H\) 是复 Hilbert 空间，\(B\in\mathcal B(H)\) 且 \(B=B^*\)。
若 \(\langle Bx,x\rangle\geq0\) 对所有 \(x\in H\) 成立，
则称 \(B\) 为\term[zheng-suanzi]{正算子}[positive operator]，记作 \(B\geq0\)。
这里允许 \(B\) 有非零核；正性不表示存在严格正的统一下界。
\end{definition}

\begin{corollary}[正平方根与复合]
\label{b1:cor:29-positive-square-root}
设 \(H\ne\{0\}\) 是复 Hilbert 空间，\(T\in\mathcal B(H)\) 且 \(T=T^*\)。
若 \(f\in C(\sigma(T);\mathbb R)\)、\(g\in C(f(\sigma(T));\mathbb C)\)，则
\[
 g(f(T))=(g\circ f)(T).
\]
此外，\(T\geq0\) 当且仅当 \(\sigma(T)\subset[0,\infty)\)；
每个有界正算子 \(T\) 有唯一的有界正平方根 \(T^{1/2}\)，满足 \((T^{1/2})^2=T\)。
\end{corollary}
\begin{proof}
复合等式对多项式 \(g\) 由乘法性质成立；
在实紧集 \(f(\sigma(T))=\sigma(f(T))\) 上一致逼近一般 \(g\)，即得所述等式。
若谱包含在非负半轴，连续演算的正性给出 \(T\geq0\)。
反之，若 \(T\geq0\)，对任意 \(t\in\sigma(T)\) 取上一节的单位近似特征向量，
便有 \(0\leq\langle Tx_n,x_n\rangle\to t\)，所以 \(t\geq0\)。
存在性由 \(t\mapsto\sqrt t\) 的连续演算得到。
若 \(B\geq0\) 且 \(B^2=T\)，复合等式给出
\(T^{1/2}=(B^2)^{1/2}=B\)，从而平方根唯一。
\end{proof}

最后记录一个以后会用到的交换性质。若 \(S\in\mathcal B(H)\) 且 \(ST=TS\)，
则 \(S\) 与每个多项式 \(p(T)\) 交换；对多项式逼近取算子范数极限，
可得 \(Sf(T)=f(T)S\) 对所有 \(f\in C(\sigma(T))\) 成立。
下一节将把这个演算与积分联系起来。

% !TeX root = ../../Book1.tex
\section{有界自伴算子的谱表示}
\label{b1:sec:2903}

连续函数演算把谱上的函数变成算子。若再固定一个向量，取算子的二次型，
就得到连续函数空间上的正泛函；\chapref{b1:ch:11}的表示定理把它化为积分。
本节沿着这一步构造一个 \(L^2\) 空间，并把原算子在相应子空间上的作用化为乘以坐标函数。
一般 Hilbert 空间可能不可分，因而最后得到的是一族这样的空间的正交直和。

\subsection{正泛函与标量谱测度}
\label{b1:subsec:290301}

\begin{proposition}[向量所确定的标量测度]
\label{b1:prop:29-scalar-spectral-measure}
设 \(H\ne\{0\}\) 是复 Hilbert 空间，\(T\in\mathcal B(H)\) 且 \(T=T^*\)，
\(K=\sigma(T)\)。对每个 \(v\in H\)，存在唯一的有限正 Radon 测度 \(\mu_v\)，使
\[
 \langle f(T)v,v\rangle=\int_K f(t)\,\mathrm{d}\mu_v(t)
 \qquad(f\in C(K;\mathbb C)).
\]
它满足 \(\mu_v(K)=\|v\|^2\)，并且
\[
 \|f(T)v\|^2=\int_K|f|^2\,\mathrm{d}\mu_v.
\]
\end{proposition}
\begin{proof}
映射 \(f\mapsto\langle f(T)v,v\rangle\) 线性，且在 \(f\geq0\) 时非负，
因为连续演算保持正性。紧空间上 \(C_c(K)=C(K)\)，
所以\cref{b1:thm:riesz-markov-positive}给出唯一的正 Radon 表示测度。
取 \(f=1\) 得总质量有限且等于 \(\|v\|^2\)。又因为
\(f(T)^*f(T)=|f|^2(T)\)，有
\[
 \|f(T)v\|^2=\langle |f|^2(T)v,v\rangle
 =\int_K|f|^2\,\mathrm{d}\mu_v.
\]
\end{proof}

测度 \(\mu_v\) 称为向量 \(v\) 关于 \(T\) 的\term[biaoliang-pu-cedu]{标量谱测度}[scalar spectral measure]。
\symidx{\(\mu_v\)：向量的标量谱测度}
它记录连续演算对这个向量的作用强度。
若 \(v\) 是单位特征向量且 \(Tv=\lambda v\)，则先对多项式、再取一致极限可得
\(f(T)v=f(\lambda)v\)，所以 \(\mu_v=\delta_\lambda\)。
非特征向量的测度则可以分布在多个点上，也可以没有原子。

\subsection{循环子空间}
\label{b1:subsec:290302}

\begin{definition}[循环子空间]
\label{b1:def:29-cyclic-subspace}
设 \(H\ne\{0\}\) 是复 Hilbert 空间，\(T\in\mathcal B(H)\) 且 \(T=T^*\)，\(v\in H\)。
闭子空间
\[
 H_v=\overline{\{p(T)v:p\text{ 是复多项式}\}}
\]
称为 \(v\) 生成的\term[xunhuan-zikongjian]{循环子空间}[cyclic subspace]。
若 \(H_v=H\)，则称 \(v\) 为 \(T\) 的\term[xunhuan-xiangliang]{循环向量}[cyclic vector]。
\end{definition}

多项式的和仍是多项式，所以闭包内的集合确实是线性子空间。
它在 \(T\) 下不变；由自伴性，\(H_v^\perp\) 也在 \(T\) 下不变。
一般地，对复 Hilbert 空间上的有界算子，若一个闭子空间及其正交补都在该算子下不变，就称它为该算子的
\term[yuehua-zikongjian]{约化子空间}[reducing subspace]。
在这种子空间上限制有界自伴算子，仍得到自伴算子。

\begin{definition}[酉算子与酉等价]
\label{b1:def:29-unitary-equivalence}
设 \(H_1,H_2\) 是复 Hilbert 空间。若线性映射 \(U:H_1\to H_2\)
满射且满足 \(\|Ux\|=\|x\|\)，则称 \(U\) 为\term[you-suanzi]{酉算子}[unitary operator]。
若 \(T_j\in\mathcal B(H_j)\)，且存在这样的 \(U\) 使 \(UT_1=T_2U\)，
则称 \(T_1,T_2\) \term[you-dengjia]{酉等价}[unitarily equivalent]。
\end{definition}

由极化恒等式，酉算子保持内积，因而 \(U^{-1}=U^*\)。
反过来，\(U^*U=I\) 和 \(UU^*=I\) 分别给出等距性与满射性。

\begin{theorem}[循环子空间上的乘法表示]
\label{b1:thm:29-cyclic-representation}
设 \(H\ne\{0\}\) 是复 Hilbert 空间，\(T\in\mathcal B(H)\) 且 \(T=T^*\)，
\(v\ne0\)，\(K=\sigma(T)\)，\(\mu_v\) 如\cref{b1:prop:29-scalar-spectral-measure}。
映射 \(f\mapsto f(T)v\) 从 \(C(K)\) 唯一延拓为酉算子
\[
 W_v:L^2(K,\mu_v)\longrightarrow H_v.
\]
它满足 \(W_v1=v\)，以及
\[
 W_v M_t=(T|_{H_v})W_v,
 \qquad (M_t f)(t)=t\cdot f(t).
\]
\end{theorem}
\begin{proof}
上一命题说明 \(\|f(T)v\|=\|f\|_{L^2(\mu_v)}\)。
因此两个连续函数若代表同一 \(L^2\) 等价类，就给出相同的向量。
由\cref{b1:prop:radon-cc-dense-lp}，\(C(K)\) 在 \(L^2(K,\mu_v)\) 中稠密，
等距映射因而唯一延拓到全空间；具体地，对 \(f_n\to f\) 于 \(L^2\)，
定义 \(W_vf=\lim_n f_n(T)v\)，等距性保证极限存在且与逼近无关。
等距映射的像闭，而其像包含全部 \(p(T)v\)，故像恰为 \(H_v\)。
对 \(f\in C(K)\)，乘法性质给出
\[
 W_v(tf)=(tf)(T)v=T f(T)v=T W_vf.
\]
坐标乘法 \(M_t\) 在 \(L^2(K,\mu_v)\) 上有界，故由稠密性，这个等式延伸到所有 \(f\)。
\end{proof}

这里的 \(L^2\) 范数恰好来自原空间中的范数，测度不是任意添入的。
例如，对\cref{b1:examp:29-multiplication-spectrum}中的 \(Mf(t)=t\cdot f(t)\)，
向量 \(1\) 的谱测度就是 \([0,1]\) 上的 Lebesgue 测度。
多项式在连续函数中一致稠密，连续函数又在 \(L^2\) 中稠密，所以 \(1\) 是循环向量。
这个模型有循环向量，却没有任何特征向量。

\subsection{正交循环分解}
\label{b1:subsec:290303}

一个向量生成的子空间未必等于整个空间。为同时容纳多个这样的表示，
给定一族复 Hilbert 空间 \((H_\alpha)_{\alpha\in J}\)，定义它们的
\term[hilbert-zhengjiao-zhihe]{Hilbert 正交直和}[Hilbert direct sum] 为
\[
 \bigoplus_{\alpha\in J}H_\alpha
 =\left\{(x_\alpha):x_\alpha\in H_\alpha,
   \ \sum_{\alpha\in J}\|x_\alpha\|^2<\infty\right\},
 \qquad
 \|(x_\alpha)\|^2=\sum_\alpha\|x_\alpha\|^2.
\]
非负和按有限部分和的上确界理解。
对每个这样的向量，非零分量至多可数：对每个正整数 \(n\)，
满足 \(\|x_\alpha\|\geq1/n\) 的指标只有有限个。
内积由 \(\sum_\alpha\langle x_\alpha,y_\alpha\rangle\) 给出，
其绝对收敛由 Cauchy--Schwarz 不等式保证。

\ProseParagraphBreak

这个空间完备。事实上，若 \((x^{(n)})\) 是 Cauchy 列，每个分量都收敛于某个 \(x_\alpha\)。
对任意有限指标集取平方和并令 \(n\to\infty\)，可知 \(\sum_\alpha\|x_\alpha\|^2<\infty\)。
固定充分大的 \(n\)，再在
\(\sum_{\alpha\in F}\|x_\alpha^{(n)}-x_\alpha^{(m)}\|^2\leq\varepsilon^2\)
中令 \(m\to\infty\)，并对有限集 \(F\) 取上确界，便有 \(\|x^{(n)}-x\|\leq\varepsilon\)。

\begin{theorem}[正交循环分解]
\label{b1:thm:29-cyclic-decomposition}
设 \(H\ne\{0\}\) 是复 Hilbert 空间，\(T\in\mathcal B(H)\) 且 \(T=T^*\)。
存在单位向量族 \((v_\alpha)_{\alpha\in J}\)，使相应循环子空间两两正交，且
\[
 H=\bigoplus_{\alpha\in J}H_{v_\alpha}.
\]
若 \(H\) 可分，可以取 \(J\) 为有限集或可数集。
\end{theorem}
\begin{proof}
考虑所有生成两两正交循环子空间的单位向量集合，按包含关系排序。
一条链的并仍有同样性质，故由 Zorn 引理存在极大集合 \(\{v_\alpha\}\)。
令 \(M\) 为这些子空间的闭线性张成空间。
每个子空间约化 \(T\)，于是 \(M\) 和 \(M^\perp\) 都在 \(T\) 下不变。
若 \(M^\perp\ne\{0\}\)，在其中取一个单位向量 \(v\)，
则整个 \(H_v\subset M^\perp\)，可以把 \(v\) 添入原集合，与极大性矛盾。
因此 \(M=H\)。

对于 \((x_\alpha)\in\bigoplus H_{v_\alpha}\)，正交性使有限部分和满足 Cauchy 估计，
其非零项至多可数，故这些和在 \(H\) 中收敛，且平方范数等于各分量平方范数之和。
所得映射等距，其像闭且含全部有限和，由 \(M=H\) 知它满射。
这便给出所述直和。若 \(H\) 可分，任何标准正交系至多可数：
取一个可数稠密集，各单位向量周围半径 \(1/3\) 的球两两不交，每球包含不同的稠密集点。
应用于 \((v_\alpha)\) 即得结论。
\end{proof}

\subsection{乘法表示与谱表示的唯一性}
\label{b1:subsec:290304}

\begin{theorem}[有界自伴算子的乘法表示]
\label{b1:thm:29-bounded-spectral-representation}
设 \(H\ne\{0\}\) 是复 Hilbert 空间，\(T\in\mathcal B(H)\) 且 \(T=T^*\)，
\(K=\sigma(T)\)。存在 \(K\) 上的一族有限正 Radon 测度 \((\mu_\alpha)_{\alpha\in J}\)
及酉算子
\[
 U:H\longrightarrow\bigoplus_{\alpha\in J}L^2(K,\mu_\alpha),
\]
使
\[
 UTU^{-1}=\bigoplus_{\alpha\in J}M_t,\qquad
 Uf(T)U^{-1}=\bigoplus_{\alpha\in J}M_f\quad(f\in C(K)).
\]
各 \(\mu_\alpha\) 可以取为概率测度。
\end{theorem}
\begin{proof}
取\cref{b1:thm:29-cyclic-decomposition}中的单位循环向量族，令 \(\mu_\alpha=\mu_{v_\alpha}\)，
并在每个分量上用\cref{b1:thm:29-cyclic-representation}构造 \(W_{v_\alpha}\)。
这些酉映射的正交直和把 \(\bigoplus_\alpha L^2(K,\mu_\alpha)\) 映到 \(H\)，取其逆为 \(U\)。
每个分量上的交织关系给出 \(UTU^{-1}=\bigoplus M_t\)。
这个关系先传到多项式，再利用一致逼近和
\(\|M_f\|\leq\|f\|_{C(K)}\) 传到一般连续函数。
由 \(\|v_\alpha\|=1\)，各谱测度的总质量为 \(1\)。
\end{proof}

这里的循环分解通常不唯一。即使 \(T=\lambda I\)，也可以选取不同的标准正交基，
把空间分解成不同的一维循环子空间。唯一性应针对由原算子决定的对象来表述。

\begin{proposition}[循环表示的唯一性]
\label{b1:prop:29-cyclic-uniqueness}
设 \(H\ne\{0\}\) 是复 Hilbert 空间，\(T\in\mathcal B(H)\) 且 \(T=T^*\)，
\(v\) 是循环向量，\(K=\sigma(T)\)。若 \(\nu\) 是 \(K\) 上的有限正 Radon 测度，
\(W:L^2(K,\nu)\to H\) 是酉算子，且 \(W1=v\)、\(WM_t=TW\)，
则 \(\nu=\mu_v\)，并且 \(W\) 就是\cref{b1:thm:29-cyclic-representation}构造的映射。
\end{proposition}
\begin{proof}
交织关系给出 \(W p=p(T)v\)，先对多项式成立，再由一致逼近对所有 \(f\in C(K)\) 成立。
于是
\[
 \int_K f\,\mathrm{d}\nu
 =\langle f,1\rangle_{L^2(\nu)}
 =\langle f(T)v,v\rangle
 =\int_K f\,\mathrm{d}\mu_v.
\]
由表示测度的唯一性，\(\nu=\mu_v\)。连续函数在这个 \(L^2\) 空间中稠密，
而两个映射在连续函数上相等，所以在整个空间上相等。
\end{proof}

对于全空间，下一节会把集合的示性函数也纳入演算，得到唯一的投影值谱测度。
这份测度与选择哪一组循环向量无关。

% !TeX root = ../../Book1.tex
\section{投影值谱测度与 Borel 函数演算}
\label{b1:sec:2904}

在乘法表示中，选取谱的一部分只需乘以该集合的示性函数。
把这个操作搬回原空间，就得到一个正交投影。
本节先构造这些投影，再把简单函数的积分推广到有界 Borel 函数。
所得谱测度是唯一的；上一节为构造它而选择的循环分解将不再出现在最终定义中。

\subsection{投影值测度}
\label{b1:subsec:290401}

\begin{definition}[投影值测度]
\label{b1:def:29-pvm}
设 \(\Omega\) 是度量空间，\(H\) 是复 Hilbert 空间。
若映射 \(E:\mathcal B(\Omega)\to\mathcal B(H)\) 满足下列条件，
则称它为\term[touying-zhi-cedu]{投影值测度}[projection-valued measure]：
\begin{enumerate}
\item \(E(\varnothing)=0\)、\(E(\Omega)=I\)，且每个 \(E(B)\) 是正交投影；
\item \(E(B\cap C)=E(B)E(C)\) 对所有 Borel 集 \(B,C\) 成立；
\item 若 \((B_n)\) 两两不交，则对每个 \(x\in H\)，有
\[
 E\left(\bigcup_{n=1}^\infty B_n\right)x
 =\sum_{n=1}^\infty E(B_n)x,
\]
其中右侧在 \(H\) 的范数中收敛。
\end{enumerate}
\end{definition}

可数可加性中的收敛是逐向量的强收敛。它通常不成立于算子范数：
在 \(\ell^2\) 上，投影到前 \(n\) 个坐标的 \(P_n\) 强收敛于 \(I\)，
但 \(\|I-P_n\|=1\) 对每个 \(n\) 都成立。
从定义还可看出，不交集合的投影有正交值域。

\begin{proposition}[投影值测度的标量化]
\label{b1:prop:29-pvm-scalarization}
设 \(E\) 是度量空间 \(\Omega\) 上、作用于复 Hilbert 空间 \(H\) 的投影值测度。
对 \(x,y\in H\)，令
\[
 \mu_{x,y}(B)=\langle E(B)x,y\rangle,\qquad
 \mu_x(B)=\mu_{x,x}(B)=\|E(B)x\|^2.
\]
则 \(\mu_x\) 是总质量为 \(\|x\|^2\) 的有限正 Borel 测度，
\(\mu_{x,y}\) 是有限复 Borel 测度，且
\( |\mu_{x,y}|(\Omega)\leq\|x\|\cdot\|y\|\)。
\end{proposition}
\begin{proof}
正交投影给出 \(\langle E(B)x,x\rangle=\|E(B)x\|^2\)。
对不交集合，投影后的向量两两正交；由可数可加性和勾股等式，\(\mu_x\) 可数可加。
取整个空间得其总质量。
对复值集合函数，可数可加性由内积的连续性得到。
又因为 \(\langle E(B)x,y\rangle=\langle E(B)x,E(B)y\rangle\)，
对任意有限 Borel 分割 \((B_j)\)，有
\[
 \sum_j|\mu_{x,y}(B_j)|
 \leq\sum_j\|E(B_j)x\|\cdot\|E(B_j)y\|
 \leq\|x\|\cdot\|y\|.
\]
对分割取上确界，得到全变差估计。
\end{proof}

\begin{proposition}[从乘法表示构造投影]
\label{b1:prop:29-pvm-construction}
设 \(H\ne\{0\}\) 是复 Hilbert 空间，\(T\in\mathcal B(H)\) 且 \(T=T^*\)，
\(K=\sigma(T)\)。取\cref{b1:thm:29-bounded-spectral-representation}中的酉算子 \(U\)，定义
\[
 E(B)=U^{-1}\left(\bigoplus_\alpha M_{\mathbf1_B}\right)U
 \qquad(B\in\mathcal B(K)).
\]
则 \(E\) 是 \(K\) 上的投影值测度。
\end{proposition}
\begin{proof}
乘以示性函数是正交投影，并且
\(M_{\mathbf1_B}M_{\mathbf1_C}=M_{\mathbf1_{B\cap C}}\)，故前两个条件成立。
若 \((B_n)\) 不交，记 \(B=\bigcup_n B_n\)。对 \(u=(u_\alpha)\) 属于表示空间，有
\[
 \left\|\left(\bigoplus_\alpha M_{\mathbf1_{B\setminus\cup_{n\leq N}B_n}}\right)u\right\|^2
 =\sum_\alpha\int_{B\setminus\cup_{n\leq N}B_n}|u_\alpha|^2\,\mathrm{d}\mu_\alpha.
\]
各项随 \(N\to\infty\) 趋于零，且受 \(\|u_\alpha\|_2^2\) 控制。
先舍去平方范数和任意小的分量尾部，再对余下有限个分量用积分的控制收敛定理，
可知整个和趋于零。用 \(U^{-1}\) 搬回原空间，即得强可数可加性。
\end{proof}

\subsection{谱积分的构造}
\label{b1:subsec:290402}

对投影值测度 \(E\) 和简单 Borel 函数 \(s=\sum_{j=1}^m a_j\cdot\mathbf1_{B_j}\)，
其中 \((B_j)\) 是 \(\Omega\) 的有限 Borel 分割，定义
\[
 s(E)=\int_\Omega s\,\mathrm{d}E=\sum_{j=1}^m a_jE(B_j).
\]
共同细分两个分割，并用投影的有限可加性，可知此定义与表示方式无关。
这称为\term[pu-jifen]{谱积分}[spectral integral]；后面同一记号也将用于无界函数。
\symidx{\(\int f\,\mathrm{d}E\)：谱积分}

\ProseParagraphBreak

不交集合的正交性给出
\[
 \|s(E)x\|^2=\sum_j|a_j|^2\cdot\|E(B_j)x\|^2
 =\int_\Omega|s|^2\,\mathrm{d}\mu_x
 \leq\|s\|_\infty^2\cdot\|x\|^2.
\]
因此对任意有界 Borel 函数 \(f\)，可以选择一致逼近它的简单 Borel 函数 \(s_n\)，
并定义 \(f(E)=\lim_n s_n(E)\)，极限取算子范数。
复函数的简单逼近可分别把实部和虚部的有界值域划分成长度趋于零的区间。
估计 \(\|s_n(E)-s_m(E)\|\leq\|s_n-s_m\|_\infty\)
保证极限存在，也保证定义不依赖逼近序列。

\begin{proposition}[有界 Borel 函数的谱积分]
\label{b1:prop:29-borel-calculus}
设 \(E\) 是度量空间 \(\Omega\) 上、作用于非零复 Hilbert 空间 \(H\) 的投影值测度。
对有界复 Borel 函数 \(f,g\) 和 \(a,b\in\mathbb C\)，有
\[
 (af+bg)(E)=af(E)+bg(E),\qquad
 (fg)(E)=f(E)g(E),\qquad
 \overline f(E)=f(E)^*.
\]
并且对 \(x,y\in H\)，
\[
 \langle f(E)x,y\rangle=\int_\Omega f\,\mathrm{d}\mu_{x,y},\qquad
 \|f(E)x\|^2=\int_\Omega|f|^2\,\mathrm{d}\mu_x.
\]
算子范数的精确值为
\[
 \|f(E)\|=\inf\{c\geq0:E(\{|f|>c\})=0\}.
\]
\end{proposition}
\begin{proof}
对简单函数，在共同分割上计算，线性、乘法和伴随等式分别来自有限和、
\(E(B)E(C)=E(B\cap C)\) 和 \(E(B)^*=E(B)\)。
内积积分式也由有限和得到，范数平方公式已在上面证明。
把 \(f,g\) 分别作一致简单逼近，有限测度上的积分、算子乘法与伴随都可取极限，
因而上述等式对一般有界 Borel 函数成立。

记最后一个下确界为 \(m\)。取一列可行的 \(c_n\downarrow m\)，
则 \(\{|f|>m\}=\bigcup_n\{|f|>c_n\}\)，所以 \(E(\{|f|>m\})=0\)。
这里可数个零投影集合的并仍给零投影，因为对每个 \(x\) 相应正测度都为零。
因此范数平方公式给出 \(\|f(E)\|\leq m\)。
若 \(0\leq c<m\)，则 \(E(B)\ne0\)，其中 \(B=\{|f|>c\}\)。
在 \(\operatorname{ran}E(B)\) 中取单位向量 \(x\)，则 \(\mu_x(B)=1\)，从而
\(\|f(E)x\|^2=\int_B|f|^2\,\mathrm{d}\mu_x\geq c^2\)。
令 \(c\uparrow m\)，得到反向不等式；\(m=0\) 时反向不等式自动成立。
\end{proof}

在这个演算中，两个函数给出同一算子，当且仅当它们在某个零投影集合之外相等。
这是由对 \(f-g\) 应用范数平方公式得到的，不能把“零测集”理解为某个任意选定的参考测度的零集。

\subsection{有界 Borel 函数演算}
\label{b1:subsec:290403}

\begin{theorem}[有界自伴算子的谱测度]
\label{b1:thm:29-bounded-pvm}
设 \(H\ne\{0\}\) 是复 Hilbert 空间，\(T\in\mathcal B(H)\) 且 \(T=T^*\)，
\(K=\sigma(T)\)。存在唯一的 \(K\) 上投影值测度 \(E\)，使
\[
 T=\int_K t\,\mathrm{d}E(t).
\]
对 \(f\in C(K)\)，谱积分 \(\int_K f\,\mathrm{d}E\) 与连续函数演算的 \(f(T)\) 相同。
对每个 \(x\in H\)，测度 \(B\mapsto\langle E(B)x,x\rangle\)
就是\cref{b1:prop:29-scalar-spectral-measure}的 \(\mu_x\)。
\end{theorem}
\begin{proof}
用\cref{b1:prop:29-pvm-construction}的 \(E\)。简单函数积分在乘法表示中为 \(\bigoplus M_s\)，
对有界函数作一致简单逼近，得到 \(f(E)=U^{-1}(\bigoplus M_f)U\)。
当 \(f\) 连续时，上一节的乘法表示已说明此算子等于 \(f(T)\)；取 \(f(t)=t\) 即得存在性。

若 \(F\) 是另一个满足题设积分式的投影值测度，
其积分由乘法性质先对多项式、再对一致极限给出
\(\int_K f\,\mathrm{d}F=f(T)\) 对所有 \(f\in C(K)\) 成立。
因而对每个 \(x\)，\(E\) 与 \(F\) 的标量正测度在所有连续函数上积分相同。
在紧度量空间上，这保证两个有限 Borel 测度相同：
对开集 \(O\ne K\)，连续函数
\(h_n(t)=\min\{1,n\operatorname{dist}(t,K\setminus O)\}\)
递增趋于 \(\mathbf1_O\)，故单调收敛给出开集测度相同；\(O=K\) 时取常函数 \(1\)。
开集对有限交封闭且生成 Borel 代数，有限测度的唯一性定理随后给出全部 Borel 集上相同。
对每个 \(B\)，两个自伴算子 \(E(B),F(B)\) 的二次型于是相同，极化恒等式给出 \(E(B)=F(B)\)。
相同的标量测度唯一性论证，把 \(E\) 的标量正测度与上一节的 \(\mu_x\) 识别起来。
\end{proof}

这个唯一的 \(E\) 称为 \(T\) 的\term[pu-cedu]{谱测度}[spectral measure]。
对有界 Borel 函数 \(f\)，以后统一记
\[
 f(T)=\int_K f\,\mathrm{d}E.
\]
这就是 \(T\) 的\term[borel-hanshu-yansuan]{Borel 函数演算}[Borel functional calculus]。
它的范数由 \(E\) 的零集决定，未必等于普通的逐点上确界。

\begin{proposition}[谱投影与特征空间]
\label{b1:prop:29-eigenprojection}
设 \(H\ne\{0\}\) 是复 Hilbert 空间，\(T\in\mathcal B(H)\) 且 \(T=T^*\)，
\(E\) 是它在 \(K=\sigma(T)\) 上的谱测度。
则 \(E(\{\lambda\}\cap K)\) 的值域等于 \(\ker(T-\lambda I)\)，其中 \(\lambda\in\mathbb R\)。
此外，\(K\) 的每个非空相对开集 \(O\) 都满足 \(E(O)\ne0\)。
\end{proposition}
\begin{proof}
范数平方公式给出
\[
 \|(T-\lambda I)x\|^2=\int_K|t-\lambda|^2\,\mathrm{d}\mu_x(t).
\]
它等于零，当且仅当 \(\mu_x\) 集中于 \(\{\lambda\}\cap K\)，
也就是 \(E(\{\lambda\}\cap K)x=x\)。
若 \(O\) 非空，取非零连续函数 \(f\) 使 \(\operatorname{supp}f\subset O\)；
这样的函数由紧空间上的连续截断构造得到。
若 \(E(O)=0\)，则 \(f(T)=f(T)E(O)=0\)，与连续演算的
\(\|f(T)\|=\|f\|_{C(K)}>0\) 矛盾。
\end{proof}

例如，趋零对角算子的 \(E(B)\) 保留所有满足 \(1/n\in B\) 的坐标，
但 \(E(\{0\})=0\)。点 \(0\) 仍属于谱，因为每个包含它的相对开邻域都有非零谱投影。
对 \(L^2(0,1)\) 上的坐标乘法算子，则有 \(E(B)f=\mathbf1_B\cdot f\)；
所有单点投影都为零。

\subsection{谱积分的极限与运算}
\label{b1:subsec:290404}

\begin{proposition}[有界逐点收敛与谱演算]
\label{b1:prop:29-borel-dominated}
设 \(E\) 是度量空间 \(\Omega\) 上、作用于复 Hilbert 空间 \(H\) 的投影值测度。
设 \(f_n,f\) 是有界复 Borel 函数，\(\sup_n\|f_n\|_\infty<\infty\)，
且在某个 \(E(N)=0\) 的 Borel 集之外有 \(f_n\to f\)。
则 \(f_n(E)x\to f(E)x\) 对所有 \(x\in H\) 成立。
\end{proposition}
\begin{proof}
对每个固定 \(x\)，零投影集合也是 \(\mu_x\) 的零测集。
于是标量控制收敛定理给出
\[
 \|(f_n(E)-f(E))x\|^2
 =\int_\Omega|f_n-f|^2\,\mathrm{d}\mu_x\longrightarrow0.
\]
\end{proof}

\begin{proposition}[与谱投影交换的判别]
\label{b1:prop:29-spectral-commutation}
设 \(H\ne\{0\}\) 是复 Hilbert 空间，\(T\in\mathcal B(H)\) 且 \(T=T^*\)，
\(E\) 是它的谱测度，\(S\in\mathcal B(H)\)。
则 \(ST=TS\) 当且仅当 \(S\) 与每个 \(E(B)\) 交换；
这也等价于 \(S\) 与所有有界 Borel 演算 \(f(T)\) 交换。
\end{proposition}
\begin{proof}
若 \(S\) 与所有谱投影交换，先对简单函数求和，再取一致极限，
得到它与所有 \(f(T)\) 交换，特别与 \(T\) 交换。
反过来，若 \(ST=TS\)，上一节已证明 \(S\) 与每个连续演算交换。
固定 Borel 集 \(B\subset K=\sigma(T)\) 和 \(x\in H\)。
有限正 Radon 测度 \(\mu_x+\mu_{Sx}\) 上的连续函数在 \(L^2\) 中稠密，
所以可取 \(h_n\in C(K)\) 使
\(h_n\to\mathbf1_B\) 于 \(L^2(\mu_x+\mu_{Sx})\)。
范数平方公式分别给出
\[
 h_n(T)x\to E(B)x,\qquad h_n(T)Sx\to E(B)Sx.
\]
在 \(S h_n(T)x=h_n(T)Sx\) 中取极限，得到 \(SE(B)x=E(B)Sx\)。
这对每个 \(x\) 成立，故算子交换。
\end{proof}

为下一节的 Cayley 变换作准备，还需要酉算子的谱表示。
这里可以把它拆成两个相互交换的自伴算子处理，而无需先建立一般正规算子的理论。

\begin{theorem}[酉算子的谱测度]
\label{b1:thm:29-unitary-spectral}
设 \(H\ne\{0\}\) 是复 Hilbert 空间，\(U\in\mathcal B(H)\) 是酉算子，
\(\mathbb T=\{z\in\mathbb C:|z|=1\}\)。
存在唯一的 \(\mathbb T\) 上投影值测度 \(F\)，使
\[
 U=\int_{\mathbb T}z\,\mathrm{d}F(z).
\]
记 \(f(U)=\int_{\mathbb T}f\,\mathrm{d}F\)。有界 Borel 函数的代数运算、伴随、
范数平方公式及强控制收敛均按上面的谱积分结论成立。
此外，\(\operatorname{ran}F(\{z\})=\ker(U-zI)\) 对 \(z\in\mathbb T\) 成立。
\end{theorem}
\begin{proof}
令
\[
 C=\frac{U+U^*}{2},\qquad D=\frac{U-U^*}{2i}.
\]
由 \(U^*U=UU^*=I\)，可知 \(C,D\) 自伴、相互交换，且
\(C^2+D^2=I\)、\(\|C\|\leq1\)、\(\|D\|\leq1\)。
取 \(D\) 的谱测度 \(E_D\)，令
\[
 Q_+=E_D([0,\infty)\cap\sigma(D)),\qquad
 Q_-=E_D((-\infty,0)\cap\sigma(D)).
\]
则 \(Q_++Q_-=I\)，其值域 \(H_+,H_-\) 正交。
由\cref{b1:prop:29-spectral-commutation}，它们同时约化 \(C,D\)。
在 \(H_+\) 上，\(D_+=D|_{H_+}\geq0\)；在 \(H_-\) 上，\(-D_-\geq0\)，
因为相应二次型是正半轴或负半轴上的坐标积分。
空的子空间直接略去。

令 \(C_\pm=C|_{H_\pm}\)。由于 \(D_\pm^2=I-C_\pm^2\)，
正平方根的唯一性给出
\[
 D_+=(I-C_+^2)^{1/2},\qquad
 D_-=-(I-C_-^2)^{1/2}.
\]
令 \(E_\pm\) 为 \(C_\pm\) 的谱测度，
\(\phi_\pm(t)=t\pm i\sqrt{1-t^2}\)（\(-1\leq t\leq1\)）。
连续演算的复合性质说明 \(U|_{H_\pm}=\phi_\pm(C_\pm)\)。
因此定义
\[
 F(B)=E_+(\phi_+^{-1}(B))\oplus E_-(\phi_-^{-1}(B))
 \qquad(B\in\mathcal B(\mathbb T)),
\]
其中逆像分别限制在 \(\sigma(C_\pm)\) 上。
逆像保持交、并和补集，故 \(F\) 是投影值测度。
简单函数积分的换元，再经一致逼近，给出
\(\int z\,\mathrm{d}F=\phi_+(C_+)\oplus\phi_-(C_-)=U\)。

若 \(F_1,F_2\) 都给出这个积分式，则它们对函数 \(z\) 的积分相同，
由伴随性质，对 \(\overline z\) 的积分也相同。
所以它们对 \(z,\overline z\) 的所有多项式积分相同。
这些函数组成含常数、分离点且对共轭封闭的代数，
由 Stone--Weierstrass 定理在 \(C(\mathbb T)\) 中稠密。
于是两者对全部连续函数的积分相同，
再按\cref{b1:thm:29-bounded-pvm}证明中的标量测度唯一性与极化论证，得到 \(F_1=F_2\)。
最后，范数平方公式给出
\[
 \|(U-zI)x\|^2=\int_{\mathbb T}|\zeta-z|^2\,\mathrm{d}\mu_x(\zeta),
\]
它为零当且仅当 \(F(\{z\})x=x\)，从而得到特征空间公式。
\end{proof}

有界自伴谱测度与酉谱测度都由算子唯一确定，但构造所选的循环子空间或正负子空间只是一种证明方法。
在下一节以后，谱积分还会作用于无界函数；这时函数相乘的代数等式必须连同定义域一起解释。

% !TeX root = ../../Book1.tex
\section{稠密定义算子与伴随}
\label{b1:sec:2905}

% 实读：Rudin1991Functional，第2版，13.1（pp.347--348）、13.3（p.349）、
% 13.7--13.12（pp.351--354）与13.16（p.356）。图像正交补和闭算子双伴随
% 与原书核对；一般可闭性、闭包伴随不变及非实点满射判据在本节直接证明，
% 不把原书13.12只陈述的闭算子版本冒称为已读到全部可闭版本。

有界算子的连续性使定义域上的极限自动与算子作用相容。求导和无界乘法不再具有这种性质：
只知道 \(x_n\to x\)，并不能决定 \(Ax_n\) 是否收敛，甚至不能保证 \(x\) 仍在定义域中。
因此，算子不只是一条作用公式，还包括允许输入的线性子空间。
本节把\secref{b1:sec:0803}的闭图像与图像范数推广到稠密定义的 Hilbert 空间算子，
再用伴随把定义域的极限问题化成正交关系。所有内积仍对第一变量线性。
稠密定义、伴随和图像方法的系统论述可对照 Rudin\cite[13.1, 13.3, 13.7--13.12]{Rudin1991Functional}。

\subsection{定义域、图像与图像范数}
\label{b1:subsec:290501}

\begin{definition}[稠密定义与算子扩张]
\label{b1:def:29-densely-defined}
设 \(H\) 为复 Hilbert 空间，\(D(A)\subseteq H\) 为线性子空间，
\(A:D(A)\to H\) 为线性映射。若 \(\overline{D(A)}=H\)，则称 \(A\) 为
\term[choumi-dingyi-suanzi]{稠密定义算子}[densely defined operator]。
若另一个线性算子 \(B:D(B)\subseteq H\to H\) 满足
\[
 D(A)\subseteq D(B),\qquad Bx=Ax\quad(x\in D(A)),
\]
则称 \(B\) 为 \(A\) 的扩张，记作 \(A\subseteq B\)。等号 \(A=B\) 同时要求定义域相等、
作用相同。算子的图像记为
\[
 G(A)=\{(x,Ax):x\in D(A)\}\subseteq H\oplus H.
\]
\symidx{\(A\subseteq B\)：线性算子的扩张关系}
\end{definition}

这里的稠密性使用原空间 \(H\) 的范数。它保证连续线性泛函在 \(D(A)\) 上的取值可以唯一决定
其在 \(H\) 上的延拓，却不保证 \(A\) 自身连续。若 \(D(A)=H\) 且 \(A\) 闭，
\cref{b1:thm:closed-graph}才给出 \(A\) 有界；无界闭算子的定义域因而必须是真子空间。

\ProseParagraphBreak

沿用\cref{b1:def:closed-operator,b1:def:graph-norm}，图像 \(G(A)\) 闭时称 \(A\) 闭，
并在 \(D(A)\) 上取图像范数
\[
 \|x\|_A=\|x\|+\|Ax\|.
\]
有时计算更适合使用平方和范数
\[
 \|x\|_{A,2}=\bigl(\|x\|^2+\|Ax\|^2\bigr)^{1/2},\qquad
 \langle x,y\rangle_A=\langle x,y\rangle+\langle Ax,Ay\rangle.
\]
二者满足 \(\|x\|_{A,2}\leq\|x\|_A\leq\sqrt2\cdot\|x\|_{A,2}\)，
所以给出相同的收敛、闭包和完备性。后面谈图像范数时仍沿用前一种记号。

\begin{proposition}[图像范数与闭性]
\label{b1:prop:29-graph-complete}
设 \(H\) 为复 Hilbert 空间，\(A:D(A)\subseteq H\to H\) 为线性算子。
则 \(A\) 闭，当且仅当 \(D(A)\) 关于 \(\|x\|_A=\|x\|+\|Ax\|\) 完备。
在这一情形，\(D(A)\) 关于内积
\(\langle x,y\rangle_A=\langle x,y\rangle+\langle Ax,Ay\rangle\) 是 Hilbert 空间。
\end{proposition}
\begin{proof}
若 \(A\) 闭，图像范数的完备性已由\cref{b1:prop:closed-operator-graph-norm}给出：
图像范数 Cauchy 列满足 \(x_n\to x\)、\(Ax_n\to y\)，闭性给出 \(x\in D(A)\) 和 \(Ax=y\)。
平方和范数与它等价，且由所列内积产生，故也完备。

反之，若 \(D(A)\) 关于图像范数完备，取 \(x_n\in D(A)\)，使
\((x_n,Ax_n)\to(x,y)\) 于 \(H\oplus H\)。两个分量都为 Cauchy 列，故 \((x_n)\)
在图像范数中为 Cauchy 列，存在 \(u\in D(A)\) 使 \(x_n\to u\)、\(Ax_n\to Au\)。
由极限唯一性，\(u=x\)、\(Au=y\)，从而 \((x,y)\in G(A)\)。
\end{proof}

例如，把恒等算子限制到 \(\ell^2\) 中的有限支撑序列空间 \(c_{00}\)，所得算子稠密定义，
却不是闭算子。有限截断 \(x^{(n)}\) 可以趋于一个不在 \(c_{00}\) 中的向量 \(x\)，
而 \((x^{(n)},x^{(n)})\to(x,x)\)。稠密定义与图像闭是两项不同的要求。

\subsection{稠密定义算子的伴随}
\label{b1:subsec:290502}

对有界算子，\cref{b1:prop:hilbert-adjoint}在每个 \(y\in H\) 上构造了伴随值。
当 \(A\) 无界时，固定 \(y\) 后得到的泛函 \(x\mapsto\langle Ax,y\rangle\)
未必关于 \(\|x\|\) 连续。伴随的定义域恰好由仍具有这种连续性的测试向量组成。

\begin{definition}[稠密定义算子的伴随]
\label{b1:def:29-unbounded-adjoint}
设 \(H\) 为复 Hilbert 空间，\(A:D(A)\subseteq H\to H\) 为稠密定义线性算子。
定义
\[
 D(A^*)=\{y\in H:\text{存在 }C_y<\infty\text{，使 }
 |\langle Ax,y\rangle|\leq C_y\cdot\|x\|\quad(x\in D(A))\}.
\]
对每个 \(y\in D(A^*)\)，以 \(A^*y\) 表示满足
\begin{equation}\label{b1:eq:29-adjoint-pairing}
 \langle Ax,y\rangle=\langle x,A^*y\rangle\qquad(x\in D(A))
\end{equation}
的唯一向量。所得线性算子 \(A^*:D(A^*)\subseteq H\to H\) 称为 \(A\) 的伴随。
\end{definition}

定义中的表示向量确实存在。给定 \(y\in D(A^*)\)，相应泛函在稠密子空间 \(D(A)\) 上有界，
所以可按范数极限唯一延拓到 \(H\)；\cref{b1:thm:hilbert-riesz}再给出唯一表示向量。
反过来，只要某个向量满足\eqref{b1:eq:29-adjoint-pairing}，Cauchy--Schwarz 不等式便给出定义中的界。
第二变量的共轭线性说明，对 \(y_1,y_2\in D(A^*)\) 和 \(\alpha,\beta\in\mathbb C\)，
表示向量为 \(\alpha A^*y_1+\beta A^*y_2\)，故 \(D(A^*)\) 为线性子空间，\(A^*\) 线性。
此时还不能断言 \(D(A^*)\) 稠密。

\begin{proposition}[伴随的闭性与扩张反向性]
\label{b1:prop:29-adjoint-closed}
设 \(H\) 为复 Hilbert 空间，\(A:D(A)\subseteq H\to H\) 为稠密定义线性算子。
则 \(A^*\) 是闭算子。若 \(B:D(B)\subseteq H\to H\) 也是稠密定义线性算子且
\(A\subseteq B\)，则 \(B^*\subseteq A^*\)。此外，对任意 \(z\in\mathbb C\)，
\[
 D((A-zI)^*)=D(A^*),\qquad (A-zI)^*=A^*-\overline zI.
\]
\end{proposition}
\begin{proof}
设 \(y_n\in D(A^*)\)、\(y_n\to y\)、\(A^*y_n\to v\)。对每个固定 \(x\in D(A)\)，
有
\[
 \langle Ax,y\rangle=\lim_n\langle Ax,y_n\rangle
 =\lim_n\langle x,A^*y_n\rangle=\langle x,v\rangle.
\]
故 \(y\in D(A^*)\)、\(A^*y=v\)，证明了闭性。
若 \(y\in D(B^*)\)，把 \(\langle Bx,y\rangle=\langle x,B^*y\rangle\)
限制到 \(D(A)\)，即得 \(y\in D(A^*)\)、\(A^*y=B^*y\)。

最后，泛函 \(x\mapsto z\langle x,y\rangle\) 对 \(\|x\|\) 有界，
所以 \(\langle(A-zI)x,y\rangle\) 有界当且仅当 \(\langle Ax,y\rangle\) 有界。
在共同定义域上，
\[
 \langle(A-zI)x,y\rangle
 =\langle x,A^*y\rangle-z\langle x,y\rangle
 =\langle x,(A^*-\overline zI)y\rangle,
\]
得到所需公式。
\end{proof}

\begin{proposition}[核、值域与图像的正交关系]
\label{b1:prop:29-adjoint-range}
设 \(H\) 为复 Hilbert 空间，\(A:D(A)\subseteq H\to H\) 为稠密定义线性算子。
对每个 \(z\in\mathbb C\)，有
\begin{equation}\label{b1:eq:29-range-kernel}
 \operatorname{ran}(A-zI)^\perp=\ker(A^*-\overline zI),\qquad
 \overline{\operatorname{ran}(A-zI)}
 =\ker(A^*-\overline zI)^\perp.
\end{equation}
在 \(H\oplus H\) 中取分量内积之和，则
\begin{equation}\label{b1:eq:29-graph-orthogonal}
 G(A)^\perp=\{(-A^*y,y):y\in D(A^*)\}.
\end{equation}
\end{proposition}
\begin{proof}
向量 \(y\) 属于 \(\operatorname{ran}(A-zI)^\perp\)，当且仅当
\(\langle(A-zI)x,y\rangle=0\) 对每个 \(x\in D(A)\) 成立。
由伴随定义，这等价于 \(y\in D((A-zI)^*)\)、\((A-zI)^*y=0\)。
使用前一命题的移位公式即得第一个等式，再由\cref{b1:thm:hilbert-projection}
中的双重正交补公式得到第二个等式。

同样，\((u,v)\in G(A)^\perp\) 当且仅当
\[
 \langle x,u\rangle+\langle Ax,v\rangle=0\qquad(x\in D(A)).
\]
这恰好要求 \(v\in D(A^*)\)、\(A^*v=-u\)，从而得到图像公式。
\end{proof}

这些关系没有把值域的闭性包含在内。核的正交补给出的通常是值域的闭包；
只有另有闭性依据，才可以去掉上面的闭包符号。

\subsection{对称算子与自伴算子}
\label{b1:subsec:290503}

\begin{definition}[对称与自伴]
\label{b1:def:29-symmetric-self-adjoint}
设 \(H\) 为复 Hilbert 空间，\(A:D(A)\subseteq H\to H\) 为稠密定义线性算子。
若
\[
 \langle Ax,y\rangle=\langle x,Ay\rangle\qquad(x,y\in D(A)),
\]
则称 \(A\) 为\term[duichen-suanzi]{对称算子}[symmetric operator]。
若 \(A=A^*\)，即 \(D(A)=D(A^*)\) 且二者在共同定义域上相等，
则称 \(A\) 为自伴算子。
\end{definition}

对称性等价于 \(A\subseteq A^*\)，因为对每个 \(y\in D(A)\)，
右侧 \(\langle x,Ay\rangle\) 关于 \(\|x\|\) 有界。
自伴性还要求伴随不带来任何新的定义域向量。
当 \(D(A)=H\) 且 \(A\) 有界时，伴随也定义在整个 \(H\) 上，
这就退化为\cref{b1:def:self-adjoint-operator}的定义；对于无界算子，两个要求必须分开检查。

\begin{lemma}[对称算子的非实参数估计]
\label{b1:lem:29-symmetric-nonreal-bound}
设 \(H\) 为复 Hilbert 空间，\(A:D(A)\subseteq H\to H\) 为稠密定义对称算子。
对 \(z=a+\mathrm ib\in\mathbb C\) 和 \(x\in D(A)\)，其中 \(a,b\in\mathbb R\)，有
\begin{equation}\label{b1:eq:29-symmetric-nonreal-bound}
 \|(A-zI)x\|^2=\|(A-aI)x\|^2+b^2\cdot\|x\|^2.
\end{equation}
因此当 \(b\ne0\) 时，\(A-zI\) 单射，且
\(\|(A-zI)x\|\geq|b|\cdot\|x\|\)。若 \(b\ne0\) 且 \(A\) 还闭，则 \(\operatorname{ran}(A-zI)\) 闭。
\end{lemma}
\begin{proof}
由对称性，\(\langle Ax,x\rangle=\overline{\langle Ax,x\rangle}\)，故它为实数。
展开 \(\|(A-aI)x-\mathrm ibx\|^2\) 时，两项交叉项的和为零，得到所列恒等式。

再设 \(A\) 闭、\(b\ne0\)，且 \(y_n=(A-zI)x_n\to y\)。下界给出
\[
 \|x_n-x_m\|\leq |b|^{-1}\cdot\|y_n-y_m\|,
\]
所以 \(x_n\to x\in H\)。又有 \(Ax_n=y_n+zx_n\to y+zx\)。
闭性给出 \(x\in D(A)\)、\(Ax=y+zx\)，故 \(y=(A-zI)x\) 属于值域。
\end{proof}

下界保证解的唯一性和对数据的稳定性。要证明自伴，仍须知道两个共轭非实参数下都能解出每个数据。

\begin{theorem}[自伴性的非实点满射判据]
\label{b1:thm:29-self-adjoint-criterion}
设 \(H\) 为复 Hilbert 空间，\(A:D(A)\subseteq H\to H\) 为稠密定义对称算子。
则 \(A\) 自伴，当且仅当
\[
 \operatorname{ran}(A-\mathrm iI)=H,
 \qquad \operatorname{ran}(A+\mathrm iI)=H.
\]
等价地，只需存在某个 \(z\in\mathbb C\setminus\mathbb R\)，使
\(\operatorname{ran}(A-zI)=\operatorname{ran}(A-\overline zI)=H\)。
若这些条件成立，则 \(A\) 闭，并且对每个 \(z\in\mathbb C\setminus\mathbb R\)，
\(A-zI:D(A)\to H\) 是双射，其逆作为 \(H\to H\) 的算子满足
\begin{equation}\label{b1:eq:29-unbounded-resolvent-bound}
 \|(A-zI)^{-1}\|\leq |\operatorname{Im}z|^{-1}.
\end{equation}
\end{theorem}
\begin{proof}
先设 \(A=A^*\)。伴随的闭性保证 \(A\) 闭。
对非实 \(z\)，前一引理说明 \(A-zI\) 单射且值域闭；
\cref{b1:prop:29-adjoint-range}又给出
\[
 \operatorname{ran}(A-zI)^\perp
 =\ker(A-\overline zI)=\{0\}.
\]
因此值域同时稠密且闭，必为 \(H\)。
下界 \(\|(A-zI)x\|\geq|\operatorname{Im}z|\cdot\|x\|\)
随后给出逆算子的范数估计，特别得到 \(z=\mathrm i,-\mathrm i\) 的满射性。

反过来，假设对于某个非实 \(z\)，两个共轭参数下的值域都等于 \(H\)。
对称性已给出 \(A\subseteq A^*\)。任取 \(y\in D(A^*)\)，
由 \(A-zI\) 满射，可以取 \(x\in D(A)\)，使
\[
 (A-zI)x=(A^*-zI)y.
\]
因为 \(A^*x=Ax\)，向量 \(w=y-x\) 满足 \((A^*-zI)w=0\)。
另一方面，
\[
 \ker(A^*-zI)=\operatorname{ran}(A-\overline zI)^\perp=\{0\}.
\]
故 \(y=x\in D(A)\)。这证明了 \(D(A^*)\subseteq D(A)\)，
于是 \(A=A^*\)。闭性和其余非实参数的结论由已经证明的正向论证得到。
取 \(z=\mathrm i\) 即得所述特例。
\end{proof}

证明的反向没有预设 \(A\) 闭。两侧满射先排除了伴随的额外定义域，再由伴随的闭性得到闭性。
而仅有 \(A\subseteq A^*\) 时，这两步都不能省略。
对于自伴 \(A\) 和实数 \(t\)，\eqref{b1:eq:29-range-kernel}同时给出
\[
 \overline{\operatorname{ran}(A-tI)}=\ker(A-tI)^\perp.
\]
它和有界自伴算子的核、值域关系具有相同形式，但不能据此断言值域总是闭的。

\subsection{闭性、可闭性与本质自伴}
\label{b1:subsec:290504}

闭算子要求极限已经留在原图像内。另一个问题是：能否通过补上缺失的图像极限得到一个闭算子，
而不使同一个输入出现两个输出？

\begin{definition}[可闭算子与闭包]
\label{b1:def:29-closable}
设 \(H\) 为复 Hilbert 空间，\(A:D(A)\subseteq H\to H\) 为线性算子。
若 \(\overline{G(A)}\subseteq H\oplus H\) 仍是某个线性算子的图像，
则称 \(A\) \term[kebi-suanzi]{可闭}[closable]，并把该算子记为 \(\overline A\)，
称为 \(A\) 的闭包。它满足 \(G(\overline A)=\overline{G(A)}\)，
是 \(A\) 的最小闭扩张。
\symidx{\(\overline A\)：可闭算子的闭包}
\end{definition}

\begin{lemma}[可闭性的序列判别]
\label{b1:lem:29-closable-sequence}
设 \(H\) 为复 Hilbert 空间，\(A:D(A)\subseteq H\to H\) 为线性算子。
则 \(A\) 可闭，当且仅当对任意 \(x_n\in D(A)\)，
\[
 x_n\to0,\quad Ax_n\to y\quad\Longrightarrow\quad y=0.
\]
\end{lemma}
\begin{proof}
图像的闭包是闭线性子空间。它成为一个单值线性映射的图像，当且仅当其中形如 \((0,y)\)
的向量只有 \((0,0)\)：必要性来自单值性；充分性来自两个具有相同第一分量的向量之差。
由于 \(H\oplus H\) 是度量空间，\((0,y)\in\overline{G(A)}\)
当且仅当存在 \((x_n,Ax_n)\to(0,y)\)，故得到所列判据。
任一闭扩张的图像包含 \(G(A)\)，从而包含其闭包，这也证明定义中的最小性。
\end{proof}

\begin{theorem}[可闭性与双伴随]
\label{b1:thm:29-double-adjoint}
设 \(H\) 为复 Hilbert 空间，\(A:D(A)\subseteq H\to H\) 为稠密定义线性算子。
则 \(A\) 可闭，当且仅当 \(D(A^*)\) 在 \(H\) 中稠密。
在这些等价条件成立时，\(\overline A\) 稠密定义，双伴随 \(A^{**}=(A^*)^*\) 有定义，并满足
\[
 A^{**}=\overline A,
 \qquad (\overline A)^*=A^*.
\]
特别地，稠密定义的闭算子满足 \(A^{**}=A\)。
\end{theorem}
\begin{proof}
记 \(E=\overline{G(A)}\)。由\eqref{b1:eq:29-graph-orthogonal}及双重正交补公式，
\((0,h)\in E\) 当且仅当
\[
 \langle(0,h),(-A^*v,v)\rangle=\langle h,v\rangle=0
 \qquad(v\in D(A^*)).
\]
所以 \(E\) 中的竖直向量恰为 \(\{0\}\times D(A^*)^\perp\)。
前一引理说明，\(E\) 是某个算子的图像当且仅当 \(D(A^*)^\perp=\{0\}\)，
也就是 \(D(A^*)\) 稠密。

在此条件下，伴随 \(A^{**}\) 可以按稠密定义算子的规则构造。
对 \(x,y\in H\)，同一正交补公式给出
\[
 \begin{split}
 (x,y)\in E
 &\iff -\langle x,A^*v\rangle+\langle y,v\rangle=0
       \quad(v\in D(A^*))\\
 &\iff \langle A^*v,x\rangle=\langle v,y\rangle
       \quad(v\in D(A^*))\\
 &\iff x\in D(A^{**}),\quad A^{**}x=y.
 \end{split}
\]
第二个等价关系取了复共轭。由此 \(G(A^{**})=E\)，故 \(A^{**}=\overline A\)。
其定义域包含原来的稠密子空间 \(D(A)\)，所以也稠密。

由 \(A\subseteq\overline A\) 和扩张反向性，先有 \((\overline A)^*\subseteq A^*\)。
任取 \(y\in D(A^*)\)、\(x\in D(\overline A)\)，选取 \(x_n\in D(A)\)，使
\(x_n\to x\)、\(Ax_n\to\overline Ax\)。在伴随配对中取极限，得
\[
 \langle\overline Ax,y\rangle=\langle x,A^*y\rangle.
\]
因此 \(y\in D((\overline A)^*)\)，且两伴随值相同，证明了另一个包含关系。
\end{proof}

稠密定义本身并不足以保证可闭。在复空间 \(\ell^2\) 上，取 \(D(A)=c_{00}\)，令
\[
 Ax=\left(\sum_{n\geq1}n x_n\right)e_1.
\]
和式对每个输入都是有限和。向量 \(x^{(n)}=e_n/n\) 满足 \(x^{(n)}\to0\)，
而 \(Ax^{(n)}=e_1\)，故 \(A\) 不可闭。
这里还可以直接算出
\[
 D(A^*)=\{y\in\ell^2:y_1=0\},\qquad A^*y=0.
\]
若 \(y_1\ne0\)，把单位向量 \(e_n\) 代入定义中的估计会得到无界的 \(n\cdot|y_1|\)；
若 \(y_1=0\)，配对恒为零。这一伴随的定义域确实不稠密。

\begin{definition}[本质自伴]
\label{b1:def:29-essentially-self-adjoint}
设 \(H\) 为复 Hilbert 空间，\(A:D(A)\subseteq H\to H\) 为稠密定义对称算子。
若其闭包 \(\overline A\) 自伴，则称 \(A\) \term[benzhi-ziban]{本质自伴}[essentially self-adjoint]。
\end{definition}

这个定义中的闭包总能构造。事实上，对称性给出 \(D(A)\subseteq D(A^*)\)，
故 \(D(A^*)\) 稠密；\cref{b1:thm:29-double-adjoint}说明 \(A\) 可闭。
又因 \(A\subseteq A^*\) 且 \(A^*\) 闭，有
\(\overline A\subseteq A^*=(\overline A)^*\)，所以闭包仍对称。
于是本质自伴等价于 \(\overline A=A^*\)，但不要求原来的 \(A\) 已经闭。

\begin{corollary}[本质自伴性的稠密值域判据]
\label{b1:cor:29-essential-self-adjoint}
设 \(H\) 为复 Hilbert 空间，\(A:D(A)\subseteq H\to H\) 为稠密定义对称算子。
则 \(A\) 本质自伴，当且仅当
\[
 \overline{\operatorname{ran}(A-\mathrm iI)}=H,
 \qquad \overline{\operatorname{ran}(A+\mathrm iI)}=H.
\]
\end{corollary}
\begin{proof}
闭包 \(\overline A\) 闭且对称，故\cref{b1:lem:29-symmetric-nonreal-bound}保证
\(\operatorname{ran}(\overline A\pm\mathrm iI)\) 闭。
图像闭包的定义又给出
\[
 \operatorname{ran}(\overline A\pm\mathrm iI)
 =\overline{\operatorname{ran}(A\pm\mathrm iI)}.
\]
这里一个包含关系来自闭值域包含原值域；另一个来自对每个
\(x\in D(\overline A)\) 选取 \(x_n\in D(A)\)，使
\(x_n\to x\)、\(Ax_n\to\overline Ax\)。
因此所列稠密性恰好等价于 \(\overline A\pm\mathrm iI\) 都满射。
对闭包使用\cref{b1:thm:29-self-adjoint-criterion}即得结论。
\end{proof}

% !TeX root = ../../Book1.tex
\section{无界自伴算子的谱表示}
\label{b1:sec:2906}

有界自伴算子的谱积分在整个 Hilbert 空间上作用；无界情形首先要回答的是哪些向量允许代入。
本节从自伴算子的 Cayley 变换出发，把单位圆上的酉谱测度移到实轴，
并由谱测度同时恢复算子作用和定义域。随后把坐标函数换成一般 Borel 函数，
逐项核对闭性、伴随、乘积以及极限。特别地，两个算子的作用公式相同，
还不足以断言它们相等，必须同时核对定义域。

以下 \(H\) 为复 Hilbert 空间，内积对第一变量线性。
对投影值测度 \(E\) 及有界 Borel 函数 \(h\)，暂记
\[
 h(E)=\int h\,\mathrm dE.
\]
这是\secref{b1:sec:2904}已经构造的有界谱积分；这个记号不预先要求有一个无界算子。
标量正测度记为
\[
 \mu_x(B)=\langle E(B)x,x\rangle=\|E(B)x\|^2.
\]
本节的 Cayley 证明路线参见\cite[Theorem 13.30, pp.~368--369]{Rudin1991Functional}；
一般谱积分的定义域、乘积与伴随参见\cite[Theorem 13.24, pp.~362--364]{Rudin1991Functional}。

\subsection{Cayley 变换}
\label{b1:subsec:290601}

\cref{b1:thm:29-self-adjoint-criterion}表明，对稠密定义的自伴算子 \(A\)，
\(A+iI\) 和 \(A-iI\) 都把 \(D(A)\) 双射到 \(H\)。
因此可以先用 \((A+iI)^{-1}\) 把任意向量送入定义域，再作用 \(A-iI\)。

\begin{proposition}[自伴算子的 Cayley 变换]
\label{b1:prop:29-cayley}
设 \(H\) 为复 Hilbert 空间，\(A:D(A)\subseteq H\to H\) 为稠密定义的自伴算子。则
\[
 U=(A-iI)(A+iI)^{-1}
   =I-2i(A+iI)^{-1}
\]
是 \(H\) 上的酉算子，称为 \(A\) 的\term[cayley-bianhuan]{Cayley 变换}[Cayley transform]。
它满足
\begin{equation}\label{b1:eq:29-cayley-domain}
 I-U=2i(A+iI)^{-1},\qquad
 \ker(I-U)=\{0\},\qquad
 \operatorname{ran}(I-U)=D(A),
\end{equation}
并且
\begin{equation}\label{b1:eq:29-cayley-inverse}
 A(I-U)y=i(I+U)y\qquad(y\in H).
\end{equation}
\end{proposition}
\begin{proof}
由对称性，\(\langle Ax,x\rangle\) 为实数，故对 \(x\in D(A)\) 有
\[
 \|(A\pm iI)x\|^2=\|Ax\|^2+\|x\|^2.
\]
两算子满射，再由这个等式可知它们的逆作为 \(H\to H\) 的映射有界，范数至多为 \(1\)。
它们的值域都恰为 \(D(A)\)。若 \(y=(A+iI)x\)，则
\[
 Uy=(A-iI)x=y-2ix,
\]
所以 \(U=I-2i(A+iI)^{-1}\) 在整个 \(H\) 上有界，且
\(\|Uy\|=\|(A-iI)x\|=\|(A+iI)x\|=\|y\|\)。
因为 \(A-iI\) 满射，\(U\) 也满射，故为酉算子。

由 \(I-U=2i(A+iI)^{-1}\)，立即得到核与值域的两个结论。
仍把 \(y\) 写成 \((A+iI)x\)，则
\[
 (I-U)y=2ix,\qquad (I+U)y=2Ax.
\]
前式保证 \((I-U)y\in D(A)\)，并给出
\(A(I-U)y=2iAx=i(I+U)y\)。这里每一次作用 \(A\) 的向量都已确定属于定义域。
\end{proof}

记单位圆为 \(\mathbb T=\{z\in\mathbb C:|z|=1\}\)。两个互逆的连续映射
\begin{equation}\label{b1:eq:29-cayley-maps}
 \phi(t)=\frac{t-i}{t+i}\quad(t\in\mathbb R),\qquad
 \theta(z)=i\frac{1+z}{1-z}\quad(z\in\mathbb T\setminus\{1\})
\end{equation}
把实轴与去掉 \(1\) 的单位圆对应起来。
对 \(|z|=1\)、\(z\ne1\)，利用 \(\overline z=z^{-1}\) 可核对
\(\overline{\theta(z)}=\theta(z)\)，故 \(\theta\) 确实取实值。
还将用到
\begin{equation}\label{b1:eq:29-cayley-weight}
 \bigl(1+|\theta(z)|^2\bigr)\cdot|1-z|^2
 =|1-z|^2+|1+z|^2=4.
\end{equation}

由酉谱定理\cref{b1:thm:29-unitary-spectral}，存在唯一的投影值测度 \(F\)，使
\[
 U=\int_{\mathbb T}z\,\mathrm dF(z).
\]
有界谱积分的乘法公式给出
\((I-U)F(\{1\})=0\)。结合\eqref{b1:eq:29-cayley-domain}的单射性，得到
\begin{equation}\label{b1:eq:29-cayley-no-atom}
 F(\{1\})=0.
\end{equation}
这个结论只排除了 \(1\) 处的谱原子。若 \(D(A)\ne H\)，则 \(I-U\) 的值域是稠密真子空间，
因而 \(1\) 仍属于 \(U\) 的谱。逆变换 \(\theta\) 在该点附近无界，
这正是需要引入无界谱积分的地方。

\subsection{无界自伴谱定理}
\label{b1:subsec:290602}

先把圆上的投影值测度推到实轴：对 Borel 集 \(B\subseteq\mathbb R\)，定义
\begin{equation}\label{b1:eq:29-real-spectral-measure}
 E(B)=F\bigl(\theta^{-1}(B)\bigr).
\end{equation}
原像保持交与不交并，故 \(E\) 保持投影乘法和强可数可加性。
由\eqref{b1:eq:29-cayley-no-atom}，又有
\(E(\mathbb R)=F(\mathbb T\setminus\{1\})=I\)，所以 \(E\) 是投影值测度。
若 \(h\) 为实轴上的有界 Borel 函数，在 \(z=1\) 处把 \(h\circ\theta\) 任意补成有限值，则
\begin{equation}\label{b1:eq:29-cayley-pushforward}
 h(E)=\int_{\mathbb T}h(\theta(z))\,\mathrm dF(z).
\end{equation}
对示性函数，这就是定义；对简单函数，利用线性性；对一般有界 Borel 函数，
用一致有界的简单函数逐点逼近，再用有界谱积分的强收敛结论即可。

对于尚未定义的坐标积分，先取有界截断
\[
 \tau_n(t)=t\cdot\mathbf1_{[-n,n]}(t),\qquad A_n=\tau_n(E).
\]
若 \(\int t^2\,\mathrm d\mu_x(t)<\infty\)，则对 \(m\ge n\)，
\[
 \|(A_m-A_n)x\|^2
 =\int_{\{n<|t|\le m\}}t^2\,\mathrm d\mu_x(t)\longrightarrow0.
\]
因此 \(A_nx\) 在 \(H\) 中有极限。下面证明，这个平方可积条件恰好刻画原来的定义域。

\begin{theorem}[无界自伴谱定理]
\label{b1:thm:29-unbounded-spectral}
设 \(H\) 为复 Hilbert 空间，\(A:D(A)\subseteq H\to H\) 为稠密定义的自伴算子。
则在 \(\mathbb R\) 的 Borel 集上存在唯一的投影值测度 \(E\)，使
\begin{equation}\label{b1:eq:29-spectral-domain}
 D(A)=\left\{x\in H:\int_{\mathbb R}t^2\,\mathrm d\mu_x(t)<\infty\right\},
 \qquad \mu_x(B)=\langle E(B)x,x\rangle,
\end{equation}
并且
\begin{equation}\label{b1:eq:29-spectral-action}
 Ax=\lim_{n\to\infty}\int_{[-n,n]}t\,\mathrm dE(t)x,
 \qquad
 \|Ax\|^2=\int_{\mathbb R}t^2\,\mathrm d\mu_x(t)
 \quad(x\in D(A)).
\end{equation}
对每个 \(w\in\mathbb C\setminus\mathbb R\)，\(A-wI:D(A)\to H\) 为双射，其逆满足
\begin{equation}\label{b1:eq:29-unbounded-resolvent}
 (A-wI)^{-1}=\int_{\mathbb R}\frac1{t-w}\,\mathrm dE(t),\qquad
 \|(A-wI)^{-1}\|\le\frac1{|\operatorname{Im}w|}.
\end{equation}
\end{theorem}
\begin{proof}
若 \(H=\{0\}\)，则零投影值测度直接给出全部结论。以下设 \(H\ne\{0\}\)。
取由\eqref{b1:eq:29-real-spectral-measure}构造的 \(E\)。
先证明 \(D(A)\) 中的向量满足平方可积条件。
由\eqref{b1:eq:29-cayley-domain}，可把任意 \(x\in D(A)\) 唯一写成
\(x=(I-U)y\)。记 \(\nu_y(C)=\langle F(C)y,y\rangle\)。
对 Borel 集 \(C\subseteq\mathbb T\)，有界谱积分的乘法与范数公式给出
\[
 \nu_x(C)=\|F(C)(I-U)y\|^2
 =\int_C|1-z|^2\,\mathrm d\nu_y(z).
\]
推前标量测度，再用\eqref{b1:eq:29-cayley-weight}，得到
\begin{equation}\label{b1:eq:29-cayley-domain-norm}
 \int_{\mathbb R}(1+t^2)\,\mathrm d\mu_x(t)
 =\int_{\mathbb T\setminus\{1\}}
   (1+|\theta(z)|^2)\cdot|1-z|^2\,\mathrm d\nu_y(z)
 =4\cdot\|y\|^2.
\end{equation}
这里无界非负函数的换元可先对简单函数验证，再由单调收敛定理得到。
特别地，\(x\) 满足所需的平方可积条件。另一方面，
\[
 A_nx
 =i\int_{\{z\ne1:\,|\theta(z)|\le n\}}
       (1+z)\,\mathrm dF(z)y
 \longrightarrow i(I+U)y=Ax.
\]
收敛来自有界谱积分的支配收敛，因为乘数的绝对值至多为 \(2\)，
而这些集合递增覆盖 \(\mathbb T\setminus\{1\}\)。

反过来，设 \(x\in H\) 满足 \(\int t^2\,\mathrm d\mu_x<\infty\)，令
\[
 y_n=\left(\frac{t+i}{2i}\cdot\mathbf1_{[-n,n]}(t)\right)(E)x.
\]
由于 \(|(t+i)/(2i)|^2=(1+t^2)/4\) 对 \(\mu_x\) 可积，
同样的范数差估计表明 \((y_n)\) 为 Cauchy 列。设其极限为 \(y\)。
由\eqref{b1:eq:29-cayley-pushforward}和 \(F(\{1\})=0\)，有 \(U=\phi(E)\)，故
\[
 (I-U)y_n
 =\left(\left(1-\frac{t-i}{t+i}\right)
          \frac{t+i}{2i}\cdot\mathbf1_{[-n,n]}(t)\right)(E)x
 =E([-n,n])x\longrightarrow x.
\]
\(I-U\) 有界，因而 \((I-U)y=x\)。
由其值域等于 \(D(A)\)，得到 \(x\in D(A)\)。这证明了定义域等式和作用公式。
最后，\(A_nx\to Ax\) 及单调收敛给出
\[
 \|Ax\|^2=\lim_n\|A_nx\|^2
 =\lim_n\int_{[-n,n]}t^2\,\mathrm d\mu_x
 =\int_{\mathbb R}t^2\,\mathrm d\mu_x.
\]

现在取 \(w\notin\mathbb R\)，置 \(r_w(t)=(t-w)^{-1}\)、\(R=r_w(E)\)。
有
\[
 |r_w(t)|\le\frac1{|\operatorname{Im}w|},\qquad
 |t r_w(t)|\le1+\frac{|w|}{|\operatorname{Im}w|}.
\]
对 \(v\in H\)，有界谱积分给出的标量测度关系为
\[
 \mu_{Rv}(B)=\int_B|r_w(t)|^2\,\mathrm d\mu_v(t).
\]
因此 \(\int t^2\,\mathrm d\mu_{Rv}<\infty\)，即 \(Rv\in D(A)\)。
再由坐标截断和有界谱积分的强收敛，
\[
 ARv=\lim_n(\tau_n r_w)(E)v=(t r_w)(E)v,
 \qquad (A-wI)Rv=v.
\]
对 \(x\in D(A)\)，\(\langle Ax,x\rangle\) 为实数，故
\[
 |\operatorname{Im}w|\cdot\|x\|^2
 =\bigl|\operatorname{Im}\langle(A-wI)x,x\rangle\bigr|
 \le\|(A-wI)x\|\cdot\|x\|.
\]
所以 \(A-wI\) 单射。结合已经证明的右逆，它为双射，且逆的范数满足所述估计。

最后证明谱测度的唯一性。设另一个实轴投影值测度 \(E'\) 也满足
\eqref{b1:eq:29-spectral-domain}和\eqref{b1:eq:29-spectral-action}。
上面的预解算子论证只用了这两个性质以及有界谱积分规则，所以同样适用于 \(E'\)。
在 \(w=-i\) 处得到
\[
 (A+iI)^{-1}=\left(\frac1{t+i}\right)(E'),\qquad U=\phi(E').
\]
在单位圆上定义 \(F'(C)=E'(\phi^{-1}(C))\)。
它是投影值测度，满足 \(F'(\{1\})=0\)，而且
\(U=\int_{\mathbb T}z\,\mathrm dF'(z)\)。
酉谱定理的唯一性给出 \(F'=F\)。
对任意 Borel 集 \(B\subseteq\mathbb R\)，\(\theta^{-1}(B)=\phi(B)\)
是 \(\mathbb T\setminus\{1\}\) 的 Borel 集，故
\[
 E'(B)=F'(\theta^{-1}(B))=F(\theta^{-1}(B))=E(B).
\]
\end{proof}

这个定理中的唯一性同时包含定义域；只在一个较小的稠密子空间上给出积分公式，
不能代替\eqref{b1:eq:29-spectral-domain}的等式。
与有界情形一样，记
\[
 \rho(A)=\{w\in\mathbb C:A-wI:D(A)\to H
       \text{ 为双射且逆在 }H\text{ 上有界}\},\qquad
 \sigma(A)=\mathbb C\setminus\rho(A).
\]
由于 \(A\) 闭，双射时逆的有界性也可由闭图像定理推出。
\eqref{b1:eq:29-unbounded-resolvent}说明 \(\sigma(A)\subseteq\mathbb R\)。
谱测度还精确说明哪些实数属于谱。

\begin{proposition}[谱测度的支集与谱]
\label{b1:prop:29-unbounded-spectral-support}
设 \(H\ne\{0\}\) 为复 Hilbert 空间，\(A:D(A)\subseteq H\to H\) 为稠密定义的自伴算子，
\(E\) 为\cref{b1:thm:29-unbounded-spectral}给出的实轴谱测度。则
\[
 \sigma(A)=\{s\in\mathbb R:
 E((s-\varepsilon,s+\varepsilon))\ne0\text{ 对每个 }\varepsilon>0\text{ 成立}\}.
\]
这个集合为非空闭集，且 \(E(\sigma(A))=I\)。
\end{proposition}
\begin{proof}
把右端记为 \(K\)。它的补集为开集。
若 \(s\notin K\)，取 \(\varepsilon>0\) 使 \(E((s-\varepsilon,s+\varepsilon))=0\)，并定义
\[
 r(t)=
 \begin{cases}
  (t-s)^{-1},&|t-s|\ge\varepsilon,\\
  0,&|t-s|<\varepsilon.
 \end{cases}
\]
\(r\) 与 \(tr\) 都有界。与预解公式的证明相同，标量测度范数公式先给出
\(r(E)H\subseteq D(A)\)，坐标截断再给出
\[
 (A-sI)r(E)v=((t-s)r)(E)v=v\qquad(v\in H).
\]
对 \(x\in D(A)\)，还可在另一侧计算
\[
 \begin{aligned}
 r(E)(A-sI)x
 &=\lim_n r(E)(A_n-sI)x\\
 &=\bigl(tr-sr\bigr)(E)x=x.
 \end{aligned}
\]
第二行的极限成立，是因为 \(r\tau_n\) 一致有界并逐点趋于 \(rt\)，
而 \((t-s)r=1\) 在 \(E\)-几乎处处成立。因此 \(s\in\rho(A)\)。

若 \(s\in K\)，则对每个 \(0<\varepsilon<1\)，可在
\(E((s-\varepsilon,s+\varepsilon))H\) 中选一个单位向量 \(x_\varepsilon\)。
它的标量谱测度集中在这个有界区间上，故 \(x_\varepsilon\in D(A)\)，且
\[
 \|(A-sI)x_\varepsilon\|^2
 =\int_{(s-\varepsilon,s+\varepsilon)}|t-s|^2
          \,\mathrm d\mu_{x_\varepsilon}(t)
 \le\varepsilon^2.
\]
这里在有界区间上，\(A\) 的作用就是有界坐标积分。
若 \((A-sI)^{-1}\) 存在且有界，就会有
\(1\le\|(A-sI)^{-1}\|\cdot\varepsilon\)，与 \(\varepsilon\downarrow0\) 矛盾。
所以 \(s\in\sigma(A)\)，得到 \(K=\sigma(A)\)。

对每个 \(s\notin K\)，可在其一个零投影邻域内选取包含 \(s\) 的有理端点开区间。
所有这些区间构成可数族，覆盖 \(\mathbb R\setminus K\)，且各自的 \(E\)-投影为零。
对每个 \(x\in H\)，标量测度的可数次可加性给出
\(\mu_x(\mathbb R\setminus K)=0\)，所以 \(E(\mathbb R\setminus K)=0\)、\(E(K)=I\)。
若 \(K\) 为空，这会与 \(E(\mathbb R)=I\ne0\) 矛盾。
\end{proof}

\subsection{无界函数演算的定义域}
\label{b1:subsec:290603}

现在固定实轴上的投影值测度 \(E\)。下面的构造只依赖 \(E\)，
因而也可以用于尚未指定自伴算子的任意这种投影值测度。
所谓 \(E\)-几乎处处，是指除去一个满足 \(E(N)=0\) 的 Borel 集 \(N\)；
这样的集合对每个 \(\mu_x\) 都为零测集。
本节只考虑处处取有限复数值的 Borel 函数，允许函数无界。

\begin{definition}[无界谱积分与最大定义域]
\label{b1:def:29-unbounded-calculus}
设 \(H\) 为复 Hilbert 空间，\(E\) 为实轴 Borel 集上的投影值测度，
\(\mu_x(B)=\langle E(B)x,x\rangle\)。设 \(f:\mathbb R\to\mathbb C\) 为有限值 Borel 函数，记
\[
 C_n=\{t:|f(t)|\le n\},\qquad P_n=E(C_n),\qquad
 f_n=f\cdot\mathbf1_{C_n}.
\]
定义谱积分 \(T_f=\int f\,\mathrm dE\) 的\term[zuida-dingyiyu]{最大定义域}[maximal domain]为
\begin{equation}\label{b1:eq:29-borel-maximal-domain}
 D_f=D(T_f)=\left\{x\in H:\int_{\mathbb R}|f(t)|^2\,\mathrm d\mu_x(t)<\infty\right\},
\end{equation}
并对 \(x\in D_f\) 置
\begin{equation}\label{b1:eq:29-borel-truncation}
 T_fx=\lim_{n\to\infty}f_n(E)x.
\end{equation}
若 \(E\) 为自伴算子 \(A\) 的谱测度，则记
\[
 f(A)=T_f=\int_{\mathbb R}f(t)\,\mathrm dE(t).
\]
\symidx{\(D(f(A))\)：无界谱函数演算的最大定义域}
\end{definition}

先验证定义确实给出稠密定义的线性算子。
由有界谱积分的范数公式及单调收敛，
\begin{equation}\label{b1:eq:29-borel-domain-test}
 \sup_n\|f_n(E)x\|^2=\int|f|^2\,\mathrm d\mu_x.
\end{equation}
所以 \(D_f\) 恰为使左侧有限的向量全体，算子范数的三角不等式说明它为线性子空间。
对 \(x\in D_f\)、\(m\ge n\)，有
\[
 \|(f_m(E)-f_n(E))x\|^2
 =\int_{\{n<|f|\le m\}}|f|^2\,\mathrm d\mu_x\longrightarrow0,
\]
故\eqref{b1:eq:29-borel-truncation}中的极限存在。
各个截断算子的线性性再说明 \(T_f\) 线性，且
\begin{equation}\label{b1:eq:29-borel-norm}
 \|T_fx\|^2=\int|f|^2\,\mathrm d\mu_x\qquad(x\in D_f).
\end{equation}
对任意 \(y\in H\)，\(P_ny\) 的标量谱测度集中在 \(C_n\) 上，因而
\(\int|f|^2\,\mathrm d\mu_{P_ny}\le n^2\cdot\|y\|^2\)，即 \(P_nH\subseteq D_f\)。
\(f\) 处处有限意味着 \(C_n\uparrow\mathbb R\)，从而 \(P_n\to I\) 强收敛，
\(D_f\) 在 \(H\) 中稠密。

若 \(f\) 有界，则充分大的 \(n\) 满足 \(f_n=f\)，所以此定义与原来的有界谱积分一致。
若 \(f\) 和 \(g\) 在 \(E\)-几乎处处相等，则各自的平方可积域相同，
截断函数的有界谱积分也相同，故 \(T_f=T_g\)。
对于坐标函数 \(f(t)=t\)，\cref{b1:thm:29-unbounded-spectral}说明得到的算子恰为 \(A\)，
包括定义域的等式。

\begin{proposition}[谱投影与标量谱测度的变换]
\label{b1:prop:29-unbounded-spectral-commutation}
设 \(H\) 为复 Hilbert 空间，\(E\) 为实轴 Borel 集上的投影值测度，
\(f\) 为有限值 Borel 函数，\(T_f=\int f\,\mathrm dE\) 取最大定义域 \(D_f\)。
设 \(B\subseteq\mathbb R\) 为 Borel 集。
则对 \(x\in D_f\)，有
\[
 E(B)x\in D_f,\qquad T_fE(B)x=E(B)T_fx,
\]
且
\begin{equation}\label{b1:eq:29-unbounded-scalar-transport}
 \mu_{T_fx}(B)=\int_B|f|^2\,\mathrm d\mu_x.
\end{equation}
若 \(f\) 在 \(B\) 上有界，则 \(E(B)H\subseteq D_f\)，并且作为整个 \(H\) 上的有界算子，
\[
 T_fE(B)=(f\cdot\mathbf1_B)(E).
\]
\end{proposition}
\begin{proof}
投影乘法给出
\(\mu_{E(B)x}(C)=\mu_x(B\cap C)\)。因此
\[
 \int|f|^2\,\mathrm d\mu_{E(B)x}=\int_B|f|^2\,\mathrm d\mu_x<\infty.
\]
每个有界截断 \(f_n(E)\) 都与 \(E(B)\) 交换，令 \(n\to\infty\) 即得作用公式。
再用\eqref{b1:eq:29-borel-norm}，
\[
 \mu_{T_fx}(B)=\|E(B)T_fx\|^2
 =\|T_fE(B)x\|^2=\int_B|f|^2\,\mathrm d\mu_x.
\]
若 \(|f|\le M\) 于 \(B\)，上述积分对任意 \(x\in H\) 至多为 \(M^2\cdot\|x\|^2\)。
并且对 \(n\ge M\)，\(f_n\cdot\mathbf1_B=f\cdot\mathbf1_B\)，所以截断定义直接给出最后的等式。
\end{proof}

\begin{theorem}[无界谱积分的闭性与伴随]
\label{b1:thm:29-unbounded-calculus-adjoint}
设 \(H\) 为复 Hilbert 空间，\(E\) 为实轴 Borel 集上的投影值测度，
\(f\) 为有限值 Borel 函数，\(T_f=\int f\,\mathrm dE\) 取最大定义域 \(D_f\)。
则 \(T_f\) 是稠密定义的闭算子，且
\begin{equation}\label{b1:eq:29-unbounded-calculus-adjoint}
 T_f^*=T_{\overline f},\qquad
 D(T_f^*)=D(T_{\overline f})=D_f.
\end{equation}
特别地，若 \(f\) 在 \(E\)-几乎处处为实值，则 \(T_f\) 自伴。
\end{theorem}
\begin{proof}
稠密定义已经验证。先直接证明闭性。
设 \(x_k\in D_f\)、\(x_k\to x\)、\(T_fx_k\to y\)。
固定 \(n\)，由上一命题有
\[
 f_n(E)x_k=P_nT_fx_k.
\]
两侧有界算子均可通过极限，得到 \(f_n(E)x=P_ny\)。于是
\(\|f_n(E)x\|\le\|y\|\) 对所有 \(n\) 成立。
\eqref{b1:eq:29-borel-domain-test}说明 \(x\in D_f\)，再令 \(n\to\infty\)，得到
\(T_fx=\lim_nP_ny=y\)，证明了闭性。

由于 \(|\overline f|=|f|\)，有 \(D_{\overline f}=D_f\)。
对 \(x\in D_f\)、\(y\in D_{\overline f}\)，有界谱积分的伴随公式给出
\[
 \langle T_fx,y\rangle
 =\lim_n\langle f_n(E)x,y\rangle
 =\lim_n\langle x,\overline{f_n}(E)y\rangle
 =\langle x,T_{\overline f}y\rangle.
\]
因此 \(T_{\overline f}\subseteq T_f^*\)。

反过来，设 \(y\in D(T_f^*)\)、\(v=T_f^*y\)。
对任意 \(u\in H\)，\(P_nu\in D_f\)，所以
\[
 \langle f_n(E)u,y\rangle
 =\langle T_fP_nu,y\rangle
 =\langle P_nu,v\rangle
 =\langle u,P_nv\rangle.
\]
由有界谱积分的伴随公式及表示向量唯一性，
\(\overline{f_n}(E)y=P_nv\)。因此
\[
 \int_{C_n}|f|^2\,\mathrm d\mu_y
 =\|\overline{f_n}(E)y\|^2\le\|v\|^2.
\]
单调收敛给出 \(y\in D_{\overline f}\)。令 \(n\to\infty\)，再得到
\(T_{\overline f}y=v\)，从而伴随等式连同定义域等式全部成立。
最后，若 \(f=\overline f\) 于 \(E\)-几乎处处，两个谱积分相同，故 \(T_f=T_f^*\)。
\end{proof}

上述闭性说明最大定义域对图像极限封闭。
若再把 \(T_f\) 限制到一个较小的稠密子空间，所得算子可能不闭，也可能不再自伴。
因此，谱积分中的平方可积条件是算子本身的一部分。

\subsection{谱演算的乘积与极限}
\label{b1:subsec:290604}

以下把记号恢复为 \(f(A)\)。所有函数都相对于同一个自伴算子 \(A\) 的谱测度计算。
算子乘积先作用右侧因子，约定
\[
 D(f(A)g(A))
 =\{x\in D(g(A)):g(A)x\in D(f(A))\}.
\]
算子和的定义域则是 \(D(f(A))\cap D(g(A))\)。

\begin{proposition}[乘积和加法的定义域]
\label{b1:prop:29-unbounded-calculus-products}
设 \(H\) 为复 Hilbert 空间，\(A:D(A)\subseteq H\to H\) 为稠密定义的自伴算子，
\(f,g:\mathbb R\to\mathbb C\) 为有限值 Borel 函数，函数演算均取最大定义域。则
\begin{equation}\label{b1:eq:29-product-domain}
 D(f(A)g(A))=D(g(A))\cap D((fg)(A)),
\end{equation}
并且
\begin{equation}\label{b1:eq:29-product-restriction}
 f(A)g(A)x=(fg)(A)x
 \qquad\bigl(x\in D(f(A)g(A))\bigr).
\end{equation}
对于加法，有
\begin{equation}\label{b1:eq:29-sum-restriction}
 D(f(A))\cap D(g(A))\subseteq D((f+g)(A)),\qquad
 (f(A)+g(A))x=(f+g)(A)x
\end{equation}
在左端定义域上成立。上述两个算子限制均可闭，且
\begin{equation}\label{b1:eq:29-product-sum-closures}
 \overline{f(A)g(A)}=(fg)(A),\qquad
 \overline{f(A)+g(A)}=(f+g)(A).
\end{equation}
\end{proposition}
\begin{proof}
对 \(x\in D(g(A))\)，由\eqref{b1:eq:29-unbounded-scalar-transport}，先对简单非负函数，
再用单调收敛，可得
\[
 \int|f|^2\,\mathrm d\mu_{g(A)x}
 =\int|fg|^2\,\mathrm d\mu_x.
\]
所以 \(g(A)x\in D(f(A))\) 当且仅当 \(x\in D((fg)(A))\)，证明了精确的乘积域。

令
\[
 B_n=\{|f|\le n,\ |g|\le n\},\qquad Q_n=E(B_n).
\]
两函数都取有限值，故 \(B_n\uparrow\mathbb R\)、\(Q_n\to I\) 强收敛。
在 \(B_n\) 上两个函数有界，可用有界谱积分的乘法公式。
对乘积域中的 \(x\)，上一小节的谱投影交换关系给出
\[
 \begin{aligned}
 Q_nf(A)g(A)x
 &=f(A)Q_ng(A)x\\
 &=(f\cdot\mathbf1_{B_n})(E)(g\cdot\mathbf1_{B_n})(E)x\\
 &=(fg\cdot\mathbf1_{B_n})(E)x
 =Q_n(fg)(A)x.
 \end{aligned}
\]
令 \(n\to\infty\) 得\eqref{b1:eq:29-product-restriction}。

若 \(x\in D(f(A))\cap D(g(A))\)，由
\(|f+g|^2\le2\cdot|f|^2+2\cdot|g|^2\)，首先有 \(x\in D((f+g)(A))\)。
再在同一 \(B_n\) 上使用有界谱积分的加法公式，得到
\[
 Q_n(f(A)x+g(A)x)
 =((f+g)\cdot\mathbf1_{B_n})(E)x
 =Q_n(f+g)(A)x.
\]
令 \(n\to\infty\)，得到加法的作用公式。

最后，这两个限制都包含 \(\bigcup_nQ_nH\) 于定义域中，故稠密定义；
又分别包含于闭算子 \((fg)(A)\) 与 \((f+g)(A)\)，因而可闭。
分别取 \(h=fg\) 或 \(h=f+g\)。对每个 \(x\in D(h(A))\)，置 \(x_n=Q_nx\)。
则 \(x_n\) 属于相应限制的定义域，并且
\[
 \|x-x_n\|^2+\|h(A)x-h(A)x_n\|^2
 =\int_{B_n^c}(1+|h|^2)\,\mathrm d\mu_x\longrightarrow0.
\]
这里被积函数对 \(\mu_x\) 可积，且 \(B_n\uparrow\mathbb R\)。
所以每个 \((x,h(A)x)\) 都在相应限制的图像闭包中。
反向包含由 \(h(A)\) 的闭性保证，证明了两个闭包等式。
\end{proof}

由\eqref{b1:eq:29-product-domain}，当右侧因子 \(g\) 有界时，确有
\(f(A)g(A)=(fg)(A)\)。仅知道左侧因子 \(f\) 有界，则一般只能得到扩张关系
\(f(A)g(A)\subseteq(fg)(A)\)。加法中的域包含也可能为严格包含。

在 \(H=\ell^2(\mathbb N)\) 上，可以把这两个现象直接写出。
定义投影值测度
\[
 E(B)(x_1,x_2,\ldots)
 =(\mathbf1_B(1)\cdot x_1,\mathbf1_B(2)\cdot x_2,\ldots).
\]
其强可数可加性由每个向量的平方可和性及标量支配收敛得到，且
\(\mu_x=\sum_{k\ge1}|x_k|^2\cdot\delta_k\)。
对坐标函数应用上一小节的谱积分，得到自伴算子
\[
 Mx=(kx_k)_{k\ge1},\qquad
 D(M)=\left\{x\in\ell^2:\sum_{k\ge1}k^2\cdot|x_k|^2<\infty\right\}.
\]
由无界自伴谱定理的唯一性，它的谱测度正是这个 \(E\)。
若 \(g(t)=t\)、\(f(t)=(1+|t|)^{-1}\)，则 \(f(M)\) 与 \((fg)(M)\) 都在整个空间上有界，
而 \(f(M)g(M)\) 的定义域仍为 \(D(M)\)。
例如 \(x=(1/k)_{k\ge1}\) 属于 \(\ell^2\) 却不属于 \(D(M)\)，所以这两个算子不相等。
若改取 \(f(t)=t\)、\(g(t)=-t\)，则 \(f(M)+g(M)\) 是定义在 \(D(M)\) 上的零算子，
而 \((f+g)(M)\) 是定义在整个 \(\ell^2\) 上的零算子。

\ProseParagraphBreak

函数列的点态极限也要连同可积的控制条件使用。

\begin{proposition}[谱积分的支配收敛]
\label{b1:prop:29-unbounded-calculus-dominated}
设 \(H\) 为复 Hilbert 空间，\(A:D(A)\subseteq H\to H\) 为稠密定义的自伴算子，
\(E\) 为其谱测度。设 \(f_n,f\) 为有限值 Borel 函数，\(h\ge0\) 为有限值 Borel 函数。
若在一个共同的 \(E\)-零测集之外，
\[
 f_n(t)\longrightarrow f(t),\qquad |f_n(t)|\le h(t)\quad(n\ge1),
\]
则对每个 \(x\in D(h(A))\)，均有 \(x\in D(f_n(A))\cap D(f(A))\)，且
\[
 f_n(A)x\longrightarrow f(A)x\quad\text{于 }H.
\]
\end{proposition}
\begin{proof}
取极限给出 \(|f|\le h\) 于共同零测集之外，故
\[
 \int|f_n|^2\,\mathrm d\mu_x\le\int h^2\,\mathrm d\mu_x<\infty,
 \qquad
 \int|f|^2\,\mathrm d\mu_x\le\int h^2\,\mathrm d\mu_x<\infty.
\]
由加法公式及范数公式，
\[
 \|f_n(A)x-f(A)x\|^2
 =\int|f_n-f|^2\,\mathrm d\mu_x.
\]
被积函数逐点趋于零，且由 \(4h^2\) 控制。
标量支配收敛定理因此给出右端趋于零。
\end{proof}

这个结论是对固定向量的范数收敛。
即使所有函数都一致有界，也不能据此推出算子范数收敛。
例如，对上面的 \(M\) 取 \(f_n=\mathbf1_{(n,\infty)}\)，有 \(f_n(M)x\to0\) 对每个 \(x\in\ell^2\) 成立，
但 \(f_n(M)e_{n+1}=e_{n+1}\)，所以 \(\|f_n(M)\|=1\) 对所有 \(n\) 成立。

最后把谱截断表述为定义域中的逼近。
这里使用图像范数下稠密的定义子空间（英文称 core）：对闭算子 \(T\)，
它指在 \(D(T)\) 中关于图像范数 \(\|x\|_T=\|x\|+\|Tx\|\) 稠密的线性子空间
\(\mathcal C\subseteq D(T)\)。
这等价于 \(T|_{\mathcal C}\) 的图像闭包等于 \(T\) 的图像。

\begin{proposition}[谱截断的图像范数逼近]
\label{b1:prop:29-spectral-graph-core}
设 \(H\) 为复 Hilbert 空间，\(A:D(A)\subseteq H\to H\) 为稠密定义的自伴算子，
\(E\) 为其谱测度，\(f\) 为有限值 Borel 函数，令 \(P_n=E(\{|f|\le n\})\)。
则 \(\mathcal C_f=\bigcup_{n\ge1}P_nH\) 在 \(D(f(A))\) 中关于图像范数稠密。
更具体地，对每个 \(x\in D(f(A))\)，有
\begin{equation}\label{b1:eq:29-spectral-graph-approximation}
 \|x-P_nx\|^2+
 \|f(A)x-f(A)P_nx\|^2
 =\int_{\{|f|>n\}}(1+|f|^2)\,\mathrm d\mu_x\longrightarrow0.
\end{equation}
特别地，\(\bigcup_{n\ge1}E([-n,n])H\) 在 \(D(A)\) 中关于图像范数稠密，且
\[
 \|x-E([-n,n])x\|^2+
 \|Ax-AE([-n,n])x\|^2
 =\int_{\{|t|>n\}}(1+t^2)\,\mathrm d\mu_x(t)\longrightarrow0
 \quad(x\in D(A)).
\]
\end{proposition}
\begin{proof}
上一小节已经证明 \(P_nH\subseteq D(f(A))\)。这些子空间递增，故其并为线性子空间。
由谱投影交换关系，
\(f(A)x-f(A)P_nx=f(A)(I-P_n)x\)。
\(I-P_n=E(\{|f|>n\})\)，所以分别对恒等乘数和 \(f\) 使用范数公式，
得到\eqref{b1:eq:29-spectral-graph-approximation}中的等式。
因为 \(1+|f|^2\) 对 \(\mu_x\) 可积，尾部积分趋于零。
平方和收敛等价于图像范数收敛，因而 \(\mathcal C_f\) 在定义域中关于图像范数稠密。
最后取 \(f(t)=t\)，得到关于 \(A\) 的结论。
\end{proof}

对一般 Borel 函数，采用 \(\{|f|\le n\}\) 作为截断集，保证每个投影的整个值域都在定义域中。
函数未必在实轴的有界区间上有界，所以不能无条件把它改为 \([-n,n]\)。
对于坐标函数，这两种截断恰好一致；所得谱子空间上的算子有界，
而图像范数逼近保证这些有界部分足以恢复原来的无界自伴算子。

% !TeX root = ../../Book1.tex
\section{乘法算子与微分算子的定义域模型}
\label{b1:sec:2907}

% 实读：Rudin1991Functional，13.4（pp.349--350）的区间一阶微分算子、
% 周期域及两端零值域，并核对其伴随关系。原书使用+i d/dx；本节固定-i d/dx，
% 用本卷弱导数/一维绝对连续代表重证，再推广到全部相位边界条件。
% 最大域实乘法模型、截断图像逼近与C_c^infty闭包均在本节直接证明。

乘法与求导提供了两种不同的定义域问题。乘法要求输出仍平方可积；求导除了要求弱导数平方可积，
还可能带有端点条件。前一种要求直接确定最大允许输入，后一种要求则让同一条微分表达式产生不同算子。
本节逐个计算这些定义域及其伴随，把\secref{b1:sec:2905}的抽象条件落实到可以检验的模型。

\subsection{无界乘法算子}
\label{b1:subsec:290701}

\begin{theorem}[实可测函数的最大域乘法算子]
\label{b1:thm:29-real-multiplication}
设 \((X,\Sigma,\mu)\) 为测度空间，\(m:X\to\mathbb R\) 为可测函数，
\(H=L^2(X,\Sigma,\mu;\mathbb C)\)。定义
\[
 D(M_m)=\{f\in H:mf\in H\},\qquad M_mf=mf.
\]
则 \(M_m\) 稠密定义、闭且自伴。对每个 \(z\in\mathbb C\setminus\mathbb R\)，
\(M_m-zI:D(M_m)\to H\) 是双射，且
\begin{equation}\label{b1:eq:29-multiplication-resolvent}
 (M_m-zI)^{-1}g=\frac{g}{m-z},\qquad
 \|(M_m-zI)^{-1}\|\leq|\operatorname{Im}z|^{-1}.
\end{equation}
这里不要求 \(\mu\) 有限或 \(\sigma\)-有限。
\symidx{\(D(M_m)=\{f\in L^2:mf\in L^2\}\)：乘法算子的最大定义域}
\end{theorem}
\begin{proof}
乘积的可测性及几乎处处等价保证算子在 \(L^2\) 等价类上良定义，定义域的线性来自乘法的线性。
记 \(E_n=\{x:|m(x)|\leq n\}\)。对每个 \(f\in H\)，函数
\(f_n=\mathbf1_{E_n}\cdot f\) 满足 \(\|mf_n\|_2\leq n\cdot\|f\|_2\)，故属于定义域。
由于 \(m\) 为有限实值函数，\(E_n\uparrow X\)，控制收敛给出 \(f_n\to f\) 于 \(H\)。
因此 \(D(M_m)\) 稠密。

对 \(f,g\in D(M_m)\)，Cauchy--Schwarz 不等式保证下列积分绝对收敛，且实值性给出
\[
 \langle M_mf,g\rangle
 =\int_Xmf\overline g\,\mathrm d\mu
 =\int_Xf\overline{mg}\,\mathrm d\mu
 =\langle f,M_mg\rangle.
\]
所以 \(M_m\) 对称。

固定非实 \(z\)，对任意 \(g\in H\) 令 \(f=g/(m-z)\)。逐点估计
\[
 \frac1{|m-z|}\leq\frac1{|\operatorname{Im}z|},\qquad
 \frac{|m|}{|m-z|}\leq1+\frac{|z|}{|\operatorname{Im}z|}
\]
说明 \(f\in H\)、\(mf\in H\)，故 \(f\in D(M_m)\)，并且 \((M_m-zI)f=g\)。
由于 \(m-z\) 处处不为零，解唯一，第一个估计又给出逆算子的范数界。
特别地，两个参数 \(z=\mathrm i,-\mathrm i\) 下均满射。
\cref{b1:thm:29-self-adjoint-criterion}说明 \(M_m\) 自伴且闭。
\end{proof}

定义域包含了所有使乘积仍属于 \(L^2\) 的函数，因而称为乘法表达式的最大定义域。
若原函数只在几乎处处有限，可先在其取无穷值的可测零集上改成零，所得算子相同。
这里的“最大”是对固定乘法表达式而言；任意取一个较小的稠密子空间作为定义域，仍会得到它的另一个限制算子。

\ProseParagraphBreak

在 \(L^2(\mathbb R;\mathbb C)\) 上取 \(m(t)=t\)，就得到一个无界自伴算子。
单位向量 \(f_n=\mathbf1_{[n,n+1]}\) 都在定义域内，却满足 \(\|M_mf_n\|_2\geq n\)。
它也没有定义在整个 \(L^2\) 上：函数 \(f(t)=(1+|t|)^{-1}\) 平方可积，
而 \(tf(t)\) 不平方可积。闭性在这里控制同时收敛的输入与输出，没有给出全空间上的有界性。

\subsection{微分算子与边界配对}
\label{b1:subsec:290702}

以下固定
\[
 H=L^2((0,1),\mathrm dx;\mathbb C),\qquad
 \langle f,g\rangle=\int_0^1f(x)\overline{g(x)}\,\mathrm dx.
\]
\cref{b1:prop:sobolev-interval-representative}保证每个 \(H^1(0,1)\) 元素都有唯一的
\([0,1]\) 上绝对连续代表。下面的 \(f(0),f(1)\) 一律指这个代表的端点值。
该命题的估计及 Hölder 不等式给出
\begin{equation}\label{b1:eq:29-endpoint-estimate}
 \max_{0\leq x\leq1}|f(x)|\leq\|f\|_2+\|f'\|_2
 \qquad(f\in H^1(0,1)).
\end{equation}
所以端点条件对 \(H^1\) 范数连续，不能给任意 \(L^2\) 代表直接附加同样的点值条件。

\begin{proposition}[一阶微分算子的边界配对]
\label{b1:prop:29-derivative-boundary}
在复 Hilbert 空间 \(H=L^2(0,1)\) 上，定义
\[
 D(P_{\max})=H^1(0,1),\qquad P_{\max}f=-\mathrm if',
\]
其中导数为弱导数。则 \(P_{\max}\) 是稠密定义的闭算子，其图像范数为
\(\|f\|_{P_{\max}}=\|f\|_2+\|f'\|_2\)，与 \(H^1\) 范数等价。
对所有 \(f,g\in H^1(0,1)\)，有
\begin{equation}\label{b1:eq:29-derivative-boundary}
 \langle P_{\max}f,g\rangle-\langle f,P_{\max}g\rangle
 =-\mathrm i\bigl(f(1)\overline{g(1)}-f(0)\overline{g(0)}\bigr).
\end{equation}
\symidx{\(P_{\max}\)：定义域为\(H^1(0,1)\)的\(-\mathrm i\,\mathrm d/\mathrm dx\)}
\end{proposition}
\begin{proof}
先说明 \(C_c^\infty(0,1)\) 在 \(H\) 中稠密。若 \(h\in H\) 与所有这类测试函数正交，
则 \(\int_0^1\varphi\overline h=0\) 对每个测试函数成立。
由正则分布的单射性\cref{b1:prop:regular-distribution-injective}，有 \(h=0\) 几乎处处。
双重正交补公式因此给出稠密性。于是包含这一空间的 \(D(P_{\max})\) 也稠密。

若 \(f_n\to f\)、\(-\mathrm if_n'\to h\) 都在 \(L^2\) 中成立，
则对任意 \(\varphi\in C_c^\infty(0,1)\)，可以在弱导数公式中取极限，得到
\[
 \int_0^1 f\varphi'
 =-\int_0^1\mathrm ih\varphi.
\]
因此 \(f'=\mathrm ih\in L^2\)，即 \(f\in H^1(0,1)\)、\(P_{\max}f=h\)，证明闭性。
图像范数公式来自 \(\|-\mathrm if'\|_2=\|f'\|_2\)，其与平方和范数的等价性已在
\secref{b1:sec:2905}说明。

最后，对所选绝对连续代表，乘积 \(f\overline g\) 绝对连续，导数为
\(f'\overline g+f\overline{g'}\)。这些项可积，因为 \(f,g\) 在闭区间上有界，
\(f',g'\in L^2(0,1)\subseteq L^1(0,1)\)。积分微积分基本定理给出
\[
 \int_0^1\bigl(f'\overline g+f\overline{g'}\bigr)\,\mathrm dx
 =f(1)\overline{g(1)}-f(0)\overline{g(0)}.
\]
乘以 \(-\mathrm i\) 即得边界配对公式，符号使用了内积第一变量线性的约定。
\end{proof}

把同一表达式先限制在内部测试函数上，伴随就会恢复全部允许的弱导数定义域。

\begin{proposition}[测试函数域的伴随]
\label{b1:prop:29-minimal-adjoint}
在 \(H=L^2((0,1),\mathrm dx;\mathbb C)\) 上，令
\[
 D(P_c)=C_c^\infty(0,1),\qquad P_cf=-\mathrm if'.
\]
则 \(P_c\) 稠密定义且对称，并满足
\[
 P_c^*=P_{\max},\qquad
 D(P_{\max}^*)=H_0^1(0,1),\qquad P_{\max}^*g=-\mathrm ig'.
\]
这里 \(P_{\max}\) 的定义域为 \(H^1(0,1)\)，
\(H_0^1(0,1)=\{g\in H^1(0,1):g(0)=g(1)=0\}\)。
\symidx{\(P_c\)：定义在\(C_c^\infty(0,1)\)上的\(-\mathrm i\,\mathrm d/\mathrm dx\)}
\end{proposition}
\begin{proof}
稠密性已在前一命题中证明。测试函数的两端值均为零，所以边界配对也给出对称性。
若 \(g\in H^1(0,1)\)，对 \(f\in C_c^\infty(0,1)\) 使用边界公式，得到
\(\langle P_cf,g\rangle=\langle f,-\mathrm ig'\rangle\)。
因此 \(g\in D(P_c^*)\)、\(P_c^*g=-\mathrm ig'\)。

反之，若 \(g\in D(P_c^*)\)、\(P_c^*g=h\in L^2\)，伴随关系写成
\[
 -\mathrm i\int_0^1\varphi'\overline g\,\mathrm dx
 =\int_0^1\varphi\overline h\,\mathrm dx
 \qquad(\varphi\in C_c^\infty(0,1)).
\]
取复共轭，再把 \(\overline\varphi\) 换记为任意测试函数 \(\psi\)，得到
\[
 \int_0^1g\psi'\,\mathrm dx=-\mathrm i\int_0^1h\psi\,\mathrm dx.
\]
按弱导数的定义，\(g'=\mathrm ih\in L^2\)，故 \(g\in H^1\)、\(-\mathrm ig'=h\)。
这证明 \(P_c^*=P_{\max}\)，包括定义域的相等。

又因 \(P_c\subseteq P_{\max}\)，扩张反向性给出
\(P_{\max}^*\subseteq P_c^*=P_{\max}\)。因此 \(P_{\max}^*\) 的定义域向量已属于 \(H^1\)。
对这样的 \(g\)，属于伴随定义域还要求
\[
 f(1)\overline{g(1)}-f(0)\overline{g(0)}=0
 \qquad(f\in H^1(0,1)).
\]
仿射函数可以独立指定 \(f(0)\) 与 \(f(1)\)，故这等价于 \(g(0)=g(1)=0\)。
反向包含也由同一个边界公式得到。
最后，\cref{b1:thm:interval-zero-boundary}在 \(p=2\) 时把这一子空间识别为 \(H_0^1(0,1)\)。
\end{proof}

由此，最大域算子虽然闭，却不对称。例如 \(f(x)=x\) 使
\(\langle P_{\max}f,f\rangle-\langle f,P_{\max}f\rangle=-\mathrm i\ne0\)。
要使微分算子自伴，需要选择恰好让边界配对消失、且伴随不再扩大定义域的边界条件。

\subsection{边界条件与自伴性}
\label{b1:subsec:290703}

\begin{theorem}[相位边界条件的一阶自伴算子]
\label{b1:thm:29-phase-derivative}
在复 Hilbert 空间 \(H=L^2((0,1),\mathrm dx;\mathbb C)\) 上，固定 \(\theta\in\mathbb R\)，定义
\[
 D(P_\theta)=\{f\in H^1(0,1):f(1)=e^{\mathrm i\theta}f(0)\},
 \qquad P_\theta f=-\mathrm if'.
\]
其中端点值取唯一的绝对连续代表。则 \(P_\theta\) 稠密定义、闭且自伴。
它的特征值恰为 \(\theta+2\pi k\)（\(k\in\mathbb Z\)），各特征空间为
\[
 \operatorname{span}\{x\mapsto e^{\mathrm i(\theta+2\pi k)x}\}.
\]
\end{theorem}
\begin{proof}
定义域包含 \(C_c^\infty(0,1)\)，所以稠密。
对 \(f,g\in D(P_\theta)\)，有
\[
 f(1)\overline{g(1)}
 =e^{\mathrm i\theta}f(0)\cdot e^{-\mathrm i\theta}\overline{g(0)}
 =f(0)\overline{g(0)}.
\]
边界公式表明 \(P_\theta\) 对称。

从 \(P_c\subseteq P_\theta\subseteq P_{\max}\) 得到
\(P_\theta^*\subseteq P_c^*=P_{\max}\)。
所以对 \(g\in D(P_\theta^*)\)，已知 \(g\in H^1\)、\(P_\theta^*g=-\mathrm ig'\)。
伴随条件进一步要求
\[
 f(0)\bigl(e^{\mathrm i\theta}\overline{g(1)}-\overline{g(0)}\bigr)=0
 \qquad(f\in D(P_\theta)).
\]
对任意 \(c\in\mathbb C\)，仿射函数
\(f(x)=c(1-x)+e^{\mathrm i\theta}cx\) 属于定义域，且 \(f(0)=c\)。
因此括号内必须为零，取复共轭后即得 \(g(1)=e^{\mathrm i\theta}g(0)\)。
故 \(D(P_\theta^*)\subseteq D(P_\theta)\)。
对称性已给出反向包含，所以 \(P_\theta=P_\theta^*\)，闭性随之成立。

若 \(P_\theta f=\lambda f\)、\(f\ne0\)，自伴性保证 \(\lambda\in\mathbb R\)。
绝对连续代表满足 \(f'=\mathrm i\lambda f\) 几乎处处，故
\((e^{-\mathrm i\lambda x}f(x))'=0\) 几乎处处。
绝对连续性推出 \(f(x)=ce^{\mathrm i\lambda x}\)，其中 \(c\ne0\)。
端点条件等价于 \(e^{\mathrm i\lambda}=e^{\mathrm i\theta}\)，
也就是 \(\lambda=\theta+2\pi k\)。反过来，每个所列指数函数都在定义域中，并直接满足特征方程。
\end{proof}

参数 \(\theta=0\) 给出周期条件，\(\theta=\pi\) 给出 \(f(1)=-f(0)\)。
两者使用同一个微分表达式，特征值却分别为 \(2\pi k\) 与 \((2k+1)\pi\)。
因此端点条件确实改变算子及其谱信息。

\begin{corollary}[测试函数算子的全部自伴扩张]
\label{b1:cor:29-derivative-extensions}
在 \(H=L^2((0,1),\mathrm dx;\mathbb C)\) 上，令
\(P_c=-\mathrm i\,\mathrm d/\mathrm dx\)，\(D(P_c)=C_c^\infty(0,1)\)。
若 \(B\) 是 \(P_c\) 的自伴扩张，则存在 \(\theta\in[0,2\pi)\)，使
\[
 D(B)=\{f\in H^1(0,1):f(1)=e^{\mathrm i\theta}f(0)\},
 \qquad Bf=-\mathrm if'.
\]
\end{corollary}
\begin{proof}
由 \(P_c\subseteq B=B^*\)，扩张反向性给出 \(B\subseteq P_c^*=P_{\max}\)。
因此 \(D(B)\subseteq H^1\)，作用为 \(-\mathrm if'\)。
又因 \(B\) 闭，\(C_c^\infty\) 在 \(H^1\) 中的闭包 \(H_0^1\) 包含于 \(D(B)\)：
对每个 \(f\in H_0^1\) 选取 \(H^1\) 逼近列，其输入及导数同时收敛，再用 \(B\) 的闭性。

记 \(S=\{(f(0),f(1)):f\in D(B)\}\subseteq\mathbb C^2\)。
由于两个端点相同的 \(H^1\) 函数之差属于 \(H_0^1\)，有
\[
 D(B)=\{f\in H^1(0,1):(f(0),f(1))\in S\}.
\]
若 \(S=\{0\}\)，则 \(D(B)=H_0^1\)。边界配对说明每个 \(g\in H^1\) 都属于
\(D(B^*)\)，特别是常数函数 \(1\)；这与 \(B=B^*\) 矛盾。
故可取非零 \((a,b)\in S\)。对称性及边界配对给出 \(|b|^2=|a|^2\)，
于是 \(a\ne0\)，存在唯一 \(\theta\in[0,2\pi)\) 使 \(b=e^{\mathrm i\theta}a\)。
对任意 \((c,d)\in S\)，同样的配对给出 \(b\overline d-a\overline c=0\)，
所以 \(d=e^{\mathrm i\theta}c\)。
因此 \(S\) 正好是一维子空间 \(\{(c,e^{\mathrm i\theta}c):c\in\mathbb C\}\)，结论随之成立。
\end{proof}

\subsection{闭包、定义域与反例}
\label{b1:subsec:290704}

闭算子的定义域往往先由一个更便于计算的子空间逼近。这里需要同时逼近函数与算子值。
对复 Hilbert 空间上的闭算子 \(A:D(A)\subseteq H\to H\)，若线性子空间
\(\mathcal D\subseteq D(A)\) 关于图像范数稠密，即对每个 \(x\in D(A)\) 都存在
\(x_n\in\mathcal D\)，使
\[
 x_n\to x,\qquad Ax_n\to Ax,
\]
则英文文献称 \(\mathcal D\) 为 \(A\) 的 core。\foreignidx{core}
这里用明确的图像逼近条件描述这一子空间，不把它与零空间 \(\ker A\) 混同。

\begin{proposition}[限制算子的闭包判据]
\label{b1:prop:29-core-criterion}
设 \(H\) 为复 Hilbert 空间，\(A:D(A)\subseteq H\to H\) 为闭线性算子，
\(\mathcal D\subseteq D(A)\) 为线性子空间。
则限制算子 \(A|_{\mathcal D}\) 可闭，并且
\[
 \overline{A|_{\mathcal D}}=A
 \quad\Longleftrightarrow\quad
 \mathcal D\text{ 在 }D(A)\text{ 中关于图像范数稠密}.
\]
\end{proposition}
\begin{proof}
由于 \(G(A|_{\mathcal D})\subseteq G(A)\) 且后者闭，图像闭包仍包含于 \(G(A)\)，
故没有非零竖直向量，限制算子可闭。
其图像闭包等于 \(G(A)\)，当且仅当对每个 \(x\in D(A)\)，
存在 \(x_n\in\mathcal D\) 使 \((x_n,Ax_n)\to(x,Ax)\)。
乘积中的这一收敛恰好是 \(\|x_n-x\|_A\to0\)。
\end{proof}

\begin{proposition}[乘法算子的截断定义子空间]
\label{b1:prop:29-multiplication-core}
设 \((X,\Sigma,\mu)\) 为测度空间，\(m:X\to\mathbb R\) 为可测函数，
\(M_m\) 是 \(L^2(X,\mu;\mathbb C)\) 上定义域为 \(\{f:mf\in L^2\}\) 的乘法算子。
记 \(E_n=\{|m|\leq n\}\)，并令
\[
 \mathcal D_m=\bigcup_{n\geq1}
 \{f\in L^2(X,\mu;\mathbb C):f=0\text{ 几乎处处于 }X\setminus E_n\}.
\]
则 \(\mathcal D_m\subseteq D(M_m)\) 为线性子空间，在 \(D(M_m)\) 中关于图像范数稠密，
从而 \(\overline{M_m|_{\mathcal D_m}}=M_m\)。
\end{proposition}
\begin{proof}
各 \(E_n\) 递增，故上面的并为线性子空间；在 \(E_n\) 上 \(|m|\leq n\)，
所以每个所列函数都在定义域中。
对 \(f\in D(M_m)\)，取 \(f_n=\mathbf1_{E_n}\cdot f\in\mathcal D_m\)。有
\[
 \|f-f_n\|_2^2+\|M_mf-M_mf_n\|_2^2
 =\int_{X\setminus E_n}(1+|m|^2)\cdot|f|^2\,\mathrm d\mu
 \longrightarrow0.
\]
被积函数的可积性来自 \(f\in L^2\) 和 \(mf\in L^2\)，而 \(E_n\uparrow X\)。
控制收敛因此给出图像范数收敛，再用前一命题。
\end{proof}

这样的截断按照乘子的大小进行。对一般可测乘子，不能直接把 \(C_c^\infty\) 当成定义域的一部分。
例如在 \((0,1)\) 上令 \(m(x)=|x-1/2|^{-1}\)（在 \(x=1/2\) 处改为零）。
若 \(\varphi\in C_c^\infty(0,1)\) 且 \(\varphi(1/2)\ne0\)，
则 \(m\varphi\notin L^2\)。是否可以使用光滑函数作为图像逼近子空间，必须结合乘子的局部可积性另行判断。

\begin{theorem}[测试函数微分算子的闭包]
\label{b1:thm:29-minimal-derivative}
在复 Hilbert 空间 \(H=L^2((0,1),\mathrm dx;\mathbb C)\) 上，
令 \(P_cf=-\mathrm if'\)、\(D(P_c)=C_c^\infty(0,1)\)。
则 \(P_c\) 可闭，其闭包为
\[
 D(P_{\min})=H_0^1(0,1),\qquad P_{\min}f=-\mathrm if'.
\]
算子 \(P_{\min}\) 闭且对称，但不自伴；并有
\[
 \overline{P_c}=P_{\min},\qquad P_{\min}^*=P_{\max},
 \qquad D(P_{\max})=H^1(0,1).
\]
因此 \(P_c\) 不是本质自伴算子。对于每个 \(\theta\in\mathbb R\)，
\(C_c^\infty(0,1)\) 在 \(D(P_\theta)\) 中关于 \(P_\theta\) 的图像范数均不稠密。
\symidx{\(P_{\min}=\overline{P_c}\)：定义域为\(H_0^1(0,1)\)的一阶微分算子}
\end{theorem}
\begin{proof}
由 \(P_c\subseteq P_{\max}\) 及 \(P_{\max}\) 闭，\(P_c\) 可闭。
它的图像范数为 \(\|f\|_2+\|f'\|_2\)，与 \(H^1\) 范数等价。
\cref{b1:def:sobolev-zero-boundary}把 \(H_0^1\) 定义为 \(C_c^\infty\) 的 \(H^1\) 闭包，
因此每个 \(f\in H_0^1\) 都有一列测试函数使输入和导数同时收敛，
故 \((f,-\mathrm if')\in\overline{G(P_c)}\)。

反过来，若 \((f_n,-\mathrm if_n')\to(f,h)\)，其中 \(f_n\in C_c^\infty\)，
\(P_{\max}\) 的闭性给出 \(f\in H^1\)、\(h=-\mathrm if'\)。
于是 \(f_n\to f\) 于 \(H^1\)，故 \(f\in H_0^1\)。这证明闭包的公式，包括其整个定义域。

由闭包不改变伴随，及\cref{b1:prop:29-minimal-adjoint}，有
\(P_{\min}^*=P_c^*=P_{\max}\)。端点均为零使 \(P_{\min}\) 对称，
但常数函数 \(1\) 属于 \(H^1\) 而不属于 \(H_0^1\)，故两个定义域不相等，\(P_{\min}\) 不自伴。
因此 \(P_c\) 的闭包不是自伴算子。

最后，\(P_\theta\) 的图像范数仍等价于 \(H^1\) 范数，
而测试函数的闭包只有 \(H_0^1\)。函数 \(x\mapsto e^{\mathrm i\theta x}\)
在 \(D(P_\theta)\) 中，却不在 \(H_0^1\) 中，所以图像稠密性失败。
\end{proof}

还可以直接看到非实点满射为何失败。解一阶方程得到
\[
 \ker(P_{\max}-\mathrm iI)=\operatorname{span}\{e^{-x}\},\qquad
 \ker(P_{\max}+\mathrm iI)=\operatorname{span}\{e^x\}.
\]
由\eqref{b1:eq:29-range-kernel}，\(\operatorname{ran}(P_{\min}+\mathrm iI)\)
与 \(e^{-x}\) 正交，\(\operatorname{ran}(P_{\min}-\mathrm iI)\) 与 \(e^x\) 正交。
二者都不是稠密值域；闭性和对称性并没有排除这两个缺失方向。

\begin{table}[htbp]
\centering
\small
\caption[一阶微分算子的定义域比较]{同一表达式 \(-\mathrm i\,\mathrm d/\mathrm dx\) 的不同定义域。
所有算子都作用于 \(L^2(0,1;\mathbb C)\)，\(\theta\in\mathbb R\)。}
\label{b1:tab:29-derivative-domains}
\begin{tabular}{lllll}
\toprule
算子 & 定义域 & 闭 & 对称 & 自伴\\
\midrule
\(P_c\) & \(C_c^\infty(0,1)\) & 否 & 是 & 否\\
\(P_{\min}\) & \(H_0^1(0,1)\) & 是 & 是 & 否\\
\(P_\theta\) & \(\{f\in H^1:f(1)=e^{\mathrm i\theta}f(0)\}\) & 是 & 是 & 是\\
\(P_{\max}\) & \(H^1(0,1)\) & 是 & 否 & 否\\
\bottomrule
\end{tabular}
\end{table}

表中第一行的非闭性也可从闭包公式读出：\(x(1-x)\in H_0^1(0,1)\)，
却不是 \((0,1)\) 内的紧支集函数，故 \(P_c\ne\overline{P_c}\)。
从较小的计算定义域补成闭图像时，保留下来的端点条件由图像范数控制；
仅在 \(L^2\) 中的稠密性不能决定这些边界信息。

% !TeX root = ../../Book1.tex
% 实读来源：Rudin1991Functional，13.34--13.35，印刷 pp.375--379；
% 13.38，印刷 pp.382--385。大学公开扫描的零基 PDF 页序为印刷页码加 16。
% 目视核对 pp.379、382；正文自证酉群情形的平滑平均、两侧 Laplace 公式与唯一性，
% 不调用该书 13.38 所采用的 Gelfand 正规半群构造。
\section{强连续酉群与 Stone 定理}
\label{b1:sec:2908}

谱积分可以把实变量函数 \(e^{\mathrm it\lambda}\) 变成一个酉算子。
当时间 \(t\) 改变时，这些算子满足群律；反过来，
一个连续而且保持内积的演化，是否一定来自某个自伴算子？
Stone 定理给出肯定回答。这里的连续是逐个初值的强连续，
零时导数则只对一个稠密子空间中的初值存在。
因此，对应关系必须同时确定算子的作用和定义域。

\ProseParagraphBreak

本节固定
\[
 U(t)=e^{\mathrm itA},\qquad
 Gx=\lim_{h\to0}\frac{U(h)x-x}{h}=\mathrm iAx.
\]
其中 \(G\) 是由差商定义的生成元，\(A=-\mathrm iG\) 是要证明自伴的算子。
关于生成元的基本性质与 Laplace 公式，可对照 Rudin 的
《Functional Analysis》\cite[Theorem~13.35, pp.~376--379; Theorem~13.38, pp.~382--385]{Rudin1991Functional}。
下面在酉群的情形逐步完成证明。

\subsection{强连续酉群}
\label{b1:subsec:290801}

按\cref{b1:def:29-unitary-equivalence}，酉算子是满射线性等距映射。
在同一个 Hilbert 空间上，这等价于 \(U^*U=UU^*=I\)。
现在用时间参数把这些保持内积的算子组织起来。

\begin{definition}[强连续酉群]
\label{b1:def:29-unitary-group}
设 \(H\) 是复 Hilbert 空间。若一族有界线性算子
\((U(t))_{t\in\mathbb R}\subset\mathcal B(H)\) 满足
\[
 U(0)=I,\qquad U(t+s)=U(t)U(s)\quad(s,t\in\mathbb R),
\]
每个 \(U(t)\) 都是酉算子，并且对每个 \(x\in H\)，
\[
 \lim_{t\to0}\|U(t)x-x\|=0,
\]
则称 \((U(t))_{t\in\mathbb R}\) 为 \(H\) 上的
\term[qiang-lianxu-you-qun]{强连续酉群}[strongly continuous unitary group]。
\end{definition}

群律和酉性给出
\[
 U(t)^*=U(-t),\qquad \|U(t)x\|=\|x\|.
\]
而且
\[
 \|U(t+h)x-U(t)x\|=\|U(h)x-x\|,
\]
所以零时的强连续自动给出每条轨道 \(t\mapsto U(t)x\) 在全部实轴上的连续性。
这没有要求算子范数 \(\|U(h)-I\|\) 趋于零。

\ProseParagraphBreak

例如，在复空间 \(\ell^2\) 上令
\[
 U(t)(x_n)_{n\geq1}=(e^{\mathrm i n t}x_n)_{n\geq1}.
\]
群律和酉性可以逐坐标核对。对固定 \(x=(x_n)\in\ell^2\)，有
\[
 \|(U(h)-I)x\|_2^2
 \leq\sum_{n=1}^N|e^{\mathrm i n h}-1|^2\cdot|x_n|^2
       +4\sum_{n>N}|x_n|^2.
\]
先选 \(N\) 控制尾和，再令 \(h\to0\) 控制有限和，便得强连续。
但令 \(h_n=\pi/n\)，并以第 \(n\) 个单位向量测试，便有
\[
 \|U(h_n)-I\|\geq\|(U(h_n)-I)e_n\|_2=2.
\]
另一方面，酉性给出上界 \(2\)，故这些范数恰等于 \(2\)。
强连续允许高频坐标随时间尺度改变，不能把固定初值的极限改成单位球上的一致极限。

\ProseParagraphBreak

证明生成元的性质需要积分连续的 \(H\) 值轨道。
这里所需积分可以直接用 Riemann 和建立。
若 \(v:[a,b]\to H\) 连续，则它一致连续；
把两个足够细的带标记分割细化到共同分割后，
相应 Riemann 和之差的范数由各小区间长度乘以 \(v\) 的连续模控制。
这些和因而构成 Cauchy 族，\(H\) 的完备性给出唯一极限，记为
\(\int_a^b v(s)\,\mathrm ds\)。从有限和取极限可得
\begin{equation}\label{b1:eq:29-vector-integral-bound}
 \left\|\int_a^b v(s)\,\mathrm ds\right\|
 \leq\int_a^b\|v(s)\|\,\mathrm ds.
\end{equation}
有界线性算子可与这个积分交换，内积对任一固定向量的测试也可与积分交换。
约定 \(\int_b^a=-\int_a^b\)，则对连续 \(v\)，
\[
 \left\|\frac1h\int_t^{t+h}v(s)\,\mathrm ds-v(t)\right\|
 \leq\sup_{|s-t|\leq|h|}\|v(s)-v(t)\|\longrightarrow0.
\]
这就是所需的强微积分基本公式。
若 \(F:[a,b]\to H\) 连续可微，对任一 \(y\in H\) 将标量微积分基本定理
用于 \(\langle F(s),y\rangle\)，便得
\(F(b)-F(a)=\int_a^bF'(s)\,\mathrm ds\)。
以下有限区间积分均取这个意义，无穷区间积分则取有限区间积分的范数极限。

\subsection{生成元、符号与定义域}
\label{b1:subsec:290802}

\begin{definition}[酉群的生成元]
\label{b1:def:29-group-generator}
设 \(H\) 是复 Hilbert 空间，\((U(t))_{t\in\mathbb R}\) 是其上的强连续酉群。
定义
\begin{equation}\label{b1:eq:29-generator-domain}
 \begin{split}
 D(G)&=\left\{x\in H:
       \lim_{h\to0}\frac{U(h)x-x}{h}\text{ 在 }H\text{ 中存在}\right\},\\
 Gx&=\lim_{h\to0}\frac{U(h)x-x}{h}\qquad(x\in D(G)).
 \end{split}
\end{equation}
线性算子 \(G:D(G)\subset H\to H\) 称为该群的
\term[shengchengyuan]{生成元}[generator]。
极限取实参数 \(h\) 的双侧极限，收敛均指 \(H\) 的范数收敛。
\end{definition}

差商的线性表明 \(D(G)\) 是线性子空间，\(G\) 在其中线性。
对群而言，只要求右侧差商收敛也会得到同一个定义域：若右极限为 \(y\)，
则对 \(h>0\)，
\[
 \frac{U(-h)x-x}{-h}
 =U(-h)\frac{U(h)x-x}{h}\longrightarrow y.
\]
最后一步可把右端与 \(y\) 的差拆成
\(U(-h)((U(h)x-x)/h-y)+(U(-h)y-y)\)，
再用酉性和强连续性分别控制两项。

\begin{proposition}[定义域不变与轨道积分式]
\label{b1:prop:29-group-orbits}
设 \(H\) 是复 Hilbert 空间，\((U(t))_{t\in\mathbb R}\) 是其上的强连续酉群，
\(G\) 为其生成元。对每个 \(t\in\mathbb R\)，有
\[
 U(t)D(G)=D(G),\qquad GU(t)x=U(t)Gx\quad(x\in D(G)).
\]
对每个 \(x\in D(G)\)，轨道是 \(\mathbb R\) 上连续可微的 \(H\) 值函数，且
\begin{equation}\label{b1:eq:29-group-orbit-derivative}
 \frac{\mathrm d}{\mathrm dt}U(t)x=U(t)Gx=GU(t)x,
 \qquad
 U(t)x-x=\int_0^tU(s)Gx\,\mathrm ds.
\end{equation}
\end{proposition}

\begin{proof}
若 \(x\in D(G)\)，群律给出
\[
 \frac{U(h)U(t)x-U(t)x}{h}
 =U(t)\frac{U(h)x-x}{h}\longrightarrow U(t)Gx.
\]
故 \(U(t)x\in D(G)\)，并得到交换公式。
再把时间换成 \(-t\)，可得定义域之间反向的包含，从而是等式。
同一差商也给出 \(t\) 时刻的轨道导数。
由于 \(s\mapsto U(s)Gx\) 强连续，导数连续，
上一小节的微积分基本公式便给出积分恒等式，包括负时间的情形。
\end{proof}

\begin{proposition}[稠密性、闭性与内积关系]
\label{b1:prop:29-generator-closed}
设 \(H\) 是复 Hilbert 空间，\((U(t))_{t\in\mathbb R}\) 是其上的强连续酉群，
\(G:D(G)\subset H\to H\) 为其生成元。
则 \(D(G)\) 在 \(H\) 中稠密，\(G\) 是闭算子，并且
\begin{equation}\label{b1:eq:29-generator-skew-symmetry}
 \langle Gx,y\rangle+\langle x,Gy\rangle=0
 \qquad(x,y\in D(G)).
\end{equation}
因此 \(A=-\mathrm iG\)、\(D(A)=D(G)\) 是稠密定义的闭对称算子。
\end{proposition}

\begin{proof}
先以平滑的时间平均逼近任意 \(x\in H\)。
取非负 \(\varphi\in C_c^1(\mathbb R)\)，支集包含于 \([-1,1]\)，且
\(\int_{\mathbb R}\varphi(s)\,\mathrm ds=1\)；
例如把 \((1-s^2)^2\) 在 \([-1,1]\) 外延为零，再作归一化。
对 \(\varepsilon>0\)，令
\[
 \varphi_\varepsilon(s)=\varepsilon^{-1}\varphi(s/\varepsilon),\qquad
 x_\varepsilon=\int_{\mathbb R}\varphi_\varepsilon(s)U(s)x\,\mathrm ds.
\]
这些积分的被积函数连续且具有紧支集。由积分估计，
\[
 \|x_\varepsilon-x\|
 \leq\int\varphi_\varepsilon(s)\|U(s)x-x\|\,\mathrm ds
 \leq\sup_{|s|\leq\varepsilon}\|U(s)x-x\|\longrightarrow0.
\]
还需证明每个 \(x_\varepsilon\) 确实属于生成元域。
将有界算子 \(U(h)\) 移入积分，再以 \(s\) 代替原变量加 \(h\)，得到
\[
 \frac{U(h)x_\varepsilon-x_\varepsilon}{h}
 =\int_{\mathbb R}
   \frac{\varphi_\varepsilon(s-h)-\varphi_\varepsilon(s)}{h}
   U(s)x\,\mathrm ds.
\]
标量差商等于
\(-\int_0^1\varphi_\varepsilon'(s-\theta h)\,\mathrm d\theta\)。
由于 \(\varphi_\varepsilon'\) 一致连续，它在 \(h\to0\) 时一致趋于
\(-\varphi_\varepsilon'(s)\)；当 \(|h|<1\) 时，所有被积函数又有共同的紧支集。
用 \(\|U(s)x\|=\|x\|\) 及积分估计，可将极限移入积分，得到
\[
 x_\varepsilon\in D(G),\qquad
 Gx_\varepsilon=-\int_{\mathbb R}\varphi_\varepsilon'(s)U(s)x\,\mathrm ds.
\]
故 \(D(G)\) 稠密。

为证闭性，设 \(x_n\in D(G)\)、\(x_n\to x\) 且 \(Gx_n\to y\)。
由\eqref{b1:eq:29-group-orbit-derivative}，
\[
 U(t)x_n-x_n=\int_0^tU(s)Gx_n\,\mathrm ds.
\]
左端趋于 \(U(t)x-x\)，右端与 \(\int_0^tU(s)y\,\mathrm ds\) 的差的范数
不超过 \(|t|\cdot\|Gx_n-y\|\)。因此对每个 \(t\in\mathbb R\)，
\[
 U(t)x-x=\int_0^tU(s)y\,\mathrm ds.
\]
除以非零 \(t\) 后令 \(t\to0\)，积分平均由强连续性趋于 \(y\)。
于是 \(x\in D(G)\)、\(Gx=y\)，这就证明了闭性。

最后，对 \(x,y\in D(G)\)，酉性给出
\(\langle U(t)x,U(t)y\rangle=\langle x,y\rangle\)。
在零时求导，便得\eqref{b1:eq:29-generator-skew-symmetry}。
本书内积对第一变量线性，所以
\[
 \langle-\mathrm iGx,y\rangle
 =-\mathrm i\langle Gx,y\rangle
 =\mathrm i\langle x,Gy\rangle
 =\langle x,-\mathrm iGy\rangle.
\]
这证明 \(A=-\mathrm iG\) 对称；其稠密性和闭性由 \(G\) 的相应性质给出。
\end{proof}

对称性尚未说明 \(A\) 自伴。
这里仍缺少伴随定义域的辨认，最后一小节将借助两个非实参数处的满射性补齐它。

\subsection{从自伴算子到酉群}
\label{b1:subsec:290803}

先从自伴算子构造演化。在这个方向，前面建立的谱演算同时给出
算子、连续性以及准确的可微初值范围。

\begin{proposition}[自伴算子的指数群]
\label{b1:prop:29-spectral-unitary-group}
设 \(H\) 是复 Hilbert 空间，\(A:D(A)\subset H\to H\) 是自伴算子，
\(E\) 是它在 \(\mathbb R\) 上的投影值谱测度。对 \(t\in\mathbb R\)，定义
\begin{equation}\label{b1:eq:29-unitary-spectral-integral}
 U_A(t)=e^{\mathrm itA}
       =\int_{\mathbb R}e^{\mathrm it\lambda}\,\mathrm dE(\lambda).
\end{equation}
则 \((U_A(t))_{t\in\mathbb R}\) 是 \(H\) 上的强连续酉群，
它的生成元恰为
\[
 D(G_A)=D(A),\qquad G_Ax=\mathrm iAx\quad(x\in D(A)).
\]
对每个 \(t\in\mathbb R\)，有
\[
 U_A(t)D(A)=D(A),\qquad AU_A(t)x=U_A(t)Ax\quad(x\in D(A)),
\]
且对每个 \(x\in D(A)\)，
\[
 \frac{\mathrm d}{\mathrm dt}U_A(t)x
 =\mathrm iAU_A(t)x=\mathrm iU_A(t)Ax
\]
是连续的强导数。
\end{proposition}

\begin{proof}
由无界自伴谱定理\cref{b1:thm:29-unbounded-spectral}及相应函数演算，
\(|e^{\mathrm it\lambda}|=1\) 保证 \(U_A(t)\) 是定义在全部 \(H\) 上的有界算子。
有界 Borel 演算的乘法与伴随规则给出
\[
 U_A(0)=I,\quad U_A(t+s)=U_A(t)U_A(s),\quad
 U_A(t)^*=U_A(-t)=U_A(t)^{-1}.
\]
对 \(x\in H\)，记 \(\mu_x(B)=\langle E(B)x,x\rangle\)。
谱范数公式给出
\[
 \|U_A(h)x-x\|^2
 =\int_{\mathbb R}|e^{\mathrm ih\lambda}-1|^2\,\mathrm d\mu_x(\lambda).
\]
被积函数逐点趋于零，且不超过 \(4\)，而 \(\mu_x(\mathbb R)=\|x\|^2\)。
支配收敛定理于是给出零时强连续，群律再给出全部时间的强连续。

若 \(x\in D(A)\)，则 \(\int\lambda^2\,\mathrm d\mu_x<\infty\)。
对实数 \(\lambda\) 及非零实数 \(h\)，标量恒等式
\[
 \frac{e^{\mathrm ih\lambda}-1}{h}
 =\mathrm i\lambda\int_0^1e^{\mathrm i\theta h\lambda}\,\mathrm d\theta
\]
给出其模不超过 \(|\lambda|\)，而点态极限为 \(\mathrm i\lambda\)。
因而
\[
 \begin{split}
 &\left\|\frac{U_A(h)x-x}{h}-\mathrm iAx\right\|^2\\
 &\qquad=\int_{\mathbb R}
  \left|\frac{e^{\mathrm ih\lambda}-1}{h}-\mathrm i\lambda\right|^2
  \,\mathrm d\mu_x(\lambda)\longrightarrow0,
 \end{split}
\]
这里误差平方由可积函数 \(4\lambda^2\) 控制。
这证明 \(D(A)\subseteq D(G_A)\)，以及在该域上的作用公式。

反过来，设 \(x\in D(G_A)\)。差商收敛使其范数在充分小的时间内有界。
取 \(h_n=1/n\)，由 Fatou 引理及谱范数公式，
\[
 \begin{split}
 \int_{\mathbb R}\lambda^2\,\mathrm d\mu_x(\lambda)
 &\leq\liminf_{n\to\infty}
  \int_{\mathbb R}
   \left|\frac{e^{\mathrm ih_n\lambda}-1}{h_n}\right|^2
   \,\mathrm d\mu_x(\lambda)\\
 &=\liminf_{n\to\infty}
  \left\|\frac{U_A(h_n)x-x}{h_n}\right\|^2<\infty.
 \end{split}
\]
故 \(x\in D(A)\)。于是定义域相等，而不只是在预先给定的域上导数相同。
上述必要性论证甚至只需要右侧差商存在。

还可直接从谱测度看出定义域为什么不变。
投影 \(E(B)\) 与 \(U_A(t)\) 交换，酉性给出
\[
 \mu_{U_A(t)x}(B)
 =\|E(B)U_A(t)x\|^2
 =\|U_A(t)E(B)x\|^2=\mu_x(B).
\]
因此 \(\int\lambda^2\,\mathrm d\mu_{U_A(t)x}
=\int\lambda^2\,\mathrm d\mu_x\)，这给出双向的域不变性。
若 \(x\in D(A)\)，令
\(A_N=\int\lambda\cdot\mathbf1_{[-N,N]}(\lambda)\,\mathrm dE(\lambda)\)。
有界乘法规则给出 \(A_NU_A(t)x=U_A(t)A_Nx\)；
两端分别趋于 \(AU_A(t)x\) 和 \(U_A(t)Ax\)，故无界算子的交换公式成立。
最后应用\cref{b1:prop:29-group-orbits}，便得到全部时间的强导数。
\end{proof}

把实乘法算子定理\cref{b1:thm:29-real-multiplication}用于 \(\mathbb N\) 上的计数测度，
开头的对角群便对应于自伴算子 \(Ax=(nx_n)\)，其最大定义域为
\[
 D(A)=\left\{x\in\ell^2:\sum_{n=1}^{\infty}n^2|x_n|^2<\infty\right\},
 \qquad Gx=(\mathrm i n x_n).
\]
这也可以由逐坐标差商与平方和的 Fatou 引理直接核对。
例如 \(x=(1/n)_{n\geq1}\) 属于 \(\ell^2\)，却不属于 \(D(A)\)。
其轨道处处连续，但在任何时间都没有强导数：
若某时刻差商收敛，左乘该时刻演化的逆，就会得到零时差商收敛，
与定义域的判别矛盾。

\subsection{Stone 对应与唯一性}
\label{b1:subsec:290804}

给定一个强连续酉群，\cref{b1:prop:29-generator-closed}已证明 \(-\mathrm iG\) 闭且对称。
要应用自伴满射判据，需要证明两个相反方向的非实参数都没有可解性缺口。
正、负时间轨道的 Laplace 积分恰好提供这两个逆。

\begin{lemma}[两侧 Laplace 公式]
\label{b1:lem:29-group-laplace}
设 \(H\) 是复 Hilbert 空间，\((U(t))_{t\in\mathbb R}\) 是其上的强连续酉群，
\(G\) 是其生成元。对 \(\lambda>0\)，定义
\[
 \begin{split}
 B_\lambda^+x&=\int_0^\infty e^{-\lambda s}U(s)x\,\mathrm ds,\\
 B_\lambda^-x&=\int_{-\infty}^0e^{\lambda s}U(s)x\,\mathrm ds
              =\int_0^\infty e^{-\lambda s}U(-s)x\,\mathrm ds
              \qquad(x\in H).
 \end{split}
\]
则 \(B_\lambda^\pm\in\mathcal B(H)\)、\(\|B_\lambda^\pm\|\leq\lambda^{-1}\)，且
\[
 (\lambda I-G)^{-1}=B_\lambda^+,\qquad
 (\lambda I+G)^{-1}=B_\lambda^-.
\]
这里 \(\lambda I\mp G:D(G)\to H\) 均为双射，
各逆的值域恰为 \(D(G)\)。
\end{lemma}

\begin{proof}
先逐个固定 \(x\in H\) 作积分。
酉性和\eqref{b1:eq:29-vector-integral-bound}给出
\[
 \left\|\int_T^R e^{-\lambda s}U(\pm s)x\,\mathrm ds\right\|
 \leq\|x\|\int_T^R e^{-\lambda s}\,\mathrm ds
 \leq\frac{e^{-\lambda T}}{\lambda}\|x\|
 \qquad(0\leq T<R).
\]
因此无穷积分按范数收敛；其线性及范数估计说明它们定义了有界算子。
这个构造没有把 \(s\mapsto U(s)\) 当作按算子范数连续的函数来积分。

令 \(b=B_\lambda^+x\)。将 \(U(h)\) 移入有限区间积分、换元，再取无穷积分极限，得到
\[
 U(h)b=e^{\lambda h}
       \left(b-\int_0^h e^{-\lambda r}U(r)x\,\mathrm dr\right).
\]
此式对正、负 \(h\) 均成立，负上限使用有向积分的约定。
于是
\[
 \frac{U(h)b-b}{h}
 =\frac{e^{\lambda h}-1}{h}b
  -\frac{e^{\lambda h}}{h}
   \int_0^h e^{-\lambda r}U(r)x\,\mathrm dr
 \longrightarrow\lambda b-x.
\]
故 \(b\in D(G)\)，且 \((\lambda I-G)B_\lambda^+x=x\)。
反向时间族 \(V(t)=U(-t)\) 仍为强连续酉群，
由双侧差商定义，它的生成元是 \(-G\)，定义域仍为 \(D(G)\)。
把刚才的计算用于 \(V\)，便得
\((\lambda I+G)B_\lambda^-x=x\)。这已经证明两个算子满射。

为证单射，对 \(y\in D(G)\)，由\eqref{b1:eq:29-generator-skew-symmetry}可知
\(\operatorname{Re}\langle Gy,y\rangle=0\)，从而
\[
 \|(\lambda I\mp G)y\|^2
 =\lambda^2\|y\|^2+\|Gy\|^2
 \geq\lambda^2\|y\|^2.
\]
所以两者均为单射，已经构造的右逆就是双侧逆。
具体地，对 \(y\in D(G)\)，将
\(B_\lambda^\pm(\lambda I\mp G)y-y\) 代入 \(\lambda I\mp G\)，结果为零，
单射性便迫使原向量为零。这也证明逆的值域恰为整个 \(D(G)\)。
\end{proof}

下面的\term[stone-dingli]{Stone 定理}[Stone's theorem]把自伴算子与强连续酉群一一对应起来。

\begin{theorem}[Stone]
\label{b1:thm:29-stone}
设 \(H\) 是复 Hilbert 空间。对 \(H\) 上的每个强连续酉群
\((U(t))_{t\in\mathbb R}\)，存在唯一自伴算子
\(A:D(A)\subset H\to H\)，使
\[
 U(t)=e^{\mathrm itA}\qquad(t\in\mathbb R),
\]
其中右端按谱演算定义。若 \(G\) 是该群的生成元，则准确地有
\[
 D(A)=D(G),\qquad A=-\mathrm iG,\qquad G^*=-G.
\]
反过来，每个自伴算子通过 \(e^{\mathrm itA}\) 给出一个强连续酉群，
其生成元为 \(\mathrm iA\)，定义域为 \(D(A)\)。
具有相同生成元的两个强连续酉群必然相同；这里相同包括定义域相同。
\end{theorem}

\begin{proof}
给定强连续酉群，令 \(A=-\mathrm iG\)、\(D(A)=D(G)\)。
\cref{b1:prop:29-generator-closed}表明它稠密定义、闭且对称。
由\cref{b1:lem:29-group-laplace}在 \(\lambda=1\) 时的满射性，
\[
 A+\mathrm iI=\mathrm i(I-G),\qquad
 A-\mathrm iI=-\mathrm i(I+G)
\]
都有值域 \(H\)。自伴满射判据\cref{b1:thm:29-self-adjoint-criterion}
于是给出 \(A=A^*\)。由于 \(G=\mathrm iA\)，伴随定义也给出
\(D(G^*)=D(A)\)、\(G^*=-\mathrm iA=-G\)。

现在由\cref{b1:prop:29-spectral-unitary-group}构造
\(V(t)=e^{\mathrm itA}\)。这个群的生成元是 \(\mathrm iA=G\)，
而且生成元定义域恰为 \(D(G)\)。
还须证明相同的零时导数确实决定整个群。
事实上，对任意具有同一生成元 \(G\) 的两个强连续酉群 \(U,V\)，
固定 \(t\in\mathbb R\)、\(x\in D(G)\)，令
\[
 F(s)=U(t-s)V(s)x\qquad(s\in\mathbb R).
\]
由定义域不变性，\(V(s)x\in D(G)\)。为核对对移动向量求导的步骤，写出
\[
 \begin{split}
 \frac{F(s+h)-F(s)}{h}
 &=U(t-s-h)\frac{V(s+h)x-V(s)x}{h}\\
 &\quad+\frac{U(t-s-h)-U(t-s)}{h}V(s)x.
 \end{split}
\]
第一项趋于 \(U(t-s)GV(s)x\)：
先用酉性控制差商与其极限的差，再对固定极限向量使用强连续性。
第二项的输入 \(V(s)x\) 已固定且属于 \(D(G)\)，
轨道导数公式给出其极限为 \(-U(t-s)GV(s)x\)。
因此 \(F'(s)=0\)。对任意固定 \(y\in H\)，
标量函数 \(\langle F(s),y\rangle\) 导数为零，故 \(F\) 是常值。
比较 \(s=0\) 与 \(s=t\) 得 \(U(t)x=V(t)x\)。
由于 \(D(G)\) 稠密，且两个算子有界，这个等式延伸到所有 \(x\in H\)。
应用于已构造的 \(V\)，就得到 \(U(t)=e^{\mathrm itA}\)。

若另一个自伴算子 \(B\) 也满足 \(U(t)=e^{\mathrm itB}\)，
\cref{b1:prop:29-spectral-unitary-group}的定义域结论给出
\(D(B)=D(G)\)、\(\mathrm iBx=Gx\)。故 \(B=A\)，包括定义域和作用都相同。
自伴算子到酉群的反方向已经由该命题完整证明。
\end{proof}

这个对应把方程
\[
 u'(t)=\mathrm iAu(t),\qquad u(0)=x\in D(A)
\]
的演化写成 \(u(t)=U(t)x\)，其范数保持为 \(\|x\|\)。
对任意 \(x\in H\)，连续轨道仍有定义；它具有强导数当且仅当 \(x\in D(A)\)。
由于 \(U(t)D(A)=D(A)\)，域外的初值在任何固定时间都不会进入该域。
若把指数约定改为 \(e^{-\mathrm itA}\)，相应的差商生成元便改为
\(-\mathrm iA\)；符号与定义域都属于对应关系的一部分。


\begin{chaptersummary}
预解算子把谱与方程的可解性联系起来。对有界自伴算子，谱是非空实紧集，
算子范数等于谱上绝对值的最大值；谱点也可以由近似特征向量识别，
因而没有特征向量并不意味着没有谱。

连续函数演算由多项式逼近建立，正测度表示把每个循环子空间化为乘法模型。
各循环部分的正交直和给出整体谱表示，而投影值测度把它写成与分解选择无关的公式
\[
 T=\int t\,\mathrm dE(t),\qquad f(T)=\int f(t)\,\mathrm dE(t).
\]
谱投影同时描述特征空间、谱区间和有界 Borel 函数的强极限。
逐点收敛经支配收敛传到每个向量上，一般不能传成算子范数收敛。

无界算子必须连同定义域讨论。稠密性使伴随良定义，图像闭性控制极限，
非实参数处的满射性区分对称算子与自伴算子。
Cayley 变换把无界自伴算子转为酉算子，谱积分再恢复原来的最大定义域：
\[
 D(f(A))=\left\{x:\int|f|^2\,\mathrm d\mu_x<\infty\right\}.
\]
这项平方可积条件决定了乘积、和与闭包的关系。
实乘法算子及区间上的微分算子表明，同一作用法则配上不同定义域，
可以得到不同的闭性、伴随和边界条件。

Stone 定理把这些工具用于演化：每个强连续酉群唯一写成 \(U(t)=e^{\mathrm itA}\)，
其中 \(A\) 自伴，且差商生成元为 \(G=\mathrm iA\)，二者定义域相同。
谱积分构造演化，轨道平均与两侧 Laplace 积分完成逆向恢复。
这使连续谱、无界算子与保持范数的时间演化进入同一个表示框架。
\end{chaptersummary}

\BookExercises

% 自拟：二维非自伴模型，区分算子范数、谱界与预解算子增长。
\begin{optionalexercise}\label{b1:ex:29-nilpotent-resolvent}
在带标准内积的 \(\mathbb C^2\) 上，固定 \(a>0\)，令
\[
 N_a=\begin{pmatrix}0&a\\0&0\end{pmatrix}.
\]
求 \(\sigma(N_a)\) 和 \(\|N_a\|\)，并证明对 \(z\ne0\)，
\[
 R_{N_a}(z)=-z^{-1}I-z^{-2}N_a,\qquad
 \|R_{N_a}(z)\|\geq\frac a{|z|^2}.
\]
当 \(a>|z|\) 时，将这个下界与 \(\operatorname{dist}(z,\sigma(N_a))^{-1}\) 比较。
\end{optionalexercise}
\begin{solution}
由 \(N_a(u,v)=(av,0)\)，有 \(\|N_a\|=a\)，且零是特征值。
又因 \(N_a^2=0\)，对非零 \(z\) 直接相乘可得
\[
 (N_a-zI)(-z^{-1}I-z^{-2}N_a)
 =(-z^{-1}I-z^{-2}N_a)(N_a-zI)=I.
\]
所以 \(\sigma(N_a)=\{0\}\)，且题给公式是预解算子。
令 \(e_2=(0,1)\)，则
\[
 R_{N_a}(z)e_2=(-az^{-2},-z^{-1}),\qquad
 \|R_{N_a}(z)\|\geq\bigl(a^2|z|^{-4}+|z|^{-2}\bigr)^{1/2}
 \geq a|z|^{-2}.
\]
当 \(a>|z|\) 时，这个下界大于 \(|z|^{-1}
=\operatorname{dist}(z,\sigma(N_a))^{-1}\)。
本例的谱中最大模为零，而算子范数为 \(a\)；
因此\cref{b1:thm:29-self-adjoint-spectral-bound}中的自伴性不能去掉。
\end{solution}

% 自拟：稠密对角系数模型，训练点谱与离散谱的区别及非紧性判断。
\begin{optionalexercise}\label{b1:ex:29-dense-diagonal-spectrum}
设 \((r_n)_{n\geq1}\) 是 \(\mathbb Q\cap(0,1)\) 的无重复枚举，
在复 Hilbert 空间 \(\ell^2\) 上定义 \(Ae_n=r_ne_n\)。证明 \(A\) 有界且自伴，并求其谱、点谱和连续谱。
证明 \(A\) 没有离散谱点，且不是紧算子。
\end{optionalexercise}
\begin{solution}
对 \(x=(x_n)\in\ell^2\)，有 \(Ax=(r_nx_n)\) 和 \(\|Ax\|_2\leq\|x\|_2\)。
每个 \(r_n\) 都是实数，故
\(\langle Ax,y\rangle=\sum_n r_nx_n\overline{y_n}=\langle x,Ay\rangle\)，
于是 \(A\) 自伴。若 \(z\notin[0,1]\)，
则 \(y\mapsto(y_n/(r_n-z))\) 给出有界双侧逆。
若 \(t\in[0,1]\)，由有理数的稠密性可以取指标列 \((n_k)\)，使 \(r_{n_k}\to t\)，从而
\[
 \|(A-tI)e_{n_k}\|_2=|r_{n_k}-t|\longrightarrow0.
\]
\cref{b1:cor:29-approximate-eigenvectors}给出 \(\sigma(A)=[0,1]\)。
坐标方程 \((r_n-t)x_n=0\) 表明特征值恰为这些有理数；
当 \(t\in[0,1]\setminus\mathbb Q\) 或 \(t=0,1\) 时，算子 \(A-tI\) 单射。
它的值域由自伴性知为稠密，又因 \(t\) 是谱点而不能满射。因此
\[
 \sigma_{\mathrm p}(A)=\mathbb Q\cap(0,1),\qquad
 \sigma_{\mathrm c}(A)=[0,1]\setminus\bigl(\mathbb Q\cap(0,1)\bigr).
\]
整个区间没有孤立点，所以离散谱为空，即使点谱在区间内稠密。

区间 \([1/2,1)\) 中有无穷多个 \(r_n\)。对相应的不同指标 \(n,m\)，
\[
 \|Ae_n-Ae_m\|_2=(r_n^2+r_m^2)^{1/2}\geq1/\sqrt2.
\]
这批单位向量的像没有 Cauchy 子列，因而没有收敛子列，故 \(A\) 不紧。
\end{solution}

% 自拟：把等距连续函数演算用于最佳预解估计，与第一题的非自伴增长比较。
\begin{optionalexercise}\label{b1:ex:29-sharp-resolvent}
设 \(H\ne\{0\}\) 为复 Hilbert 空间，\(T\in\mathcal B(H)\) 自伴。
证明对每个 \(z\in\rho(T)\)，
\[
 \|R_T(z)\|=\frac1{\operatorname{dist}(z,\sigma(T))}.
\]
若 \((T-zI)x=b\)、\((T-zI)\widetilde x=\widetilde b\)，
由此给出 \(\|x-\widetilde x\|\) 的最佳统一估计；
说明这里的“最佳”不要求存在达到等号的单位向量。
\end{optionalexercise}
\begin{solution}
在紧集 \(K=\sigma(T)\) 上，函数 \(r(t)=(t-z)^{-1}\) 连续。
连续函数演算的乘法性给出
\[
 (T-zI)r(T)=r(T)(T-zI)=I.
\]
所以 \(r(T)=R_T(z)\)。等距性又给出
\[
 \|R_T(z)\|=\max_{t\in K}|t-z|^{-1}
 =\operatorname{dist}(z,K)^{-1}.
\]
两方程相减，得到
\[
 \|x-\widetilde x\|
 \leq\frac{\|b-\widetilde b\|}{\operatorname{dist}(z,K)}.
\]
若一个常数 \(C\) 对所有右端差都满足这项估计，取
\(\widetilde b=0\)、\(b=y\)，便有 \(\|R_T(z)y\|\leq C\|y\|\)，
故 \(C\geq\|R_T(z)\|\)。算子范数是单位球上取值的上确界；
可选单位向量使这些取值任意接近该数，而不必在某个向量上取到它。
\end{solution}

% 自拟：重复的乘法表示，训练循环向量不存在时的正交向量构造。
\begin{optionalexercise}\label{b1:ex:29-cyclic-multiplicity}
在复 Hilbert 空间 \(H=L^2(0,1)\oplus L^2(0,1)\) 上定义
\[
 T(f,g)=(tf,tg).
\]
证明每个单独分量上的乘法算子 \(M_t\) 都有循环向量，而 \(T\) 没有循环向量。
\end{optionalexercise}
\begin{solution}
常函数 \(1\) 属于 \(L^2(0,1)\)。多项式在 \(C([0,1])\) 中一致稠密，
连续函数又在 \(L^2(0,1)\) 中稠密，故 \(\{p(M_t)1:p\text{ 为多项式}\}\)
的线性闭包为整个 \(L^2(0,1)\)。所以 \(1\) 是每个分量的循环向量。

任取非零 \(v=(f,g)\in H\)，令 \(w=(\overline g,-\overline f)\)。
则 \(w\in H\)、\(\|w\|=\|v\|>0\)，且对每个复多项式 \(p\)，
\[
 \langle p(T)v,w\rangle
 =\int_0^1 p(t)f(t)g(t)\,\mathrm dt
   -\int_0^1 p(t)g(t)f(t)\,\mathrm dt=0.
\]
这些积分存在，因为 \(p\) 在 \([0,1]\) 上有界，\(fg\in L^1(0,1)\)。
非零向量 \(w\) 与 \(v\) 生成的循环子空间正交，故该子空间不等于 \(H\)。
零向量也不能循环，因此 \(T\) 没有循环向量。
\end{solution}

% 自拟：用连续函数逼近单点谱投影，辨析强收敛与算子范数收敛。
\begin{optionalexercise}\label{b1:ex:29-atom-approximation}
设 \(H\ne\{0\}\) 为复 Hilbert 空间，\(T\in\mathcal B(H)\) 自伴，\(E\) 为其谱测度。
固定 \(\lambda\in\mathbb R\)，证明
\[
 (I+n(T-\lambda I)^2)^{-1}x\longrightarrow E(\{\lambda\})x
 \qquad(x\in H).
\]
这里把 \(E\) 在谱以外延拓为零。对 \(L^2(0,1)\) 上的 \(T=M_t\)
及任意 \(\lambda\in[0,1]\)，证明左侧算子强收敛于零，但其范数恒为 \(1\)。
\end{optionalexercise}
\begin{solution}
令 \(f_n(t)=(1+n(t-\lambda)^2)^{-1}\)。在 \(\sigma(T)\) 上，
\(f_n\) 连续且满足 \(0<f_n\leq1\)。连续函数演算的乘法性说明
\(f_n(T)=(I+n(T-\lambda I)^2)^{-1}\)。逐点地，
\(f_n\to\mathbf1_{\{\lambda\}}\)，故有界 Borel 演算的支配收敛给出所需强极限。

对 \(M_t\)，相应谱投影就是乘以 \(\mathbf1_{\{\lambda\}}\)，
该函数在 Lebesgue 几乎处处为零，所以极限算子为零。
另一方面，\(\sigma(M_t)=[0,1]\)，由连续函数演算的等距性，
\[
 \|f_n(M_t)\|=\max_{t\in[0,1]}\frac1{1+n(t-\lambda)^2}=1.
\]
端点 \(\lambda=0,1\) 也满足这个等式；附近小区间仍有正测度，
因此端点处的最大值并未改变相应乘子的本质上确界。
\end{solution}

% 自拟：用两侧非实满射判据与Neumann级数证明有界自伴扰动的稳定性。
\begin{optionalexercise}\label{b1:ex:29-bounded-perturbation}
设 \(H\) 为复 Hilbert 空间，\(A:D(A)\subseteq H\to H\) 为稠密定义的自伴算子，
\(K\in\mathcal B(H)\) 自伴。证明 \(B=A+K\) 在定义域 \(D(B)=D(A)\) 上自伴。
\end{optionalexercise}
\begin{solution}
\(B\) 稠密定义，并且对 \(x,y\in D(A)\)，
\[
 \langle Bx,y\rangle
 =\langle Ax,y\rangle+\langle Kx,y\rangle
 =\langle x,Ay\rangle+\langle x,Ky\rangle
 =\langle x,By\rangle.
\]
所以 \(B\) 对称。取实数 \(r>\|K\|\)，分别令 \(z=\mathrm ir\) 和 \(z=-\mathrm ir\)。
无界自伴算子的预解估计给出
\[
 \|K(A-zI)^{-1}\|\leq\|K\|\cdot\|(A-zI)^{-1}\|
 \leq\frac{\|K\|}{r}<1.
\]
于是 \(I+K(A-zI)^{-1}\) 由 Neumann 级数可逆。对 \(x\in D(A)\)，
\[
 (B-zI)x=\bigl(I+K(A-zI)^{-1}\bigr)(A-zI)x.
\]
右侧两个因子分别把 \(D(A)\) 双射到 \(H\)、再把 \(H\) 双射到自身，
故 \(B-\mathrm irI\) 与 \(B+\mathrm irI\) 均满射。
由\cref{b1:thm:29-self-adjoint-criterion}的共轭非实参数满射判据，\(B\) 自伴。
这个论证同时验证了原定义域上的对称性与共轭非实参数处的满射性，因而控制了算子相加后的自伴性。
\end{solution}

% 自拟：由演化的周期性识别离散的谱支集，并写出无界生成算子的精确域。
\begin{optionalexercise}\label{b1:ex:29-periodic-group}
设 \(H\) 为复 Hilbert 空间，\(A:D(A)\subseteq H\to H\) 为稠密定义的自伴算子，
\(E\) 为其谱测度，\(U(t)=e^{\mathrm itA}\)。固定 \(\tau>0\)，证明
\[
 U(\tau)=I
 \quad\Longleftrightarrow\quad
 E\bigl(\mathbb R\setminus(2\pi/\tau)\mathbb Z\bigr)=0.
\]
在这些等价条件成立时，令 \(P_k=E(\{2\pi k/\tau\})\)，写出 \(I\)、\(A\) 的投影级数及 \(D(A)\)。
给出一个满足 \(U(2\pi)=I\) 而 \(A\) 无界的例子。
\end{optionalexercise}
\begin{solution}
记 \(S=(2\pi/\tau)\mathbb Z\)。谱积分的范数公式给出
\[
 \|(U(\tau)-I)x\|^2
 =\int_{\mathbb R}|e^{\mathrm i\tau t}-1|^2\,\mathrm d\mu_x(t).
\]
若 \(E(\mathbb R\setminus S)=0\)，被积函数对每个 \(\mu_x\) 几乎处处为零，
故 \(U(\tau)=I\)。反过来，若左侧对每个 \(x\) 为零，置
\(B_m=\{t:|e^{\mathrm i\tau t}-1|\geq1/m\}\)，则
\[
 0\geq m^{-2}\mu_x(B_m),\qquad \mu_x(B_m)=\|E(B_m)x\|^2.
\]
所以 \(E(B_m)=0\)。由于 \(\mathbb R\setminus S=\bigcup_{m\geq1}B_m\)，
标量测度的可数次可加性给出 \(E(\mathbb R\setminus S)=0\)。

由谱测度的强可数可加性，
\[
 x=\sum_{k\in\mathbb Z}P_kx\quad(x\in H),\qquad
 D(A)=\left\{x\in H:\sum_{k\in\mathbb Z}
       \left(\frac{2\pi k}{\tau}\right)^2\|P_kx\|^2<\infty\right\},
\]
且
\[
 Ax=\sum_{k\in\mathbb Z}\frac{2\pi k}{\tau}P_kx\qquad(x\in D(A)).
\]
两种级数均在 \(H\) 中按有限部分和收敛；第一种由正交分解成立，
第二种由所列平方可和条件及谱积分的截断定义成立。
取 \(H=\ell^2(\mathbb N)\)、\(A(x_n)=(nx_n)\)，
\(D(A)=\{x:\sum n^2|x_n|^2<\infty\}\)。这是实乘法模型，故自伴；
\(\|Ae_n\|=n\) 说明它无界，而 \(U(2\pi)(x_n)=(e^{2\pi\mathrm in}x_n)=(x_n)\)。
\end{solution}

% 自拟：用轨道平均连接算子范数连续性与生成元有界性，补充强连续反例。
\begin{optionalexercise}\label{b1:ex:29-norm-continuous-group}
设 \(H\) 为复 Hilbert 空间，\((U(t))_{t\in\mathbb R}\) 为强连续酉群，
\(A\) 为 Stone 定理给出的自伴算子，满足 \(U(t)=e^{\mathrm itA}\)。证明
\[
 \lim_{t\to0}\|U(t)-I\|=0
 \quad\Longleftrightarrow\quad
 D(A)=H\text{ 且 }A\in\mathcal B(H).
\]
\end{optionalexercise}
\begin{solution}
若 \(A\) 处处定义且有界，则连续函数演算和
\(|e^{\mathrm it\lambda}-1|\leq |t|\cdot|\lambda|\) 给出
\[
 \|U(t)-I\|\leq |t|\cdot\|A\|\longrightarrow0.
\]
零空间的情形直接成立；以下讨论一般的逆向结论。
设 \(\|U(t)-I\|\to0\)。群律及酉性给出
\(\|U(t+h)-U(t)\|=\|U(h)-I\|\)，故轨道在算子范数下连续。
对 \(h>0\)，可在 Banach 空间 \(\mathcal B(H)\) 中作 Riemann 积分，令
\[
 J_h=\frac1h\int_0^h U(s)\,\mathrm ds.
\]
有
\[
 \|J_h-I\|\leq\sup_{0\leq s\leq h}\|U(s)-I\|\longrightarrow0.
\]
取足够小的 \(h\)，使 \(\|J_h-I\|<1\)，则 \(J_h\) 由 Neumann 级数可逆。
设 \(G\) 为差商生成元。对任意 \(x\in H\) 和非零实数 \(r\)，平移积分得到
\[
 \frac{(U(r)-I)J_hx}{r}
 =\frac1h\left(\frac1r\int_h^{h+r}U(s)x\,\mathrm ds
       -\frac1r\int_0^rU(s)x\,\mathrm ds\right).
\]
由轨道连续性，右侧当 \(r\to0\) 时趋于 \((U(h)x-x)/h\)。
所以 \(J_hH\subseteq D(G)\)，且 \(GJ_h=(U(h)-I)/h\)。
但 \(J_hH=H\)，故 \(D(G)=H\)，并且
\[
 G=\frac{U(h)-I}{h}J_h^{-1}\in\mathcal B(H).
\]
Stone 定理给出 \(A=-\mathrm iG\) 以及 \(D(A)=D(G)\)，因而 \(A\) 处处定义且有界。
\end{solution}
